% Les libertés — Allemand Tle A — Cours signé OnBuch+
% Programme officiel MINESEC. Compiler avec : tectonic AL25-les-libertes.tex
\newcommand{\DOCMATIERE}{Allemand}
\newcommand{\DOCNIVEAU}{Tle A}
\newcommand{\DOCMODULE}{Chapitre 5 : Citoyenneté}
\newcommand{\DOCLECON}{AL25}
\newcommand{\DOCTITRE}{Les libertés}
% OnBuch+ — préambule « cours signés » ENRICHI (compilé avec tectonic / XeLaTeX)
% Système visuel « Pop » d'OnBuch+ : fond crème (jamais blanc), bordures encre
% épaisses, ombres « dures » décalées, titres Archivo Black en capitales.
% Version MATHÉMATIQUES : graphes (pgfplots/tikz), boîtes pédagogiques
% (définition, propriété, méthode, attention, le savais-tu, exercice résolu,
% exercice, TP, bilan), page de garde et signature OnBuch+ ancrée en bas.
%
% Macros attendues dans main.tex AVANT \input{preamble} :
%   \DOCMATIERE  \DOCNIVEAU  \DOCMODULE  \DOCLECON (numéro)  \DOCTITRE
\documentclass[11pt]{article}

\usepackage{fontspec}
\usepackage[ngerman,french]{babel} % français = langue principale ; l'allemand via \all{..} / textebox
\usepackage[a4paper,margin=2cm,top=3.4cm,bottom=2.8cm,headheight=1.4cm,footskip=1.3cm]{geometry}
\usepackage{amsmath,amssymb,amsfonts}
\usepackage{mathtools} % \xmapsto, \coloneqq... souvent utilisés par les modèles
\usepackage{cancel} % simplifications barrées (bilans de forces)
% \abs{z} (module, valeur absolue) et \norm{u} (norme), utilisés par les modèles
\providecommand{\abs}{}\let\abs\relax\DeclarePairedDelimiter{\abs}{\lvert}{\rvert}
\providecommand{\norm}{}\let\norm\relax\DeclarePairedDelimiter{\norm}{\lVert}{\rVert}
\usepackage[table]{xcolor} % [table] : fournit aussi \rowcolors (alternance de lignes)
\usepackage{pagecolor}
\usepackage{tikz}
\usetikzlibrary{arrows.meta,positioning,calc,shapes.geometric,shapes.misc,shapes.arrows,decorations.pathmorphing,decorations.pathreplacing,decorations.markings,patterns,shadows,mindmap,trees,fit,backgrounds,matrix,intersections,angles,quotes}
\usepackage{pgfplots}
\usepackage[version=4]{mhchem} % équations nucléaires \ce{^{235}_{92}U}
\usepackage[european,siunitx]{circuitikz} % schémas électriques
\pgfplotsset{compat=1.18}
\usepgfplotslibrary{fillbetween,groupplots} % aires entre courbes (calcul intégral)
\usepackage{enumitem}
\usepackage{fancyhdr}
% [most] charge toutes les bibliothèques tcolorbox (listings, minted, raster,
% documentation...) et gonfle inutilement la mémoire TeX sur un document long
% (~300 boîtes) : on ne charge que ce qui est réellement utilisé.
\usepackage[skins,breakable]{tcolorbox}
\usepackage{graphicx}
\usepackage{colortbl}
\usepackage{booktabs}
\usepackage{tabularx}
\usepackage{diagbox} % case d'en-tête diagonale des tableaux à double entrée
% Colonnes extensibles centrée / gauche / droite (C, L, R), souvent utilisées par les modèles
\newcolumntype{C}{>{\centering\arraybackslash}X}
\newcolumntype{L}{>{\raggedright\arraybackslash}X}
\newcolumntype{R}{>{\raggedleft\arraybackslash}X}
\usepackage{array}
\usepackage{multicol}
\usepackage{multirow}
\usepackage{siunitx}
% parse-numbers=false : accepte aussi \SI{2,5 \times 10^{-3}}{...}
\sisetup{locale=FR,output-decimal-marker={,},group-minimum-digits=4,parse-numbers=false}
\DeclareSIUnit\ohm{\ensuremath{\symup{\Omega}}} % U+2126 absent de la police
\DeclareSIUnit\division{div} % réglages d'oscilloscope (ms/div, V/div)
\DeclareSIUnit\tr{tr} % tours (vitesse de rotation, ex. \SI{1500}{\tr\per\minute})
\DeclareSIUnit\molar{M} % concentration molaire (ex. \SI{3}{\molar}, \SI{1}{\micro\molar})
\DeclareSIUnit\an{an} % année (le modèle confond parfois avec \year, un compteur TeX) — ex. \SI{5}{\kilo\meter\per\an}
\DeclareSIUnit\FCFA{FCFA} % franc CFA (ex. \SI{500}{\FCFA\per\kilo\gram})
\DeclareSIUnit\degreeBrix{\textdegree Brix} % teneur en sucre d'un sirop/jus (réfractométrie)
\DeclareSIUnit\month{mois} % le modèle utilise parfois \month, un compteur TeX, comme unité de durée
\usepackage{qrcode}
% \degree : commande fréquemment utilisée par les modèles pour un angle ou une
% température en dehors de \SI{}{} (non fournie par les paquets chargés).
\providecommand{\degree}{^{\circ}}
% \qed (fin de démonstration), utilisé par les modèles sans amsthm
\providecommand{\qed}{\hfill$\square$}
% \barre : double barre des valeurs interdites dans un tableau de variations (tkz-tab)
\providecommand{\barre}{\ensuremath{\|}}
% environnement proof (amsthm non chargé) utilisé par les modèles
\ifdefined\proof\else\newenvironment{proof}[1][Démonstration]{\par\noindent\textit{#1.}\ }{\qed\par}\fi
\usepackage{lastpage}
\usepackage{hyperref}
\hypersetup{hidelinks}

% ============================================================
% Palette « Pop » OnBuch+ — mêmes hex que la classe P (plus_ui.dart)
% ============================================================
\definecolor{poporange}{HTML}{F59321}
\definecolor{popdark}{HTML}{E07A0C}
\definecolor{popcream}{HTML}{FFF6E8}
\definecolor{popink}{HTML}{1C1714}
\definecolor{poppurple}{HTML}{7A5AE0}
\definecolor{popgreen}{HTML}{2E9E6A}
\definecolor{poppink}{HTML}{E0457B}
\definecolor{popblue}{HTML}{2F7DE1}
\definecolor{popgold}{HTML}{C98A12}
\definecolor{popmuted}{HTML}{8A7F76}
\definecolor{popline}{HTML}{EADFD2}
\definecolor{popgray}{HTML}{8A7F76}
\colorlet{popgrey}{popgray}
% Alias tolérants (le modèle nomme parfois les couleurs différemment)
\colorlet{popwhite}{white}
\colorlet{popblack}{popink}
\colorlet{popred}{poppink}
\colorlet{popyellow}{popgold}
% Teintes claires (fonds de tableaux/encadrés) — le modèle nomme parfois
% les couleurs avec un suffixe L (ex. popblueL) sans les avoir définies.
\colorlet{poporangeL}{poporange!20}
\colorlet{popdarkL}{popdark!20}
\colorlet{poppurpleL}{poppurple!20}
\colorlet{popgreenL}{popgreen!20}
\colorlet{poppinkL}{poppink!20}
\colorlet{popblueL}{popblue!20}
\colorlet{popgoldL}{popgold!20}
\colorlet{popredL}{popred!20}
\colorlet{popgrayL}{popgray!20}
% Teintes claires pour fonds de tableaux / zones
\colorlet{popblueL}{popblue!10!white}
\colorlet{poporangeL}{poporange!14!white}
\colorlet{popgreenL}{popgreen!12!white}
\colorlet{poppurpleL}{poppurple!12!white}
\colorlet{poppinkL}{poppink!12!white}
\colorlet{popgoldL}{popgold!14!white}

\pagecolor{popcream}

% ============================================================
% Typographie
% ============================================================
\defaultfontfeatures{Ligatures=TeX}
\setmainfont{PlusJakartaSans-Medium.ttf}[Path=../../../fonts/,BoldFont=PlusJakartaSans-Bold.ttf]
% Maths en sans-serif (Fira Math) avec chiffres et lettres droites de Plus
% Jakarta, pour que les formules se fondent dans le texte.
\usepackage[math-style=ISO,bold-style=ISO]{unicode-math}
\setmathfont{FiraMath-Regular.otf}[Path=../../../fonts/]
\setmathfont{PlusJakartaSans-Medium.ttf}[Path=../../../fonts/,range={up/num,up/{Latin,latin},\mathup}]
\usepackage{icomma} % virgule décimale sans espace en mode math
% Caractères mathématiques Unicode absents des polices de texte (Plus Jakarta,
% Archivo Black) : rendus par leur équivalent mathématique, en texte comme en
% formule — sinon ils s'affichent en carré vide (ex. « Arithmétique dans ℤ »).
\usepackage{newunicodechar}
\newunicodechar{ℝ}{\ensuremath{\mathbb{R}}}
\newunicodechar{ℤ}{\ensuremath{\mathbb{Z}}}
\newunicodechar{ℂ}{\ensuremath{\mathbb{C}}}
\newunicodechar{ℕ}{\ensuremath{\mathbb{N}}}
\newunicodechar{ℚ}{\ensuremath{\mathbb{Q}}}
\newunicodechar{𝔻}{\ensuremath{\mathbb{D}}}
\newunicodechar{α}{\ensuremath{\alpha}}
\newunicodechar{β}{\ensuremath{\beta}}
\newunicodechar{γ}{\ensuremath{\gamma}}
\newunicodechar{δ}{\ensuremath{\delta}}
\newunicodechar{ε}{\ensuremath{\varepsilon}}
\newunicodechar{θ}{\ensuremath{\theta}}
\newunicodechar{λ}{\ensuremath{\lambda}}
\newunicodechar{μ}{\ensuremath{\mu}}
\newunicodechar{σ}{\ensuremath{\sigma}}
\newunicodechar{τ}{\ensuremath{\tau}}
\newunicodechar{φ}{\ensuremath{\varphi}}
\newunicodechar{ψ}{\ensuremath{\psi}}
\newunicodechar{ω}{\ensuremath{\omega}}
\newunicodechar{Σ}{\ensuremath{\Sigma}}
% majuscules produites par \MakeUppercase dans les titres (α-aminés -> Α)
\newunicodechar{Α}{\ensuremath{\alpha}}
\newunicodechar{Β}{\ensuremath{\beta}}
\newunicodechar{Γ}{\ensuremath{\Gamma}}
\newunicodechar{Δ}{\ensuremath{\Delta}}
\newunicodechar{Ε}{\ensuremath{\varepsilon}}
\newunicodechar{Θ}{\ensuremath{\Theta}}
\newunicodechar{Λ}{\ensuremath{\Lambda}}
\newunicodechar{Μ}{\ensuremath{\mu}}
\newunicodechar{Τ}{\ensuremath{\tau}}
\newunicodechar{Φ}{\ensuremath{\Phi}}
\newunicodechar{Ψ}{\ensuremath{\Psi}}
\newunicodechar{Ω}{\ensuremath{\Omega}}
\newunicodechar{↦}{\ensuremath{\mapsto}}
\newunicodechar{∈}{\ensuremath{\in}}
\newunicodechar{∉}{\ensuremath{\notin}}
\newunicodechar{∧}{\ensuremath{\wedge}}
\newunicodechar{⇔}{\ensuremath{\Leftrightarrow}}
\newunicodechar{⟺}{\ensuremath{\Longleftrightarrow}}
\newunicodechar{⇒}{\ensuremath{\Rightarrow}}
\newunicodechar{⟹}{\ensuremath{\Longrightarrow}}
\newunicodechar{∘}{\ensuremath{\circ}}
\newunicodechar{∪}{\ensuremath{\cup}}
\newunicodechar{∩}{\ensuremath{\cap}}
\newunicodechar{≡}{\ensuremath{\equiv}}
\newunicodechar{⊂}{\ensuremath{\subset}}
\newunicodechar{⊥}{\ensuremath{\perp}}
\newunicodechar{∃}{\ensuremath{\exists}}
\newunicodechar{∀}{\ensuremath{\forall}}
\newunicodechar{∛}{\ensuremath{\sqrt[3]{\,}}}
\newunicodechar{∜}{\ensuremath{\sqrt[4]{\,}}}
\newunicodechar{⃗}{\ensuremath{{}^{\to}}}
\newunicodechar{ⁿ}{\ifmmode^{n}\else\textsuperscript{\ensuremath{n}}\fi}
\newunicodechar{ˣ}{\ifmmode^{x}\else\textsuperscript{\ensuremath{x}}\fi}
\newunicodechar{ᵇ}{\ifmmode^{b}\else\textsuperscript{\ensuremath{b}}\fi}
\newunicodechar{ʳ}{\ifmmode^{r}\else\textsuperscript{\ensuremath{r}}\fi}
\newunicodechar{ᵘ}{\ifmmode^{u}\else\textsuperscript{\ensuremath{u}}\fi}
\newunicodechar{ʸ}{\ifmmode^{y}\else\textsuperscript{\ensuremath{y}}\fi}
\newunicodechar{ᵃ}{\ifmmode^{a}\else\textsuperscript{\ensuremath{a}}\fi}
\newunicodechar{ᵉ}{\ifmmode^{e}\else\textsuperscript{\ensuremath{e}}\fi}
\newunicodechar{ᵏ}{\ifmmode^{k}\else\textsuperscript{\ensuremath{k}}\fi}
\newunicodechar{ᵖ}{\ifmmode^{p}\else\textsuperscript{\ensuremath{p}}\fi}
\newunicodechar{ᴺ}{\ifmmode^{N}\else\textsuperscript{\ensuremath{N}}\fi}
\newunicodechar{ᶜ}{\ifmmode^{c}\else\textsuperscript{\ensuremath{c}}\fi}
\newunicodechar{ʰ}{\ifmmode^{h}\else\textsuperscript{\ensuremath{h}}\fi}
\newunicodechar{ᵠ}{\ifmmode^{q}\else\textsuperscript{\ensuremath{q}}\fi}
\newunicodechar{ₙ}{\ifmmode_{n}\else\textsubscript{\ensuremath{n}}\fi}
\newunicodechar{ₐ}{\ifmmode_{a}\else\textsubscript{\ensuremath{a}}\fi}
\newunicodechar{ᵢ}{\ifmmode_{i}\else\textsubscript{\ensuremath{i}}\fi}
\newunicodechar{ᵣ}{\ifmmode_{r}\else\textsubscript{\ensuremath{r}}\fi}
\newunicodechar{ₖ}{\ifmmode_{k}\else\textsubscript{\ensuremath{k}}\fi}
\newunicodechar{ᵦ}{\ifmmode_{\beta}\else\textsubscript{\ensuremath{\beta}}\fi}
\newunicodechar{ₓ}{\ifmmode_{x}\else\textsubscript{\ensuremath{x}}\fi}
\newfontfamily\popdisplay{ArchivoBlack-Regular.ttf}[Path=../../../fonts/]
\newfontfamily\poplogo{SpaceGrotesk-Bold.ttf}[Path=../../../fonts/]
\newfontfamily\popsemi{PlusJakartaSans-SemiBold.ttf}[Path=../../../fonts/]
\newfontfamily\popxbold{PlusJakartaSans-ExtraBold.ttf}[Path=../../../fonts/]
\newfontfamily\popxboldls{PlusJakartaSans-ExtraBold.ttf}[Path=../../../fonts/,LetterSpace=3.0]

\setlist[enumerate]{leftmargin=1.7em,itemsep=5pt,topsep=4pt}
\setlist[itemize]{leftmargin=1.7em,itemsep=4pt,topsep=4pt,label={\color{poporange}\textbullet}}
\setlength{\parindent}{0pt}
\setlength{\parskip}{5pt}
\renewcommand{\arraystretch}{1.25}

% pgfplots : style Pop par défaut
\pgfplotsset{
  every axis/.append style={axis line style={popink,line width=1pt},tick style={popink},
    grid style={popline,line width=0.5pt},label style={font=\small},tick label style={font=\footnotesize}},
  popcurve/.style={poporange,line width=1.8pt,no markers,smooth},
}
% Même style disponible dans un \draw TikZ simple (hors pgfplots)
\tikzset{popcurve/.style={poporange,line width=1.8pt}}
\tikzset{popgrid/.style={popline,very thin}}
\colorlet{popcurve}{poporange} % le nom du style est parfois utilisé comme couleur

% ============================================================
% Pill / squiggle
% ============================================================
\newcommand{\poppill}[3]{%
  \tikz[baseline=-0.55ex]\node[draw=popink,line width=1.2pt,rounded corners=2.6pt,%
    fill=#2,inner xsep=6.5pt,inner ysep=3pt]{%
    \popxboldls\fontsize{7}{7}\selectfont\color{#3}\MakeUppercase{#1}};%
}
\newcommand{\popsquiggle}[1]{%
  \tikz[baseline=-0.4ex]{
    \draw[#1,line width=2.6pt,line cap=round]
      plot [smooth] coordinates {(0,0.09) (0.34,0.01) (0.68,0.21) (1.02,0.01) (1.36,0.10)};
  }%
}
% Mot-clé surligné (marqueur orange sous le texte)
\newcommand{\cle}[1]{{\popxbold\color{popdark}#1}}

% ============================================================
% En-tête / pied
% ============================================================
\pagestyle{fancy}
\fancyhf{}
\renewcommand{\headrulewidth}{0pt}
\renewcommand{\footrulewidth}{0pt}
\lhead{\raisebox{-0.15ex}{\poplogo\fontsize{13}{13}\selectfont\color{popink}OnBuch\color{poporange}+}}
\rhead{\poppill{\DOCMATIERE\ \textbullet\ \DOCNIVEAU\ \textbullet\ Leçon \DOCLECON}{white}{popink}}
\lfoot{\popxbold\fontsize{7.6}{7.6}\selectfont\color{popmuted}\MakeUppercase{Cours signé OnBuch+}\ \textbullet\ {\popsemi\DOCTITRE}}
\rfoot{\tikz[baseline=-0.6ex]\node[rounded corners=8.5pt,fill=popink,minimum height=17pt,minimum width=17pt,inner xsep=5pt,inner sep=0]{\popxbold\fontsize{8}{8}\selectfont\color{popcream}\thepage\,/\,\pageref*{LastPage}};}

% ============================================================
% Titres
% ============================================================
\newcommand{\chaptitre}[2]{%
  \begin{tcolorbox}[enhanced,colback=poporange,colframe=popink,boxrule=2.6pt,arc=11pt,
      drop shadow=popink,left=15pt,right=15pt,top=12pt,bottom=11pt,before skip=3pt,after skip=18pt]
    \poppill{#2}{popink}{popcream}\par\vspace{9pt}
    {\raggedright\hyphenpenalty=10000\exhyphenpenalty=10000\popdisplay\fontsize{22}{24}\selectfont\color{popcream}\MakeUppercase{#1}\par}
  \end{tcolorbox}
}
% Section numérotée : pastille orange avec le numéro + titre Archivo
\newcounter{popsec}
\newcommand{\coursec}[1]{%
  \stepcounter{popsec}\setcounter{popsub}{0}%
  \par\vspace{16pt}\noindent
  \begin{minipage}{\linewidth}
  \tikz[baseline=-0.7ex]\node[draw=popink,line width=1.6pt,fill=poporange,rounded corners=4pt,minimum size=22pt,inner sep=0]{\popdisplay\fontsize{12}{12}\selectfont\color{popcream}\thepopsec};%
  \hspace{8pt}{\popdisplay\fontsize{15}{17}\selectfont\color{popink}\MakeUppercase{#1}}\par
  \vspace{4pt}\noindent{\color{poporange}\rule{36pt}{3.6pt}}
  \end{minipage}\par\nopagebreak\vspace{8pt}
}
\newcounter{popsub}
\newcommand{\courssub}[1]{%
  \stepcounter{popsub}%
  \par\vspace{9pt}\noindent{\popxbold\fontsize{11.5}{13}\selectfont\color{popink}{\color{poporange}\thepopsec.\thepopsub}\hspace{6pt}#1}\par\nopagebreak\vspace{4pt}
}

% ============================================================
% Boîtes pédagogiques — même recette : fond blanc, bordure épaisse colorée,
% ombre dure encre, badge plein. Titre optionnel [#1].
% ============================================================
\tcbset{popbase/.style={breakable,enhanced,colback=white,boxrule=2.3pt,arc=9pt,
  shadow={3pt}{-3pt}{0pt}{popink},left=14pt,right=14pt,top=11pt,bottom=11pt,before skip=11pt,after skip=15pt,
  fontupper=\popsemi\fontsize{10.6}{14.5}\selectfont\color{popink},
  fontlower=\popsemi\fontsize{10.6}{14.5}\selectfont\color{popink}}}
% Coche dessinée (le glyphe ✓ n'existe pas dans les polices du document)
\newcommand{\popcheck}[1]{\tikz[baseline=-0.5ex]\draw[#1,line width=1.6pt,line cap=round,line join=round](-0.13,0.0)--(-0.04,-0.09)--(0.13,0.1);}
\providecommand{\coche}{\popcheck{popgreen}}
\providecommand{\resultat}[1]{\par\smallskip\noindent\fcolorbox{popgreen}{popgreenL}{\parbox{\dimexpr\linewidth-2\fboxsep-2\fboxrule\relax}{\textbf{Résultat.} #1}}\par\smallskip}
\newcommand{\popboxhead}[3]{\poppill{#1}{#2}{white}\ifx\relax#3\relax\else\hspace{6pt}{\popxbold\fontsize{10.6}{12}\selectfont\color{popink}#3}\fi\par\vspace{7pt}}

\newtcolorbox{aretenir}[1][]{popbase,colframe=popgreen,before upper={\popboxhead{À retenir}{popgreen}{#1}}}
% Texte en allemand (document de lecture, dialogue, extrait) : cadre
% d'étude. Un titre optionnel (source, auteur) s'affiche dans l'en-tête.
\newtcolorbox{textebox}[1][]{popbase,colframe=popblue,before upper={\popboxhead{Text}{popblue}{#1}\begin{otherlanguage}{ngerman}},after upper={\end{otherlanguage}}}
% Marques ✓ / ✗ (forme correcte / fautive) souvent utilisées par les modèles
\providecommand{\cross}{\textcolor{poppink}{\ensuremath{\times}}}
\providecommand{\xmark}{\cross}\providecommand{\crossmark}{\cross}\providecommand{\cmark}{\ensuremath{\checkmark}}
% Ligne de réponse / justification à compléter (inventée par les modèles)
\providecommand{\justification}[1]{\par\noindent\textit{Justification~:} \dotfill\par}
% \textipa (package tipa non chargé) : la police ne couvre pas l'API -> simple passage
\providecommand{\textipa}[1]{#1}
% Soulignement (ulem non chargé) : \uline, \ul
\providecommand{\uline}[1]{\underline{#1}}\providecommand{\ul}[1]{\underline{#1}}
% \glossaire{\item ...} inventée par les modèles : liste de glossaire
\providecommand{\glossaire}[1]{\begin{itemize}#1\end{itemize}}
% \alert (beamer) inventée par les modèles : texte mis en évidence
\providecommand{\alert}[1]{\textbf{\textcolor{poppink}{#1}}}
% variantes ASCII d'ordinaux écrites en macro
\providecommand{\iere}{\textsuperscript{re}}\providecommand{\ieme}{\textsuperscript{e}}\providecommand{\ere}{\textsuperscript{re}}\providecommand{\eme}{\textsuperscript{e}}\providecommand{\er}{\textsuperscript{er}}\providecommand{\ie}{\textsuperscript{e}}
% Ordinaux français écrits en macro par les modèles : 1\ière, 3\ième
\newcommand{\ière}{\textsuperscript{re}}\newcommand{\ième}{\textsuperscript{e}}
% \exemple (inventée par les modèles) : étiquette « Exemple : »
\providecommand{\exemple}{\textbf{Exemple~:}\ }
% \square utilisable aussi hors mode math (cases à cocher)
\AtBeginDocument{\let\squareMath\square\renewcommand{\square}{\ensuremath{\squareMath}}}
% Mot ou phrase en allemand à l'intérieur d'un paragraphe français
\newcommand{\all}[1]{\ifmmode\text{\textit{#1}}\else\begingroup\selectlanguage{ngerman}\itshape #1\endgroup\fi}
\let\esp\all % alias toléré
\newtcolorbox{exemplebox}[1][]{popbase,colframe=poporange,before upper={\popboxhead{Exemple}{poporange}{#1}}}
\newtcolorbox{definition}[1][]{popbase,colframe=popblue,before upper={\popboxhead{Définition}{popblue}{#1}}}
\newtcolorbox{propriete}[1][]{popbase,colframe=poppurple,before upper={\popboxhead{Propriété}{poppurple}{#1}}}
\newtcolorbox{methode}[1][]{popbase,colframe=popgold,before upper={\popboxhead{Méthode}{popgold}{#1}}}
\newtcolorbox{attention}[1][]{popbase,colframe=poppink,before upper={\popboxhead{Attention}{poppink}{#1}}}
\newtcolorbox{experience}[1][]{popbase,colframe=popblue,colback=popblueL,before upper={\popboxhead{Expérience}{popblue}{#1}}}
% « Le savais-tu ? » : carte violette pleine (contexte, culture, Cameroun)
\newtcolorbox{savaistu}[1][]{popbase,colback=poppurple,colframe=popink,
  fontupper=\popsemi\fontsize{10.4}{14.2}\selectfont\color{white},
  before upper={\poppill{Le savais-tu\,?}{white}{poppurple}\ifx\relax#1\relax\else\hspace{6pt}{\popxbold\color{white}#1}\fi\par\vspace{7pt}}}
% Exercice résolu : énoncé en haut, \tcblower puis la solution sur fond crème
\newcounter{popexo}\newcounter{popexr}
\newtcolorbox{exoresolu}[1][]{popbase,colframe=poporange,colbacklower=popcream,
  segmentation style={popink,line width=1.2pt,dashed},
  before upper={\stepcounter{popexr}\popboxhead{Exercice résolu \thepopexr}{poporange}{#1}},
  before lower={\poppill{Solution}{popink}{popcream}\par\vspace{6pt}}}
% Exercice (non résolu en place) — la correction est dans la partie Corrigés
\newtcolorbox{exercice}[1][]{popbase,colframe=popink,
  before upper={\stepcounter{popexo}\popboxhead{Exercice \thepopexo}{popink}{#1}}}
% Corrigé
\newtcolorbox{corrige}[1][]{popbase,colframe=popgreen,colback=popgreenL,
  before upper={\popboxhead{Corrigé}{popgreen}{#1}}}
% Objectifs / prérequis / bilan
\newtcolorbox{objectifs}{popbase,colframe=popink,colback=poporangeL,
  before upper={\popboxhead{Objectifs de la leçon}{popink}{}}}
\newtcolorbox{prerequis}{popbase,colframe=popink,colback=popblueL,
  before upper={\popboxhead{Prérequis}{popblue}{}}}
\newtcolorbox{bilan}{popbase,colframe=popink,colback=white,boxrule=2.8pt,
  before upper={\popboxhead{Fiche bilan}{poporange}{L'essentiel à connaître pour l'examen}}}
% Figure encadrée avec légende
\newtcolorbox{popfigure}[1][]{enhanced,breakable=false,colback=white,colframe=popline,boxrule=1.4pt,arc=8pt,
  left=10pt,right=10pt,top=10pt,bottom=8pt,before skip=10pt,after skip=12pt,halign=center,
  lower separated=false,halign lower=center,
  fontlower=\popsemi\fontsize{9}{11.5}\selectfont\color{popmuted},
  #1}
\newcommand{\legende}[1]{\par\vspace{6pt}{\popsemi\fontsize{9}{11.5}\selectfont\color{popmuted}#1}}

% ============================================================
% Signature OnBuch+
% ============================================================
\newcommand{\poplogoinline}[1]{{\poplogo\fontsize{#1}{#1}\selectfont\color{popink}OnBuch\color{poporange}+}}
\newcommand{\signatureonbuch}{%
  \begin{tcolorbox}[enhanced,colback=popink,colframe=popink,boxrule=2.2pt,arc=10pt,
      drop shadow=poporange,left=14pt,right=14pt,top=10pt,bottom=10pt,before skip=0pt,after skip=0pt]
    \tikz[baseline=-0.7ex]\node[circle,fill=popgreen,draw=popcream,line width=1.3pt,minimum size=22pt,inner sep=0]{\popcheck{white}};%
    \hspace{9pt}\begin{minipage}[c]{0.62\linewidth}
      {\popxbold\fontsize{12}{13}\selectfont\color{popcream}Signé\ \textbullet\ {\poplogo OnBuch\color{poporange}+}}\par\vspace{2pt}
      {\raggedright\popsemi\fontsize{8}{10.5}\selectfont\color{popline}\MakeUppercase{Équipe pédagogique OnBuch+}\\\MakeUppercase{Conforme au programme officiel MINESEC}\par}
    \end{minipage}\hfill
    {\poplogo\fontsize{22}{22}\selectfont\color{popcream}OnBuch\color{poporange}+}
  \end{tcolorbox}%
}

% ============================================================
% Page de garde
% #1 titre · #2 matière · #3 niveau · #4 module · #5 n° leçon · #6 durée ·
% #7 description / objectifs (texte ou itemize)
% ============================================================
\newcommand{\pagegarde}[7]{%
  \thispagestyle{empty}
  \newgeometry{a4paper,margin=1.7cm,top=1.6cm,bottom=1.4cm}%
  \begin{tikzpicture}[remember picture,overlay]
    \fill[poporange!18] ([xshift=-3.2cm,yshift=-3.6cm]current page.north east) circle (4.6cm);
    \fill[poppurple!14] ([xshift=2.4cm,yshift=7.6cm]current page.south west) circle (3.8cm);
  \end{tikzpicture}%
  {\poplogo\fontsize{30}{30}\selectfont\color{popink}OnBuch\color{poporange}+}\hfill
  \raisebox{0.6ex}{\poppill{Cours signé \textbullet\ Premium}{popink}{popcream}}\par
  \vspace{1pt}\popsquiggle{poppurple}\par
  \vspace{8pt}
  \begin{tcolorbox}[enhanced,colback=poporange,colframe=popink,boxrule=2.8pt,arc=13pt,
      drop shadow=popink,left=18pt,right=18pt,top=12pt,bottom=13pt,after skip=10pt]
    \poppill{#2 \textbullet\ #3}{popink}{popcream}\hspace{5pt}\poppill{#4}{white}{popink}\par\vspace{8pt}
    {\popdisplay\fontsize{14}{14}\selectfont\color{popink}LEÇON #5}\par\vspace{4pt}
    {\raggedright\hyphenpenalty=10000\exhyphenpenalty=10000\popdisplay\fontsize{25}{27}\selectfont\color{popcream}\MakeUppercase{#1}\par}
    \vspace{8pt}
    {\popxbold\fontsize{9.5}{9.5}\selectfont\color{popink}\MakeUppercase{Durée conseillée : #6}}
  \end{tcolorbox}
  \begin{tcolorbox}[enhanced,colback=white,colframe=popink,boxrule=2.2pt,arc=10pt,
      drop shadow=popink,left=16pt,right=16pt,top=10pt,bottom=10pt,after skip=8pt]
    \poppill{Dans cette leçon}{poppurple}{white}\par\vspace{5pt}
    \popsemi\fontsize{9.3}{12.6}\selectfont\color{popink}#7
  \end{tcolorbox}
  \noindent\begin{minipage}[t]{0.487\linewidth}
    \begin{tcolorbox}[enhanced,colback=popgreen,colframe=popink,boxrule=2.2pt,arc=10pt,
        drop shadow=popink,left=11pt,right=11pt,top=10pt,bottom=10pt,height=3.1cm,valign=center]
      \tcbox[colback=white,colframe=popink,boxrule=1.3pt,arc=5pt,boxsep=0pt,nobeforeafter,
        left=3pt,right=3pt,top=3pt,bottom=3pt,valign=center]{\qrcode[height=1.9cm]{https://play.google.com/store/apps/details?id=com.luvvix.onbuch&hl=fr}}%
      \hspace{8pt}\begin{minipage}[c]{0.5\linewidth}
        {\popdisplay\fontsize{11}{12.5}\selectfont\color{white}\MakeUppercase{Télécharge OnBuch}}\par\vspace{4pt}
        {\raggedright\popsemi\fontsize{8.4}{11}\selectfont\color{white}Gratuit \textbullet\ Google Play\\Tuteur IA, annales, concours, résultats\par}
      \end{minipage}
    \end{tcolorbox}
  \end{minipage}\hfill
  \begin{minipage}[t]{0.487\linewidth}
    \begin{tcolorbox}[enhanced,colback=poppurple,colframe=popink,boxrule=2.2pt,arc=10pt,
        drop shadow=popink,left=11pt,right=11pt,top=10pt,bottom=10pt,height=3.1cm,valign=center]
      \tikz[baseline=-0.7ex]\node[circle,fill=white,draw=popink,line width=1.3pt,minimum size=1.6cm,inner sep=0]{\popdisplay\fontsize{24}{24}\selectfont\color{poppurple}+};%
      \hspace{8pt}\begin{minipage}[c]{0.6\linewidth}
        {\popdisplay\fontsize{11}{12.5}\selectfont\color{white}\MakeUppercase{Passe à OnBuch+}}\par\vspace{4pt}
        {\raggedright\popsemi\fontsize{8.4}{11}\selectfont\color{white}Tous les cours signés, Tuteur IA illimité, fiches PDF — onglet OnBuch+ de l'app\par}
      \end{minipage}
    \end{tcolorbox}
  \end{minipage}
  \vspace{8pt}
  \popcontenu
  \vfill
  \signatureonbuch
  \restoregeometry
  \newpage
}

% Rangée « Ce cours contient » (page de garde)
\newcommand{\poptile}[3]{% #1 couleur #2 gros texte #3 légende
  \begin{tcolorbox}[enhanced,colback=#1,colframe=popink,boxrule=2pt,arc=9pt,shadow={2.5pt}{-2.5pt}{0pt}{popink},
      left=6pt,right=6pt,top=9pt,bottom=9pt,halign=center,valign=center,height=1.75cm,width=\linewidth,nobeforeafter]
    {\popdisplay\fontsize{12.5}{13}\selectfont\color{white}\MakeUppercase{#2}}\par\vspace{3pt}
    {\centering\popsemi\fontsize{7.8}{9.5}\selectfont\color{white}#3\par}
  \end{tcolorbox}}
\newcommand{\popcontenu}{%
  \par\noindent{\popxboldls\fontsize{7.5}{7.5}\selectfont\color{popmuted}CE COURS SIGNÉ CONTIENT}\par\vspace{7pt}
  \noindent\begin{minipage}[t]{0.235\linewidth}\poptile{popblue}{Cours}{complet, illustré, pas à pas}\end{minipage}\hfill
  \begin{minipage}[t]{0.235\linewidth}\poptile{poporange}{Méthodes}{et exercices résolus}\end{minipage}\hfill
  \begin{minipage}[t]{0.235\linewidth}\poptile{poppink}{Exercices}{type examen + corrigés}\end{minipage}\hfill
  \begin{minipage}[t]{0.235\linewidth}\poptile{popgreen}{Fiche bilan}{l'essentiel à retenir}\end{minipage}\par}

% Sommaire
\newtcolorbox{sommaire}{enhanced,colback=white,colframe=popink,boxrule=2.2pt,arc=10pt,
  drop shadow=popink,left=18pt,right=18pt,top=13pt,bottom=13pt,before skip=6pt,after skip=20pt,
  before upper={\poppill{Sommaire}{poppurple}{white}\par\vspace{9pt}\popsemi\fontsize{10.6}{14.5}\selectfont\color{popink}}}

% Signature de fin de cours — poussée tout en bas de la dernière page
\newcommand{\findecours}{%
  \par\vfill\nopagebreak
  \begin{center}\popsquiggle{poporange}\end{center}\vspace{4pt}
  \signatureonbuch
}
\AtBeginDocument{\def\llbracket{\mathopen{[\mkern-3.5mu[}}\def\rrbracket{\mathclose{]\mkern-3.5mu]}}} % intervalles d'entiers (glyphe ⟦ absent de la police)
% Glyphes absents de Fira Math : redéfinis à partir de glyphes disponibles
\AtBeginDocument{%
  \def\setminus{\mathbin{\backslash}}\let\smallsetminus\setminus
  \def\vdots{\mathord{\vbox{\baselineskip3.2pt\lineskiplimit0pt\kern4pt\hbox{.}\hbox{.}\hbox{.}}}}%
  \def\ddots{\mathinner{\mkern1mu\raise6.5pt\hbox{.}\mkern2mu\raise3.6pt\hbox{.}\mkern2mu\raise0.7pt\hbox{.}\mkern1mu}}%
  \def\triangle{\mathord{\scalebox{0.9}{$\symup{\Delta}$}}}%
  \def\nabla{\mathord{\raisebox{\depth}{\scalebox{1}[-1]{$\symup{\Delta}$}}}}%
  \def\checkmark{\popcheck{popdark}}%
  \def\star{\mathbin{\ast}}\def\bigstar{\mathord{\ast}}%
  \def\diamond{\mathbin{\scalebox{0.6}{$\Diamond$}}}%
  \def\bigcup{\mathop{\mathchoice{\raisebox{-0.3em}{\scalebox{1.7}{$\cup$}}}{\scalebox{1.25}{$\cup$}}{\cup}{\cup}}\displaylimits}%
  \def\bigcap{\mathop{\mathchoice{\raisebox{-0.3em}{\scalebox{1.7}{$\cap$}}}{\scalebox{1.25}{$\cap$}}{\cap}{\cap}}\displaylimits}%
  \def\lesseqqgtr{\mathrel{\lessgtr}}\def\bigoplus{\mathop{\oplus}\displaylimits}%
}
\providecommand{\ssi}{si et seulement si} % abréviation fréquente
\AtBeginDocument{\providecommand{\wideparen}{\overparen}} % arc de cercle

%% FIGFIX-BEGIN v5 — correctifs globaux des figures (docs/onbuch-generation/tools/figfix)
\usepackage{adjustbox}\usepackage{etoolbox}\tcbuselibrary{raster}
% toute figure TikZ/pgfplots plus large que la ligne est réduite à la largeur disponible
\BeforeBeginEnvironment{tikzpicture}{\begin{adjustbox}{max width=\linewidth}}
\AfterEndEnvironment{tikzpicture}{\end{adjustbox}}
% légendes pgfplots lisibles par-dessus les courbes
\pgfplotsset{every axis/.append style={legend style={fill=white,fill opacity=0.92,text opacity=1,draw=black!15,inner sep=3pt}}}
% étiquettes d'arêtes (syntaxe edge["..."]) avec fond blanc
\tikzset{every edge quotes/.append style={fill=white,inner sep=1.2pt}}
%% FIGFIX-END
\begin{document}

\pagegarde{Les libertés}{Allemand}{Tle A}{Chapitre 5 : Citoyenneté}{AL25}{2 h}{Cette leçon de type « texte » propose la lecture et le commentaire d'un texte allemand sur les grandes libertés (expression, presse, religion, réunion). Elle permet d'acquérir le vocabulaire thématique des droits et devoirs, de réviser les verbes de modalité et d'apprendre à exprimer la concession avec obwohl et trotzdem.
\vspace{4pt}
\begin{multicols}{2}\small
\begin{itemize}
\item À la fin de cette leçon, je sais comprendre globalement puis en détail un texte allemand sur les libertés fondamentales.
\item À la fin de cette leçon, je sais employer correctement le vocabulaire des droits et devoirs (die Freiheit, das Recht, die Pflicht, erlauben, verbieten) avec articles et pluriels.
\item À la fin de cette leçon, je sais conjuguer et placer correctement les verbes de modalité (dürfen, müssen, können, sollen, wollen) dans une phrase.
\item À la fin de cette leçon, je sais exprimer la concession avec obwohl (subordonnée) et trotzdem (phrase principale).
\item À la fin de cette leçon, je sais répondre par écrit à des questions de compréhension en phrases complètes.
\item À la fin de cette leçon, je sais résumer un texte en quelques phrases avec mes propres mots.
\end{itemize}\end{multicols}}

\begin{sommaire}
\begin{multicols}{2}
\begin{enumerate}
\item Texte support : « Freiheiten in einer Demokratie »
\item Compréhension du texte
\item Lexique thématique : droits, libertés et devoirs
\item Grammaire 1 : les verbes de modalité
\item Grammaire 2 : la concession (obwohl, trotzdem)
\item Méthode du commentaire et commentaire modèle
\item Production guidée écrite et orale
\item Activité d'intégration
\item Méthodes et exercices résolus
\item Exercices
\item Corrigés des exercices
\item Fiche bilan
\end{enumerate}
\end{multicols}
\end{sommaire}

\begin{objectifs}
\begin{itemize}
\item À la fin de cette leçon, je sais comprendre globalement puis en détail un texte allemand sur les libertés fondamentales.
\item À la fin de cette leçon, je sais employer correctement le vocabulaire des droits et devoirs (die Freiheit, das Recht, die Pflicht, erlauben, verbieten) avec articles et pluriels.
\item À la fin de cette leçon, je sais conjuguer et placer correctement les verbes de modalité (dürfen, müssen, können, sollen, wollen) dans une phrase.
\item À la fin de cette leçon, je sais exprimer la concession avec obwohl (subordonnée) et trotzdem (phrase principale).
\item À la fin de cette leçon, je sais répondre par écrit à des questions de compréhension en phrases complètes.
\item À la fin de cette leçon, je sais résumer un texte en quelques phrases avec mes propres mots.
\item À la fin de cette leçon, je sais rédiger un court paragraphe d'opinion sur une liberté fondamentale.
\item À la fin de cette leçon, je sais comparer les libertés au Cameroun et dans les pays germanophones à l'oral.
\end{itemize}
\end{objectifs}

\begin{prerequis}
\begin{itemize}
\item La phrase simple allemande : place du verbe conjugué en 2e position et de l'infinitif en fin de phrase (Seconde/Première)
\item Les verbes de modalité au présent, niveau élémentaire (dürfen, müssen, können)
\item Les subordonnées avec weil, dass : verbe conjugué en fin de proposition
\item Le vocabulaire de la vie en société vu en Première (die Gesellschaft, der Bürger)
\item Les articles définis et indéfinis au nominatif et à l'accusatif
\end{itemize}
\end{prerequis}

\newpage

\coursec{Texte support : « Freiheiten in einer Demokratie »}

\courssub{Situation d'accroche et problématique}

En février, un débat anime un lycée de Yaoundé : des élèves de Terminale veulent publier dans le journal scolaire un article critique sur l'organisation de la cantine. Le chef d'établissement hésite : a-t-il le droit d'interdire cet article ? Les élèves, eux, affirment : « Nous avons le droit de donner notre opinion ! » Qui a raison ?

En Allemagne, la réponse est claire : l'article 5 de la Loi fondamentale (\all{das Grundgesetz}) garantit à tous les citoyens, y compris aux jeunes, la liberté d'expression et la liberté de la presse. En Autriche et en Suisse, ces libertés sont également inscrites dans la constitution. Mais attention : toute liberté a des \cle{limites} (\all{Grenzen}) et implique des \cle{devoirs} (\all{Pflichten}). On ne peut pas tout dire, tout écrire, tout faire.

\textbf{Problématique :} quelles sont les grandes libertés dans une démocratie, quelles sont leurs limites, et comment en parler correctement en allemand ? Pour répondre à ces questions, nous allons lire un article de journal destiné à la jeunesse allemande.

\courssub{Présentation du texte}

\begin{definition}[Présentation du document]
\begin{itemize}
\item \textbf{Nature :} article de journal pour la jeunesse (texte fictif, type magazine scolaire allemand).
\item \textbf{Titre :} \all{Freiheiten in einer Demokratie} (« Les libertés dans une démocratie »).
\item \textbf{Thème :} les libertés fondamentales : expression, presse, religion, réunion.
\item \textbf{Longueur :} environ 12 lignes, niveau A2/B1 (accessible en Terminale A).
\item \textbf{Objectif :} repérer les quatre libertés citées, comprendre les limites des droits, découvrir les verbes de modalité et les connecteurs de concession.
\end{itemize}
\end{definition}

\courssub{Lecture globale : stratégies}

\begin{methode}[Comment aborder une première lecture]
\begin{enumerate}
\item \textbf{Lecture silencieuse} (2 minutes) : lis le texte sans dictionnaire. Ne t'arrête pas sur les mots inconnus.
\item \textbf{Repérage des mots transparents} : souligne les mots qui ressemblent au français : \all{Demokratie} (démocratie), \all{Presse} (presse), \all{Religion} (religion), \all{Politik} (politique), \all{friedlich} (pacifique). Ces mots te donnent déjà le thème général.
\item \textbf{Écoute} : le professeur lit le texte à voix haute. Suis du doigt et note la prononciation : \all{Freiheit} se prononce « fraï-haït », \all{Bürger} « biur-gueur », \all{Meinung} « maï-noung », \all{versammeln} « fèr-zam-meln ».
\item \textbf{Hypothèse globale} : de quoi parle le texte ? Formule une phrase en français avant toute traduction.
\end{enumerate}
\end{methode}

\courssub{Texte intégral}

\begin{textebox}[Freiheiten in einer Demokratie — article de journal pour la jeunesse]
In einer Demokratie haben die Bürger wichtige Freiheiten. Jeder darf seine Meinung frei sagen: das ist die Meinungsfreiheit. Die Presse darf über Politik berichten, ohne Angst zu haben: das nennt man Pressefreiheit. Außerdem hat jeder das Recht, seine Religion frei zu wählen, und die Bürger dürfen sich friedlich versammeln: das ist die Versammlungsfreiheit. Obwohl diese Freiheiten wichtig sind, haben sie auch Grenzen: man darf zum Beispiel niemanden beleidigen. Trotzdem müssen wir für unsere Rechte kämpfen, denn Freiheit ist ein kostbares Gut.
\end{textebox}

\textbf{Glossaire (lexique en marge) :}

\begin{center}
\begin{tabularx}{\linewidth}{|l|X|l|}\hline
\rowcolor{popblueL} \textbf{Mot allemand} & \textbf{Traduction} & \textbf{Remarque} \\ \hline
\all{die Freiheit (-en)} & la liberté & féminin \\ \hline
\all{der Bürger (-)} & le citoyen & pluriel identique \\ \hline
\all{die Meinung (-en)} & l'opinion, l'avis & \all{seine Meinung sagen} = dire son avis \\ \hline
\all{die Presse (-n)} & la presse & mot transparent \\ \hline
\all{das Recht (-e)} & le droit & \all{das Recht auf} + accusatif = le droit à/de \\ \hline
\all{die Religion (-en)} & la religion & mot transparent \\ \hline
\all{sich versammeln} & se rassembler, se réunir & verbe réfléchi, particule inséparable \\ \hline
\all{friedlich} & pacifiquement, pacifique & adverbe / adjectif \\ \hline
\all{die Grenze (-n)} & la limite, la frontière & féminin \\ \hline
\all{beleidigen} & insulter, offenser & verbe faible, transitif \\ \hline
\all{kämpfen} & lutter, combattre & \all{für} + accusatif (\all{für unsere Rechte}) \\ \hline
\all{kostbar} & précieux & adjectif \\ \hline
\all{das Gut (\-"er)} & le bien (économique/moral) & pluriel \all{die Güter} \\ \hline
\all{die Angst (-e)} & la peur & \all{ohne Angst zu haben} = sans avoir peur \\ \hline
\all{obwohl} & bien que & conjonction de concession \\ \hline
\all{trotzdem} & pourtant, néanmoins & adverbe de concession \\ \hline
\end{tabularx}
\end{center}

\courssub{Repérage des quatre libertés citées}

Le texte nomme explicitement quatre libertés fondamentales. Observons leur formation (mots composés) :

\begin{exemplebox}[Les quatre libertés du texte]
\begin{itemize}
\item \all{die Meinungsfreiheit} : \all{Jeder darf seine Meinung frei sagen.} « Chacun a le droit d'exprimer librement son opinion. »
\item \all{die Pressefreiheit} : \all{Die Presse darf über Politik berichten, ohne Angst zu haben.} « La presse peut rendre compte de la politique sans avoir peur. »
\item \all{die Religionsfreiheit} : \all{Jeder hat das Recht, seine Religion frei zu wählen.} « Chacun a le droit de choisir librement sa religion. »
\item \all{die Versammlungsfreiheit} : \all{Die Bürger dürfen sich friedlich versammeln.} « Les citoyens ont le droit de se rassembler pacifiquement. »
\end{itemize}
\end{exemplebox}

\begin{definition}[La liberté d'expression]
\all{die Meinungsfreiheit} = \all{das Recht, seine Meinung frei zu sagen oder zu schreiben} : le droit d'exprimer librement son opinion, à l'oral ou à l'écrit. C'est la première liberté citée dans le texte, car elle est à la base de toutes les autres.
\end{definition}

\begin{popfigure}
\begin{center}\tcbox[colback=popblue,colframe=popblue,coltext=white,arc=9pt,boxsep=4pt,left=6pt,right=6pt]{\bfseries\small die Freiheiten}\end{center}
\vspace{-2pt}
\begin{tcbraster}[raster columns=2,raster equal height=rows,raster column skip=6pt,raster row skip=6pt,enhanced,blankest]
\begin{tcolorbox}[enhanced,colback=poporange,colframe=poporange,coltext=white,boxrule=0.9pt,arc=3pt,left=4pt,right=4pt,top=3pt,bottom=3pt,boxsep=1pt,halign=left]\bfseries\scriptsize die Meinungs- freiheit (expression)\end{tcolorbox}
\begin{tcolorbox}[enhanced,colback=popgreen,colframe=popgreen,coltext=white,boxrule=0.9pt,arc=3pt,left=4pt,right=4pt,top=3pt,bottom=3pt,boxsep=1pt,halign=left]\bfseries\scriptsize die Presse- freiheit (presse)\end{tcolorbox}
\begin{tcolorbox}[enhanced,colback=popgold,colframe=popgold,coltext=white,boxrule=0.9pt,arc=3pt,left=4pt,right=4pt,top=3pt,bottom=3pt,boxsep=1pt,halign=left]\bfseries\scriptsize die Religions- freiheit (religion)\end{tcolorbox}
\begin{tcolorbox}[enhanced,colback=poppurple,colframe=poppurple,coltext=white,boxrule=0.9pt,arc=3pt,left=4pt,right=4pt,top=3pt,bottom=3pt,boxsep=1pt,halign=left]\bfseries\scriptsize die Versammlungs- freiheit (réunion)\end{tcolorbox}
\end{tcbraster}

\legende{Carte mentale : les quatre libertés fondamentales du texte}
\end{popfigure}

\courssub{Familles de mots et dérivation (Vocabulaire étendu)}

\begin{propriete}[Familles de mots : Freiheit, Recht, Pflicht]
\begin{center}
\begin{tabularx}{\linewidth}{|X|X|X|X|}\hline
\rowcolor{popblueL} \textbf{Verbe} & \textbf{Nom (article + pluriel)} & \textbf{Adjectif / Adverbe} & \textbf{Français} \\ \hline
\all{frei sein} & \all{die Freiheit (-en)} & \all{frei} & être libre / la liberté / libre \\ \hline
\all{meinen} & \all{die Meinung (-en)} & \all{meiner Meinung nach} & penser / l'opinion / à mon avis \\ \hline
\all{berichten} & \all{der Bericht (-e)} & \all{berichtet} & rapporter / le rapport / rapporté \\ \hline
\all{wählen} & \all{die Wahl (-en)} & \all{wählbar} & choisir / le choix, l'élection / éligible \\ \hline
\all{versammeln} & \all{die Versammlung (-en)} & \all{versammelt} & rassembler / la réunion / réuni \\ \hline
\all{beleidigen} & \all{die Beleidigung (-en)} & \all{beleidigend} & insulter / l'insulte / insultant \\ \hline
\all{kämpfen} & \all{der Kampf (\"e)} & \all{kämpferisch} & lutter / la lutte / combatif \\ \hline
\all{verbieten} & \all{das Verbot (-e)} & \all{verboten} & interdire / l'interdiction / interdit \\ \hline
\all{erlauben} & \all{die Erlaubnis (-se)} & \all{erlaubt} & permettre / la permission / permis \\ \hline
\all{pflichen} & \all{die Pflicht (-en)} & \all{pflichtbewusst} & (verbe rare) / le devoir / consciencieux \\ \hline
\end{tabularx}
\end{center}
\textbf{Attention :} \all{das Recht} (le droit) $\neq$ \all{richtig} (correct) $\neq$ \all{rechts} (à droite). \all{die Pflicht} (le devoir) s'oppose à \all{das Recht} (le droit).
\end{propriete}

\courssub{Repérage grammatical dans le texte}

Avant d'étudier la grammaire en détail (sections 4 et 5), repérons dans le texte deux phénomènes essentiels.

\textbf{1. Les verbes de modalité.} Souligne dans le texte : \all{darf} (2 fois), \all{dürfen}, \all{müssen}. Ces verbes expriment la \cle{permission} (\all{dürfen} = pouvoir, avoir le droit de) et l'\cle{obligation} (\all{müssen} = devoir). Ils sont au cœur du thème des libertés : une liberté, c'est précisément ce qu'on \emph{a le droit de} faire.

\textbf{2. Les connecteurs de concession.} Encadre : \all{obwohl} (bien que) et \all{trotzdem} (pourtant, néanmoins). Ils introduisent une opposition : les libertés sont importantes, \emph{mais} elles ont des limites ; elles ont des limites, \emph{mais} il faut quand même lutter pour elles.

\begin{center}
\begin{tabularx}{\linewidth}{|X|X|X|}\hline
\rowcolor{popblueL} \textbf{Mot allemand} & \textbf{Français} & \textbf{Fonction} \\ \hline
\all{dürfen} (darf, dürfen) & avoir le droit de, pouvoir & permission (verbe modal) \\ \hline
\all{müssen} & devoir & obligation (verbe modal) \\ \hline
\all{obwohl} + subordonnée & bien que & concession (verbe en fin) \\ \hline
\all{trotzdem} & pourtant, néanmoins & concession (adverbe, verbe 2e) \\ \hline
\end{tabularx}
\end{center}

\coursec{Grammaire 1 : Les verbes de modalité (Modalverben)}

\begin{propriete}[Conjugaison au présent de l'indicatif (Präsens)]
Les verbes de modalité sont \textbf{irréguliers} aux 3 personnes du singulier (changement de voyelle pour \all{können, dürfen, müssen, wollen, sollen}, pas pour \all{mögen}). L'infinitif principal se place \textbf{en fin de phrase}.
\begin{center}
\begin{tabular}{|l|c|c|c|c|c|c|}\hline
\rowcolor{popblueL} \textbf{Personne} & \textbf{dürfen} & \textbf{können} & \textbf{müssen} & \textbf{wollen} & \textbf{sollen} & \textbf{mögen} \\ \hline
ich & darf & kann & muss & will & soll & mag \\ \hline
du & darfst & kannst & musst & willst & sollst & magst \\ \hline
er/sie/es & \textbf{darf} & \textbf{kann} & \textbf{muss} & \textbf{will} & \textbf{soll} & mag \\ \hline
wir & dürfen & können & müssen & wollen & sollen & mögen \\ \hline
ihr & dürft & könnt & müsst & wollt & sollt & mögt \\ \hline
sie/Sie & dürfen & können & müssen & wollen & sollen & mögen \\ \hline
\end{tabular}
\end{center}
\textbf{Rappel :} pas de \all{-t} à la 3e pers. sg. pour \all{darf, kann, muss, will, soll, mag}. \all{mögen} n'est presque pas utilisé comme modal au présent (surtout \all{möchte} = Konjunktiv II).
\end{propriete}

\begin{propriete}[Sens et emploi des modaux (niveau Terminale)]
\begin{itemize}
\item \all{\textbf{dürfen}} : \textbf{permission / autorisation} (droit objectif). \all{Man darf hier rauchen.} (C'est permis). \textbf{Négation :} \all{nicht dürfen} = \textbf{interdiction} (\all{verboten}). \all{Man darf hier nicht rauchen.}
\item \all{\textbf{können}} : \textbf{capacité / possibilité} (physique, technique, circonstances). \all{Ich kann Deutsch sprechen.} / \all{Heute kann ich nicht kommen.}
\item \all{\textbf{müssen}} : \textbf{obligation / nécessité} (contrainte extérieure). \all{Ich muss lernen.} \textbf{Négation :} \all{nicht müssen} = \textbf{absence d'obligation} (pas \"interdiction\" !). \all{Du musst nicht kommen = Tu n'es pas obligé de venir.} (Interdiction = \all{du darfst nicht kommen}).
\item \all{\textbf{sollen}} : \textbf{obligation morale, conseil, ordre transmis}. \all{Du sollst deine Hausaufgaben machen.} (On attend de toi que...).
\item \all{\textbf{wollen}} : \textbf{volonté, intention}. \all{Ich will Arzt werden.}
\end{itemize}
\end{propriete}

\begin{exemplebox}[Structure de phrase avec modal (Präsens)]
\begin{tabular}{llll}
\all{Subjekt} & \textbf{Modal (conjugué, 2e pos.)} & \textbf{Compléments} & \textbf{Infinitif (fin)} \\ \hline
\all{Ich} & \all{darf} & \all{meine Meinung} & \all{sagen.} \\
\all{Die Presse} & \all{darf} & \all{über Politik} & \all{berichten.} \\
\all{Wir} & \all{müssen} & \all{für unsere Rechte} & \all{kämpfen.} \\
\all{Man} & \all{darf} & \all{niemanden} & \all{beleidigen.}
\end{tabular}
\end{exemplebox}

\begin{attention}[Pièges majeurs pour francophones]
\begin{enumerate}
\item \textbf{Place de l'infinitif :} En français, l'infinitif suit le modal (\emph{Je dois partir}). En allemand, l'infinitif va \textbf{à la toute fin de la phrase} (ou de la proposition), après tous les compléments. \all{Ich muss heute nach der Schule nach Hause gehen.}
\item \textbf{\all{nicht müssen} vs \all{nicht dürfen} :} \all{Du musst nicht} = « Tu n'as pas à / Tu n'es pas obligé de » (choix). \all{Du darfst nicht} = « Tu n'as pas le droit de / Il est interdit de » (interdiction). \textbf{Erreur fatale :} traduire « Tu ne dois pas fumer » par \all{Du musst nicht rauchen} (signifie : « Tu n'es pas obligé de fumer »). Correct : \all{Du darfst nicht rauchen.}
\item \textbf{\all{können} vs \all{dürfen} :} « Je peux venir » (capacité/possibilité) = \all{Ich kann kommen}. « Je peux venir » (permission/autorisation) = \all{Ich darf kommen}.
\end{enumerate}
\end{attention}

\coursec{Grammaire 2 : La concession (Konzession)}

\begin{propriete}[\all{obwohl} vs \all{trotzdem} : deux structures pour une même logique]
\begin{center}
\begin{tabularx}{\linewidth}{|X|X|}\hline
\rowcolor{popblueL} \textbf{\all{obwohl} (conjonction de subordination)} & \textbf{\all{trotzdem} (adverbe / conjonction de coordination)} \\ \hline
Introduit une \textbf{proposition subordonnée}. & Relie deux \textbf{propositions principales} (ou commence une phrase). \\
Le verbe conjugué va \textbf{à la fin} de la subordonnée. & Le verbe conjugué reste en \textbf{2e position} (sujet après si \all{trotzdem} est en tête). \\
\all{Obwohl es regnet, gehen wir spazieren.} & \all{Es regnet. Trotzdem gehen wir spazieren.} \\
(Bien qu'il pleuve, nous nous promenons.) & (Il pleut. Pourtant, nous nous promenons.) \\ \hline
\end{tabularx}
\end{center}
\textbf{Ordre des mots avec \all{trotzdem} en tête :} \all{Trotzdem} (Pos 1) + \all{Verbe} (Pos 2) + \all{Sujet} + ... \all{Trotzdem \textbf{gehen} \textbf{wir} spazieren.}
\end{propriete}

\begin{exemplebox}[Exemples de concession (traduits)]
\begin{itemize}
\item \all{Obwohl die Freiheiten wichtig sind, haben sie Grenzen.} « Bien que les libertés soient importantes, elles ont des limites. » (Verbe \all{sind} à la fin).
\item \all{Die Freiheiten sind wichtig. Trotzdem haben sie Grenzen.} « Les libertés sont importantes. Pourtant, elles ont des limites. » (Verbe \all{haben} en 2e).
\item \all{Er ist jung, trotzdem darf er wählen.} « Il est jeune, pourtant il a le droit de voter. »
\item \all{Obwohl er müde ist, lernt er weiter.} « Bien qu'il soit fatigué, il continue d'étudier. »
\end{itemize}
\end{exemplebox}

\begin{attention}[Erreurs fréquentes : place du verbe]
\begin{enumerate}
\item Ne pas mettre le verbe en 2e position après \all{obwohl}. \textbf{Faux :} \all{Obwohl die Freiheiten sind wichtig...} $\rightarrow$ \textbf{Correct :} \all{Obwohl die Freiheiten wichtig \textbf{sind}...}
\item Oublier l'inversion sujet/verbe après \all{trotzdem} en début de phrase. \textbf{Faux :} \all{Trotzdem wir haben Rechte.} $\rightarrow$ \textbf{Correct :} \all{Trotzdem \textbf{haben} \textbf{wir} Rechte.}
\item Confondre \all{trotzdem} (pourtant) et \all{trotz} (malgré + génitif). \all{Trotz des Regens gehen wir raus.} (Malgré la pluie...).
\end{enumerate}
\end{attention}

\coursec{Traduction guidée phrase par phrase (Version)}

\begin{methode}[Traduire un texte allemand en version]
\begin{enumerate}
\item Découpe la phrase en groupes de sens (sujet / verbe conjugué / compléments / infinitif ou participe final).
\item Repère le verbe conjugué (2e position en principale, fin en subordonnée \all{obwohl}) et l'infinitif (fin de phrase).
\item Traduis groupe par groupe, puis reformule en français naturel (attention aux faux amis \all{Recht/Grenze}).
\item Vérifie : aucun mot allemand ne doit être oublié ; le sens logique (concession, modalité) doit être rendu.
\end{enumerate}
\end{methode}

\begin{exemplebox}[Traduction commentée du texte support]
\begin{itemize}
\item \all{In einer Demokratie haben die Bürger wichtige Freiheiten.} « Dans une démocratie, les citoyens ont des libertés importantes. » (\all{haben} 2e, sujet \all{die Bürger} postposé).
\item \all{Jeder darf seine Meinung frei sagen.} « Chacun a le droit d'exprimer librement son opinion. » (\all{darf} modal 2e, \all{sagen} infinitif fin).
\item \all{Die Presse darf über Politik berichten, ohne Angst zu haben.} « La presse peut rendre compte de la politique sans avoir peur. » (\all{berichten} infinitif fin ; \all{ohne ... zu} infinitif).
\item \all{Außerdem hat jeder das Recht, seine Religion frei zu wählen.} « De plus, chacun a le droit de choisir librement sa religion. » (\all{hat} 2e ; infinitif \all{zu wählen} fin).
\item \all{Die Bürger dürfen sich friedlich versammeln.} « Les citoyens ont le droit de se rassembler pacifiquement. » (\all{dürfen} modal pluriel ; \all{sich versammeln} réfléchi).
\item \all{Obwohl diese Freiheiten wichtig sind, haben sie auch Grenzen.} « Bien que ces libertés soient importantes, elles ont aussi des limites. » (\all{obwohl} $\rightarrow$ verbe \all{sind} fin ; principale \all{haben} 2e).
\item \all{Man darf zum Beispiel niemanden beleidigen.} « On n'a pas le droit, par exemple, d'insulter quelqu'un. » (\all{niemanden} accusatif ; \all{darf nicht} = interdiction).
\item \all{Trotzdem müssen wir für unsere Rechte kämpfen, denn Freiheit ist ein kostbares Gut.} « Néanmoins, nous devons lutter pour nos droits, car la liberté est un bien précieux. » (\all{trotzdem} $\rightarrow$ inversion \all{müssen wir} ; \all{denn} $\rightarrow$ pas d'inversion).
\end{itemize}
\end{exemplebox}

\begin{exemplebox}[Exemples supplémentaires traduits (Thème / Version)]
\all{Jeder darf seine Meinung sagen.} « Chacun a le droit d'exprimer son opinion. »

\all{Die Presse muss die Wahrheit berichten.} « La presse doit rapporter la vérité. »

\all{Obwohl er jung ist, darf er wählen.} « Bien qu'il soit jeune, il a le droit de voter. »

\all{Wir dürfen demonstrieren, trotzdem brauchen wir eine Erlaubnis.} « Nous avons le droit de manifester, mais nous avons besoin d'une autorisation. »

\all{Man darf nicht lügen.} « Il est interdit de mentir / On n'a pas le droit de mentir. » (\all{nicht dürfen} = interdiction).

\all{Du musst nicht kommen.} « Tu n'es pas obligé de venir. » (\all{nicht müssen} = absence d'obligation).
\end{exemplebox}

\begin{attention}[Freiheit ou Recht ?]
Ne confonds pas \all{die Freiheit} (la liberté, état, concept) et \all{das Recht} (le droit, garantie juridique). La \emph{liberté} est l'état ; le \emph{droit} est ce que la loi garantit. On dit : \all{das Recht auf Meinungsfreiheit} « le droit à la liberté d'expression ». Autre piège : \all{das Recht} ne signifie jamais « correct » (c'est \all{richtig}) ni « à droite » (c'est \all{rechts}).
\end{attention}

\coursec{Compréhension du texte et Commentaire}

\courssub{Questions de compréhension globale et détaillée}

\begin{exoresolu}[Questions de compréhension (Antworten auf Deutsch)]
\textbf{Énoncé.} Réponds en allemand, par des phrases complètes :
1. \all{Was ist Meinungsfreiheit?}
2. \all{Welche vier Freiheiten nennt der Text?}
3. \all{Haben die Freiheiten Grenzen? Nenne ein Beispiel.}
4. \all{Was müssen die Bürger trotzdem tun?}
5. \all{Warum ist Freiheit ein "kostbares Gut"? (Interprétation libre)}
\tcblower
\textbf{Solution détaillée.}
1. \all{Meinungsfreiheit ist das Recht, seine Meinung frei zu sagen oder zu schreiben.}
2. \all{Der Text nennt Meinungsfreiheit, Pressefreiheit, Religionsfreiheit und Versammlungsfreiheit.}
3. \all{Ja, sie haben Grenzen. Zum Beispiel darf man niemanden beleidigen.}
4. \all{Wir müssen für unsere Rechte kämpfen.}
5. \all{Freiheit ist ein kostbares Gut, weil sie nicht selbstverständlich ist und man sie schützen muss.}
\textbf{Méthode :} Reprends les mots de la question (\all{Was} $\rightarrow$ \all{Das ist...} ; \all{Welche} $\rightarrow$ \all{Der Text nennt...}). Verbe conjugué en 2e position.
\end{exoresolu}

\courssub{Méthode du commentaire de texte (niveau Terminale)}

\begin{methode}[Comment rédiger un commentaire structuré (en allemand ou français)]
\begin{enumerate}
\item \textbf{Introduction} (3-4 lignes) : Présenter le document (nature, titre, source, thème). Annoncer le plan (ex. : 1. Les libertés énumérées, 2. Leurs limites, 3. L'engagement citoyen).
\item \textbf{Développement} (2-3 paragraphes) : Un argument par paragraphe. Citer le texte (\all{Im Text heißt es: „...“}) et expliquer. Utiliser les connecteurs logiques (\all{zuerst, außerdem, jedoch, schließlich}).
\item \textbf{Conclusion} (2-3 lignes) : Résumer l'idée principale. Ouvrir sur le contexte actuel (Cameroun, Allemagne, réseaux sociaux) ou donner son avis personnel (\all{Meiner Meinung nach...}).
\end{enumerate}
\end{methode}

\begin{exemplebox}[Commentaire modèle (en allemand, niveau B1 visé)]
\textbf{Thema: Die Freiheiten in einer Demokratie}

\textbf{Einleitung:} Der Zeitungsartikel „Freiheiten in einer Demokratie“ thematisiert die Grundrechte der Bürger in einem demokratischen Staat. Er nennt vier wichtige Freiheiten und zeigt, dass diese nicht grenzenlos sind.

\textbf{Hauptteil:} \textbf{Zuerst} werden die vier Freiheiten vorgestellt: die Meinungsfreiheit, die Pressefreiheit, die Religionsfreiheit und die Versammlungsfreiheit. Der Text erklärt kurz, was jede Freiheit bedeutet: Man darf seine Meinung sagen, die Presse darf berichten, man kann seine Religion wählen und sich friedlich treffen. \textbf{Allerdings} (oder: \textbf{Obwohl}) diese Rechte wichtig sind, \textbf{haben sie Grenzen}. Als Beispiel nennt der Text: Man darf niemanden beleidigen. Das zeigt: Freiheit endet dort, wo die Rechte anderer beginnen. \textbf{Schließlich} appelliert der Text an die Bürger: \textbf{Trotzdem müssen wir für unsere Rechte kämpfen}, denn Freiheit ist ein kostbares Gut.

\textbf{Schluss:} Meiner Meinung nach ist der Text sehr wichtig für junge Leute. In Kamerun wie in Deutschland müssen wir wissen, dass Rechte Pflichten beinhalten. Man darf kritisieren, aber man muss respektvoll bleiben.
\end{exemplebox}

\coursec{Production guidée écrite et orale}

\courssub{Expression écrite : Résumé et prise de position}

\begin{experience}[Production écrite guidée : « Mein Recht – Meine Pflicht »]
\textbf{Consigne :} Rédige un paragraphe d'environ 80-100 mots en allemand sur le sujet suivant :
„In der Schule wollen Schüler eine Zeitung gründen. Der Direktor verbietet einen kritischen Artikel. Dürfen die Schüler protestieren? Was sind ihre Rechte und Pflichten?“

\textbf{Aide au rédaction (Bausteine) :}
\begin{itemize}
\item \all{Die Schüler wollen eine Schülerzeitung gründen.}
\item \all{Der Direktor verbietet einen Artikel. / Der Direktor darf das nicht, weil...}
\item \all{Die Schüler haben das Recht auf Meinungsfreiheit (Artikel 5 GG).}
\item \all{Aber sie haben auch Pflichten: sie dürfen nicht lügen / beleidigen.}
\item \all{Meiner Meinung nach sollen die Schüler mit dem Direktor sprechen.}
\item \all{Man muss für seine Rechte kämpfen, aber friedlich.}
\end{itemize}
\textbf{Critères de réussite :} Utiliser au moins 3 verbes de modalité (\all{dürfen, müssen, sollen, können}) et un connecteur de concession (\all{obwohl} ou \all{trotzdem}). Respecter l'ordre des mots (verbe 2e / infinitif fin / verbe fin en subordonnée).
\end{experience}

\begin{exemplebox}[Modèle de production écrite (Corrigé type)]
\all{Die Schüler wollen eine Schülerzeitung gründen. Der Direktor verbietet einen kritischen Artikel. Obwohl der Direktor die Verantwortung hat, dürfen die Schüler ihre Meinung sagen. Die Pressefreiheit und die Meinungsfreiheit sind Grundrechte. Man darf aber niemanden beleidigen oder lügen. Die Schüler haben die Pflicht, fair zu bleiben. Trotzdem müssen sie für ihr Recht kämpfen. Meiner Meinung nach sollen sie das Gespräch mit dem Direktor suchen. Man muss friedlich für seine Rechte kämpfen, denn Freiheit ist ein kostbares Gut.}
\end{exemplebox}

\courssub{Expression orale : Débat et jeu de rôle}

\begin{experience}[Jeu de rôle : « Schulleiter vs. Schülerzeitung »]
\textbf{Situation :} Un groupe de 3 élèves (rédacteurs) rencontre le directeur (un élève) pour négocier la publication de l'article critique sur la cantine.
\textbf{Rôles :}
\begin{itemize}
\item \textbf{Élèves (2-3) :} Argumentent avec \all{dürfen, müssen, sollen}. Ex: \all{Wir dürfen unsere Meinung sagen! Wir sollen die Wahrheit schreiben!}
\item \textbf{Directeur :} Argumente avec \all{müssen, nicht dürfen, Grenzen}. Ex: \all{Ihr dürft nicht beleidigen. Ihr müsst die Schulordnung respektieren.}
\item \textbf{Observateur :} Note les arguments, vérifie la grammaire (place du verbe, modaux), donne un feedback.
\end{itemize}
\textbf{Durée :} 5 min préparation, 5 min jeu, 3 min feedback. Langue cible : allemand uniquement.
\end{experience}

\begin{exercice}[Entraînement final (Thème / Grammaire / Vocabulaire)]
\begin{enumerate}
\item \textbf{Vocabulaire :} Complète le tableau avec les mots de la même famille.
\begin{center}
\begin{tabular}{|l|l|l|}\hline
\rowcolor{popblueL} \textbf{Verbe} & \textbf{Nom (art. + pl.)} & \textbf{Adjectif} \\ \hline
\all{beleidigen} & \all{die Beleidigung (-en)} & \all{beleidigend} \\ \hline
\all{versammeln} & \ldots & \ldots \\ \hline
\ldots & \all{die Erlaubnis (-se)} & \all{erlaubt} \\ \hline
\all{kämpfen} & \ldots & \ldots \\ \hline
\ldots & \all{die Pflicht (-en)} & \all{pflichtbewusst} \\ \hline
\end{tabular}
\end{center}
\item \textbf{Grammaire (Modaux) :} Conjugue le verbe entre parenthèses au présent et place l'infinitif correctement.
\begin{enumerate}
\item Die Schüler \ldots{} (dürfen) die Zeitung \ldots{} (lesen / herausgeben).
\item Der Direktor \ldots{} (müssen) die Regeln \ldots{} (erklären).
\item Wir \ldots{} (können) heute nicht \ldots{} (kommen), weil wir \ldots{} (müssen) lernen.
\item \ldots{} (dürfen) ich hier \ldots{} (rauchen)? – Nein, man \ldots{} (nicht dürfen) hier \ldots{} (rauchen).
\end{enumerate}
\item \textbf{Grammaire (Concession) :} Relie les deux phrases par \all{obwohl} (subordonnée) PUIS par \all{trotzdem} (principale).
\begin{enumerate}
\item Die Freiheiten sind wichtig. Sie haben Grenzen.
\item Er ist erst 16 Jahre alt. Er darf wählen.
\item Wir haben Angst. Wir demonstrieren friedlich.
\end{enumerate}
\item \textbf{Thème (Français $\rightarrow$ Allemand) :} Traduis en allemand.
\begin{enumerate}
\item Dans une démocratie, les citoyens ont des droits et des devoirs.
\item Bien que la liberté soit importante, elle a des limites.
\item On n'a pas le droit d'insulter qui que ce soit.
\item Néanmoins, nous devons lutter pour la liberté de la presse.
\item À mon avis, la liberté est un bien précieux.
\end{enumerate}
\end{enumerate}
\end{exercice}

\begin{corrige}[Exercice : Entraînement final]
\textbf{1. Vocabulaire :}
\all{versammeln} – \all{die Versammlung (-en)} – \all{versammelt} ; \all{erlauben} – \all{die Erlaubnis (-se)} – \all{erlaubt} ; \all{kämpfen} – \all{der Kampf (\"e)} – \all{kämpferisch} ; \all{pflichen} (verbe rare) – \all{die Pflicht (-en)} – \all{pflichtbewusst}.

\textbf{2. Modaux :}
1. \all{Die Schüler dürfen die Zeitung herausgeben.} (ou \all{lesen}) – 2. \all{Der Direktor muss die Regeln erklären.} – 3. \all{Wir können heute nicht kommen, weil wir lernen müssen.} (subordonnée \all{weil} $\rightarrow$ verbe fin) – 4. \all{Darf ich hier rauchen? – Nein, man darf hier nicht rauchen.}

\textbf{3. Concession :}
1. \all{Obwohl die Freiheiten wichtig sind, haben sie Grenzen.} / \all{Die Freiheiten sind wichtig. Trotzdem haben sie Grenzen.}
2. \all{Obwohl er erst 16 Jahre alt ist, darf er wählen.} / \all{Er ist erst 16 Jahre alt. Trotzdem darf er wählen.}
3. \all{Obwohl wir Angst haben, demonstrieren wir friedlich.} / \all{Wir haben Angst. Trotzdem demonstrieren wir friedlich.}

\textbf{4. Thème :}
1. \all{In einer Demokratie haben die Bürger Rechte und Pflichten.}
2. \all{Obwohl die Freiheit wichtig ist, hat sie Grenzen.} (ou \all{Obwohl die Freiheit wichtig ist, gibt es Grenzen.})
3. \all{Man darf niemanden beleidigen.}
4. \all{Trotzdem müssen wir für die Pressefreiheit kämpfen.}
5. \all{Meiner Meinung nach ist die Freiheit ein kostbares Gut.}
\end{corrige}

\begin{savaistu}[Le Grundgesetz de 1949 et le contexte camerounais]
En Allemagne, les libertés fondamentales sont protégées par la Loi fondamentale (\all{das Grundgesetz}) de 1949. L'article 5 protège la liberté d'expression et de la presse, l'article 4 la liberté de religion, l'article 8 la liberté de réunion. Ces articles ont été écrits après la dictature nazie : les rédacteurs voulaient que jamais plus les citoyens ne soient réduits au silence.

Au Cameroun, la Constitution de 1996 (révisée en 2008) garantit également les libertés d'expression, de presse, de conscience et de réunion (Préambule et articles). Cependant, l'exercice de ces libertés est encadré par des lois (ex. loi sur la cybercriminalité, loi sur la presse) qui posent parfois débat sur l'équilibre entre sécurité et liberté. Comparer les deux contextes permet de mieux comprendre la notion de \all{Verhältnismäßigkeit} (proportionnalité) des limites.
\end{savaistu}

\begin{aretenir}[L'essentiel de la leçon AL25]
\begin{itemize}
\item \textbf{4 libertés fondamentales :} \all{die Meinungsfreiheit, die Pressefreiheit, die Religionsfreiheit, die Versammlungsfreiheit} (toutes féminines, composées).
\item \textbf{Verbes de modalité (Präsens) :} \all{dürfen} (permission / interdiction avec \all{nicht}), \all{müssen} (obligation / absence d'obligation avec \all{nicht}), \all{können} (capacité), \all{sollen} (devoir moral), \all{wollen} (volonté). \cle{Infinitif en fin de phrase.}
\item \textbf{Concession :} \all{obwohl} + subordonnée (verbe \textbf{à la fin}) vs \all{trotzdem} + principale (verbe \textbf{en 2e position}, inversion si \all{trotzdem} en tête).
\item \textbf{Vocabulaire clé :} \all{das Recht} (le droit), \all{die Pflicht} (le devoir), \all{die Grenze} (la limite), \all{verbieten} (interdire), \all{erlauben} (permettre), \all{kämpfen für} (lutter pour).
\item \textbf{Compétence Bac :} Comprendre un texte d'actualité, répondre en allemand, résumer, exprimer son avis avec modaux et concession, traduire thème/version.
\end{itemize}
\end{aretenir}

\coursec{Compréhension du texte}

Dans la section précédente, nous avons lu et écouté le texte \all{„Freiheiten in einer Demokratie“}, qui présente les grandes libertés — \all{die Meinungsfreiheit}, \all{die Pressefreiheit}, \all{die Religionsfreiheit} et \all{die Versammlungsfreiheit} — ainsi que leurs limites. Nous allons maintenant vérifier et approfondir notre compréhension : d'abord de manière \cle{globale} (de quoi parle le texte ?), puis de manière \cle{détaillée} (questions précises), avant d'apprendre la \cle{méthode} pour répondre correctement à une question de compréhension au baccalauréat.

\courssub{Rappel : le texte support}

\begin{textebox}[Freiheiten in einer Demokratie (rappel)]
In einer Demokratie haben alle Bürger Rechte und Freiheiten. Die wichtigste Freiheit ist die Meinungsfreiheit: Jeder Mensch darf seine Meinung frei sagen, schreiben oder im Internet veröffentlichen. Niemand darf wegen seiner Meinung ins Gefängnis kommen.

Auch die Pressefreiheit ist sehr wichtig. Journalisten dürfen über Politik, Wirtschaft und Gesellschaft berichten, ohne Angst zu haben. Eine freie Presse kontrolliert die Regierung und informiert die Bürger. Ohne Pressefreiheit gibt es keine echte Demokratie.

Die Religionsfreiheit bedeutet: Jeder darf an Gott glauben oder nicht glauben. Christen, Muslime und Anhänger anderer Religionen dürfen ihre Religion frei praktizieren, in der Kirche, in der Moschee oder zu Hause.

Schließlich gibt es die Versammlungsfreiheit: Die Bürger dürfen sich friedlich versammeln, zum Beispiel bei einer Demonstration oder bei einem Treffen im Jugendzentrum.

Aber die Freiheiten haben Grenzen. Man darf andere Menschen nicht beleidigen, nicht bedrohen und nicht zu Gewalt aufrufen. Meine Freiheit endet dort, wo die Freiheit des anderen beginnt.

Freiheit ist ein Geschenk, aber auch eine Aufgabe: Die Bürger müssen ihre Rechte kennen, sie benutzen und sie verteidigen. Denn Freiheiten sind nicht für immer garantiert — man muss jeden Tag dafür kämpfen, in Deutschland, in Österreich, in der Schweiz und überall auf der Welt.
\end{textebox}

\textbf{Glossaire}

\begin{center}
\begin{tabularx}{\linewidth}{|X|X|}\hline
\rowcolor{popblueL} \textbf{Deutsch} & \textbf{Français} \\ \hline
der Bürger, die Bürger & le citoyen \\ \hline
die Meinungsfreiheit, -en & la liberté d'expression \\ \hline
veröffentlichen (hat veröffentlicht) & publier \\ \hline
das Gefängnis, -se & la prison \\ \hline
die Pressefreiheit, -en & la liberté de la presse \\ \hline
berichten (hat berichtet) & rapporter, informer \\ \hline
die Religionsfreiheit, -en & la liberté de religion \\ \hline
praktizieren (hat praktiziert) & pratiquer \\ \hline
die Versammlungsfreiheit, -en & la liberté de réunion \\ \hline
sich versammeln (hat sich versammelt) & se réunir \\ \hline
die Grenze, -n & la limite, la frontière \\ \hline
beleidigen (hat beleidigt) & insulter, offenser \\ \hline
bedrohen (hat bedroht) & menacer \\ \hline
verteidigen (hat verteidigt) & défendre \\ \hline
garantieren (hat garantiert) & garantir \\ \hline
\end{tabularx}
\end{center}

\courssub{Compréhension globale : de quoi parle le texte ?}

Avant de répondre aux questions détaillées, il faut saisir l'\cle{idée générale} du texte. Pose-toi trois questions simples : \emph{de quoi parle le texte ?} (le thème), \emph{qui est concerné ?} (les personnes), \emph{où ?} (le lieu ou le contexte).

\begin{exemplebox}[Compréhension globale — réponses modèles]
\textbf{Worum geht es im Text?} (De quoi parle le texte ?)

\all{Der Text handelt von den Freiheiten in einer Demokratie: Meinungsfreiheit, Pressefreiheit, Religionsfreiheit und Versammlungsfreiheit.} « Le texte traite des libertés dans une démocratie : la liberté d'expression, la liberté de la presse, la liberté de religion et la liberté de réunion. »

\textbf{Wer hat Rechte und Freiheiten?} (Qui a des droits et des libertés ?)

\all{Alle Bürger haben Rechte und Freiheiten.} « Tous les citoyens ont des droits et des libertés. »

\textbf{Wo gelten diese Freiheiten?} (Où ces libertés sont-elles valables ?)

\all{Sie gelten in einer Demokratie, zum Beispiel in Deutschland, in Österreich und in der Schweiz.} « Elles s'appliquent dans une démocratie, par exemple en Allemagne, en Autriche et en Suisse. »
\end{exemplebox}

\begin{exercice}[Richtig oder falsch ? — Compréhension globale]
Lis le texte et coche : \all{richtig} (vrai) ou \all{falsch} (faux).
\begin{enumerate}
\item \all{Der Text handelt von den Freiheiten in einer Demokratie.}
\item \all{Nur Journalisten haben Rechte.}
\item \all{Die Meinungsfreiheit ist die wichtigste Freiheit im Text.}
\item \all{Die Freiheiten haben keine Grenzen.}
\item \all{Man muss für die Freiheiten kämpfen.}
\end{enumerate}
\end{exercice}

\begin{corrige}[Exercice : Richtig oder falsch ?]
\begin{enumerate}
\item \all{richtig} — première phrase du texte : \all{„In einer Demokratie haben alle Bürger Rechte und Freiheiten.“}
\item \all{falsch} — \all{„alle Bürger“} ont des droits, pas seulement les journalistes.
\item \all{richtig} — \all{„Die wichtigste Freiheit ist die Meinungsfreiheit.“}
\item \all{falsch} — \all{„Aber die Freiheiten haben Grenzen.“}
\item \all{richtig} — dernière phrase : \all{„man muss jeden Tag dafür kämpfen“}.
\end{enumerate}
\end{corrige}

\courssub{Compréhension détaillée : les questions sur le texte}

Passons maintenant aux questions de détail. Au baccalauréat, ces questions portent sur un passage précis : il faut localiser le passage, puis répondre avec une phrase complète.

\begin{exemplebox}[Questions détaillées — réponses modèles]
\textbf{1. Welche Freiheiten werden im Text genannt?} (Quelles libertés sont mentionnées dans le texte ?)

\all{Im Text werden vier Freiheiten genannt: die Meinungsfreiheit, die Pressefreiheit, die Religionsfreiheit und die Versammlungsfreiheit.} « Quatre libertés sont mentionnées : la liberté d'expression, la liberté de la presse, la liberté de religion et la liberté de réunion. »

\textbf{2. Was bedeutet Pressefreiheit?} (Que signifie la liberté de la presse ?)

\all{Pressefreiheit bedeutet, dass Journalisten über Politik, Wirtschaft und Gesellschaft berichten dürfen, ohne Angst zu haben.} « La liberté de la presse signifie que les journalistes peuvent informer sur la politique, l'économie et la société sans avoir peur. »

\textbf{3. Haben die Freiheiten Grenzen? Welche?} (Les libertés ont-elles des limites ? Lesquelles ?)

\all{Ja, die Freiheiten haben Grenzen: Man darf andere Menschen nicht beleidigen, nicht bedrohen und nicht zu Gewalt aufrufen.} « Oui, les libertés ont des limites : on ne doit pas insulter les autres, ni les menacer, ni appeler à la violence. »

\textbf{4. Warum müssen die Bürger für ihre Rechte kämpfen?} (Pourquoi les citoyens doivent-ils lutter pour leurs droits ?)

\all{Die Bürger müssen für ihre Rechte kämpfen, weil Freiheiten nicht für immer garantiert sind.} « Les citoyens doivent lutter pour leurs droits parce que les libertés ne sont pas garanties pour toujours. »
\end{exemplebox}

Observe bien la structure des réponses : chacune \cle{reprend les mots de la question} (\all{Pressefreiheit bedeutet...}, \all{Die Bürger müssen... kämpfen, weil...}) et s'appuie sur le vocabulaire du texte. C'est exactement ce qu'attend le correcteur.

\courssub{Méthode : répondre à une question de compréhension}

\begin{methode}[Répondre à une question de compréhension en 4 étapes]
\begin{enumerate}
\item \textbf{Souligner les mots-clés de la question} : le mot interrogatif (\all{wer}, \all{was}, \all{warum}, \all{welche}, \all{wo}) et le nom ou le verbe important (ex. \all{Grenzen}, \all{kämpfen}).
\item \textbf{Retrouver le passage correspondant} dans le texte : cherche le même mot ou un mot de la même famille (ex. question : \all{Grenzen} → texte : \all{„die Freiheiten haben Grenzen“}).
\item \textbf{Reformuler} : réponds avec une phrase complète en reprenant le vocabulaire du texte, \emph{sans recopier tout le paragraphe}.
\item \textbf{Rédiger et vérifier} : majuscule au début, verbe conjugué en deuxième position, point final. Si possible, justifie par une courte citation entre guillemets.
\end{enumerate}
\end{methode}

\begin{popfigure}
\begin{tikzpicture}[
node distance=0.55cm,
etape/.style={rectangle, rounded corners=4pt, draw=popblue, fill=popblueL, text width=6.2cm, align=center, font=\small, inner sep=6pt},
fleche/.style={-{Stealth[length=3mm]}, thick, popdark}]
\node[etape] (e1) {\textbf{1. Lire la question}\newline souligner les mots-clés (\all{wer, was, warum...})};
\node[etape, below=of e1] (e2) {\textbf{2. Repérer le passage}\newline chercher le mot-clé ou sa famille dans le texte};
\node[etape, below=of e2] (e3) {\textbf{3. Reformuler}\newline reprendre les mots de la question + vocabulaire du texte};
\node[etape, below=of e3, draw=popgreen, fill=popgreenL] (e4) {\textbf{4. Rédiger une phrase complète}\newline majuscule, verbe en position 2, point final};
\draw[fleche] (e1) -- (e2);
\draw[fleche] (e2) -- (e3);
\draw[fleche] (e3) -- (e4);
\end{tikzpicture}
\legende{Organigramme de la méthode de réponse à une question de compréhension}
\end{popfigure}

\begin{exemplebox}[Application de la méthode, étape par étape]
\textbf{Frage :} \all{Was dürfen die Bürger?} (Qu'ont le droit de faire les citoyens ?)

\textbf{Étape 1} — mots-clés : \all{was} + \all{dürfen} + \all{die Bürger}.

\textbf{Étape 2} — passage du texte : \all{„Jeder Mensch darf seine Meinung frei sagen“} et \all{„Die Bürger dürfen sich friedlich versammeln“}.

\textbf{Étapes 3 et 4} — réponse rédigée :

\all{Die Bürger dürfen ihre Meinung frei sagen und sich friedlich versammeln.} « Les citoyens ont le droit d'exprimer librement leur opinion et de se réunir pacifiquement. »
\end{exemplebox}

\begin{attention}[Les erreurs qui coûtent des points au Bac]
\begin{itemize}
\item \textbf{Ne réponds jamais seulement par} \all{ja} \textbf{ou} \all{nein}. À la question \all{„Haben die Freiheiten Grenzen?“}, il faut écrire : \all{Ja, die Freiheiten haben Grenzen: man darf niemanden beleidigen.} — la réponse complète montre que tu as compris le texte.
\item \textbf{Majuscule et sujet obligatoires} : une réponse allemande commence par une majuscule et contient un sujet et un verbe conjugué. N'écris pas \all{„weil Freiheiten nicht garantiert sind“} seul : c'est une subordonnée incomplète. Écris : \all{Die Bürger müssen kämpfen, weil Freiheiten nicht für immer garantiert sind.}
\item \textbf{Ne recopie pas tout le paragraphe} : le correcteur veut voir que tu sais \emph{sélectionner} l'information. Une ou deux phrases suffisent.
\item \textbf{Attention au verbe en position 2} : pas de calque du français. On écrit \all{Im Text werden vier Freiheiten genannt} (le verbe \all{werden} reste en deuxième position après le complément \all{im Text}).
\end{itemize}
\end{attention}

\courssub{Vrai ou faux justifié : entraînement}

L'exercice « vrai ou faux \emph{justifié} » est plus exigeant : il ne suffit pas de dire \all{falsch}, il faut \cle{corriger} la phrase et citer le texte.

\begin{exoresolu}[Vrai ou faux justifié]
\textbf{Énoncé.} Juge la phrase suivante et justifie ta réponse par une citation courte du texte :

\all{„Man darf andere Leute beleidigen.“}

\tcblower
\textbf{Solution détaillée.}

\textbf{Étape 1} — mot-clé : \all{beleidigen} (insulter). Je repère le passage : \all{„Man darf andere Menschen nicht beleidigen“}.

\textbf{Étape 2} — la phrase de l'énoncé dit le \emph{contraire} du texte : elle est donc \all{falsch}.

\textbf{Étape 3} — je rédige la réponse complète avec la correction et la justification :

\all{Das ist falsch. Man darf niemanden beleidigen. Im Text steht: „Man darf andere Menschen nicht beleidigen, nicht bedrohen und nicht zu Gewalt aufrufen.“}

« C'est faux. On ne doit insulter personne. Dans le texte, il est écrit : "On ne doit pas insulter les autres, ni les menacer, ni appeler à la violence." »

Remarque la formulation \all{Im Text steht: „...“} (« dans le texte, il est écrit... ») : c'est la manière standard de citer un texte en allemand.
\end{exoresolu}

\begin{exercice}[À ton tour : vrai ou faux justifié]
Juge les phrases suivantes (\all{richtig} ou \all{falsch}) et justifie par une citation courte.
\begin{enumerate}
\item \all{Journalisten dürfen nicht über Politik berichten.}
\item \all{Jeder darf an Gott glauben oder nicht glauben.}
\item \all{Freiheiten sind für immer garantiert.}
\end{enumerate}
\end{exercice}

\begin{corrige}[Exercice : vrai ou faux justifié]
\begin{enumerate}
\item \all{Falsch. Journalisten dürfen über Politik, Wirtschaft und Gesellschaft berichten, ohne Angst zu haben.}
\item \all{Richtig. Im Text steht: „Jeder darf an Gott glauben oder nicht glauben.“}
\item \all{Falsch. Freiheiten sind nicht für immer garantiert — man muss jeden Tag dafür kämpfen.}
\end{enumerate}
\end{corrige}

\courssub{Recherche dans le texte : les familles de mots}

Un bon lecteur sait retrouver dans un texte les mots de la même \cle{famille}. Cela aide à comprendre le vocabulaire nouveau et à repérer rapidement les passages utiles. En allemand, les familles se construisent souvent par \cle{composition} (deux noms collés) : \all{die Meinung} + \all{die Freiheit} → \all{die Meinungsfreiheit}.

\begin{center}
\begin{tabularx}{\linewidth}{|X|X|X|}\hline
\rowcolor{popblueL} \textbf{Mot de base} & \textbf{Mot composé du texte} & \textbf{Français} \\ \hline
frei (libre) & die Freiheit, -en & la liberté \\ \hline
die Meinung, -en (opinion) & die Meinungsfreiheit & la liberté d'expression \\ \hline
die Presse, -n (la presse) & die Pressefreiheit & la liberté de la presse \\ \hline
die Religion, -en & die Religionsfreiheit & la liberté de religion \\ \hline
die Versammlung, -en (réunion) & die Versammlungsfreiheit & la liberté de réunion \\ \hline
das Recht, -e (le droit) & das Bürgerrecht, -e & le droit du citoyen / les droits du citoyen \\ \hline
der Bürger, die Bürger (citoyen) & — & les citoyens (pluriel) \\ \hline
\end{tabularx}
\end{center}

\begin{aretenir}[Le genre du mot composé]
En allemand, un nom composé (\cle{Determinativkompositum}) prend \textbf{toujours le genre et le pluriel du dernier élément} (le \cle{noyau}) : \all{die Meinung} + \all{die Freiheit} → \all{die Meinungsfreiheit} (féminin, comme \all{Freiheit}). C'est une règle sûre, sans exception : \all{der Bürger} + \all{das Recht} → \all{das Bürgerrecht} (neutre, comme \all{Recht}).
\end{aretenir}

\courssub{Questions d'ouverture : et au Cameroun ?}

Pour préparer l'expression orale (section 7), le professeur te posera des questions d'\cle{ouverture} qui te font passer du texte à ta propre expérience. Voici des modèles de réponses courtes.

\begin{exemplebox}[Und in Kamerun? — réponses possibles]
\textbf{Haben wir in Kamerun dieselben Freiheiten?} (Avons-nous les mêmes libertés au Cameroun ?)

\all{Ja, in Kamerun haben wir auch Meinungsfreiheit und Religionsfreiheit. In Yaoundé und Douala gibt es viele Kirchen und Moscheen, und die Menschen praktizieren ihre Religion frei.} « Oui, au Cameroun nous avons aussi la liberté d'expression et la liberté de religion. À Yaoundé et à Douala, il y a beaucoup d'églises et de mosquées, et les gens pratiquent librement leur religion. »

\textbf{Warum ist Pressefreiheit wichtig für dein Land?} (Pourquoi la liberté de la presse est-elle importante pour ton pays ?)

\all{Pressefreiheit ist wichtig, weil eine freie Presse die Bürger informiert und die Regierung kontrolliert.} « La liberté de la presse est importante parce qu'une presse libre informe les citoyens et contrôle le gouvernement. »
\end{exemplebox}

\begin{attention}[Faux amis : \all{die Presse} et \all{die Demonstration}]
\begin{itemize}
\item \all{Die Presse} ne signifie pas « être pressé » ! \all{Die Presse} = les journaux, les médias (la presse). Pour dire « je suis pressé », on dit \all{ich habe es eilig} ou \all{ich bin in Eile}.
\item \all{Die Demonstration} = la manifestation (pacifique ou revendicative), pas une « démonstration » mathématique ou logique (ce serait \all{der Beweis} ou \all{die Demonstration} seulement dans un contexte très formel/académique).
\end{itemize}
\end{attention}

\courssub{Résumé guidé : le texte en trois phrases}

Résumer (\all{zusammenfassen}) est une compétence clé du baccalauréat. Il s'agit de garder uniquement les \cle{idées principales} en trois ou quatre phrases maximum.

\begin{methode}[Réussir un résumé court]
\begin{enumerate}
\item Repère les idées principales : \emph{qui ?} (les citoyens), \emph{quoi ?} (ils ont des droits et des libertés), \emph{quelles libertés ?} (les quatre libertés), \emph{quelles limites ?} (ne pas insulter, menacer, appeler à la violence).
\item Rédige 3 à 4 phrases maximum, avec des connecteurs : \all{außerdem} (de plus), \all{aber} (mais), \all{deshalb} (c'est pourquoi).
\item N'utilise pas de citation, pas d'exemple, pas de détail : seulement l'essentiel.
\end{enumerate}
\end{methode}

\begin{exoresolu}[Résumé à trous]
\textbf{Énoncé.} Complète le résumé avec les mots suivants : \all{Grenzen — kämpfen — Freiheiten — versammeln — Pressefreiheit}.

\all{In einer Demokratie haben alle Bürger Rechte und \_\_\_\_\_\_\_\_\_\_\_\_ : Sie dürfen ihre Meinung frei sagen und sich friedlich \_\_\_\_\_\_\_\_\_\_\_\_. Außerdem ist die \_\_\_\_\_\_\_\_\_\_\_\_ sehr wichtig, weil eine freie Presse die Regierung kontrolliert. Aber die Freiheiten haben \_\_\_\_\_\_\_\_\_\_\_\_ : Man darf niemanden beleidigen. Deshalb müssen die Bürger jeden Tag für ihre Rechte \_\_\_\_\_\_\_\_\_\_\_\_.}

\tcblower
\textbf{Solution détaillée.}

\all{In einer Demokratie haben alle Bürger Rechte und \textbf{Freiheiten}: Sie dürfen ihre Meinung frei sagen und sich friedlich \textbf{versammeln}. Außerdem ist die \textbf{Pressefreiheit} sehr wichtig, weil eine freie Presse die Regierung kontrolliert. Aber die Freiheiten haben \textbf{Grenzen}: Man darf niemanden beleidigen. Deshalb müssen die Bürger jeden Tag für ihre Rechte \textbf{kämpfen}.}

« Dans une démocratie, tous les citoyens ont des droits et des libertés : ils peuvent exprimer librement leur opinion et se réunir pacifiquement. De plus, la liberté de la presse est très importante, car une presse libre contrôle le gouvernement. Mais les libertés ont des limites : on ne doit insulter personne. C'est pourquoi les citoyens doivent lutter chaque jour pour leurs droits. »

Vérifie la logique des connecteurs : \all{außerdem} ajoute une idée, \all{aber} introduit la limite, \all{deshalb} conclut.
\end{exoresolu}

\begin{exercice}[Bilan de la section Compréhension]
\begin{enumerate}
\item Réponds par une phrase complète : \all{Was bedeutet Religionsfreiheit?}
\item Vrai ou faux justifié : \all{„Man darf zu Gewalt aufrufen.“}
\item Trouve dans le texte deux mots composés avec \all{-freiheit} et donne leur traduction.
\item Résume le texte en deux phrases avec \all{aber} et \all{deshalb}.
\end{enumerate}
\end{exercice}

\begin{corrige}[Exercice : bilan de la section Compréhension]
\begin{enumerate}
\item \all{Religionsfreiheit bedeutet, dass jeder an Gott glauben oder nicht glauben darf.} « La liberté de religion signifie que chacun peut croire en Dieu ou ne pas croire. »
\item \all{Falsch. Man darf nicht zu Gewalt aufrufen. Im Text steht: „Man darf andere Menschen nicht beleidigen, nicht bedrohen und nicht zu Gewalt aufrufen.“}
\item Exemples : \all{die Meinungsfreiheit} (la liberté d'expression), \all{die Versammlungsfreiheit} (la liberté de réunion).
\item Modèle : \all{Alle Bürger haben Freiheiten, aber die Freiheiten haben Grenzen. Deshalb müssen die Bürger für ihre Rechte kämpfen.}
\end{enumerate}
\end{corrige}

\coursec{Grammaire 1 : Les verbes de modalité}

Les verbes de modalité (\cle{Modalverben}) expriment la manière dont une action est envisagée : capacité, permission, obligation, volonté, intention. Ils sont \textbf{indispensables} pour parler de droits (\all{dürfen}), de devoirs (\all{müssen}), de capacités (\all{können}) ou de volonté (\all{wollen}). Au programme de Terminale, on les maîtrise au \cle{Präsens}, au \cle{Präteritum} (forme narrative) et au \cle{Perfekt} (forme orale, avec double infinitif).

\begin{propriete}[Conjugaison au Présent (Präsens) — Les 6 modaux]
Les verbes \all{können, müssen, dürfen, wollen, sollen} ont une particularité : \textbf{pas de -e à la 1ère et 3ème personne du singulier} et \textbf{changement de voyelle} (sauf \all{sollen}) aux 3 personnes du singulier. \all{Möchten} (subjonctif II de \all{mögen}, « voudrais/aimerais ») se conjugue comme un verbe régulier.

\begin{center}
\begin{tabular}{|l|c|c|c|c|c|c|}\hline
\rowcolor{popblueL} \textbf{Personne} & \textbf{können} & \textbf{müssen} & \textbf{dürfen} & \textbf{wollen} & \textbf{sollen} & \textbf{möchten} \\ \hline
ich & kann & muss & darf & will & soll & möchte \\ \hline
du & kannst & musst & darfst & willst & sollst & möchtest \\ \hline
er/sie/es & kann & muss & darf & will & soll & möchte \\ \hline
wir & können & müssen & dürfen & wollen & sollen & möchten \\ \hline
ihr & könnt & müsst & dürft & wollt & sollt & möchtet \\ \hline
sie/Sie & können & müssen & dürfen & wollen & sollen & möchten \\ \hline
\end{tabular}
\end{center}

\textbf{Participe passé (Perfekt) :} \all{gekonnt, gemusst, gedurft, gewollt, gesollt, gemocht} (auxiliaire \all{haben}). \textbf{Attention :} au Perfekt, on utilise le \cle{double infinitif} : \all{Ich habe das machen können} (et non *\all{gekonnt machen}).
\end{propriete}

\begin{propriete}[Emploi et nuances — Contexte « Libertés et Devoirs »]
\begin{center}
\begin{tabularx}{\linewidth}{|l|X|X|}\hline
\rowcolor{popblueL} \textbf{Verbe} & \textbf{Sens principal} & \textbf{Exemple du thème (Libertés)} \\ \hline
\textbf{dürfen} & Permission, droit, autorisation & \all{Journalisten dürfen berichten.} (Les journalistes ont le droit d'informer.) \\ \hline
\textbf{müssen} & Obligation, contrainte (interne ou externe) & \all{Bürger müssen Rechte verteidigen.} (Les citoyens doivent défendre leurs droits.) \\ \hline
\textbf{können} & Capacité, possibilité, savoir-faire & \all{Wir können wählen gehen.} (Nous pouvons / sommes en mesure d'aller voter.) \\ \hline
\textbf{sollen} & Recommandation, attente, devoir moral & \all{Man soll die Meinung anderer respektieren.} (On devrait / on est censé respecter l'opinion d'autrui.) \\ \hline
\textbf{wollen} & Volonté, intention ferme & \all{Die Jugend will Veränderung.} (La jeunesse veut du changement.) \\ \hline
\textbf{möchten} & Souhait poli, désir (Konjunktiv II) & \all{Ich möchte meine Meinung sagen.} (J'aimerais exprimer mon opinion.) \\ \hline
\end{tabularx}
\end{center}
\end{propriete}

\begin{attention}[Pièges fréquents pour francophones]
\begin{itemize}
\item \textbf{Position de l'infinitif} : Le verbe modal se conjugue en \textbf{position 2}, l'infinitif du verbe principal va à la \textbf{fin de la phrase} (position finale). \all{Ich \textbf{muss} heute \textbf{lernen}.} (Pas *\all{Ich muss lernen heute}).
\item \textbf{\all{müssen} vs \all{sollen}} : \all{müssen} = contrainte forte (loi, nécessité) ; \all{sollen} = recommandation, ordre indirect, « on dit que ». \all{Du musst Steuern zahlen} (loi) vs \all{Du sollst pünktlich sein} (règle/conseil).
\item \textbf{\all{nicht müssen} vs \all{nicht dürfen}} : \all{nicht müssen} = « ne pas avoir à » (pas d'obligation) ; \all{nicht dürfen} = « interdiction » (pas le droit). \all{Du musst nicht kommen} = Tu n'es pas obligé de venir. \all{Du darfst nicht kommen} = Tu n'as pas le droit de venir.
\item \textbf{\all{möchten} n'a pas de Präteritum ni de Perfekt standard} : on utilise \all{wollte} / \all{hat gewollt} ou on reformule.
\end{itemize}
\end{attention}

\begin{exemplebox}[Exemples authentiques tirés du texte et prolongés]
\begin{itemize}
\item \all{„Jeder Mensch \textbf{darf} seine Meinung frei \textbf{sagen}.“} (Permission/Droit — Présent)
\item \all{„Die Bürger \textbf{müssen} ihre Rechte \textbf{kennen}.“} (Obligation — Présent)
\item \all{„Man \textbf{darf} andere \textbf{nicht beleidigen}.“} (Interdiction — \all{nicht dürfen})
\item \all{„Man \textbf{kann} die Regierung \textbf{kritisieren}.“} (Possibilité — prolongement)
\item \all{„Wir \textbf{wollen} in Freiheit \textbf{leben}.“} (Volonté — prolongement)
\item \all{„Man \textbf{sollte} (Konj. II) tolerant \textbf{sein}.“} (Recommandation — prolongement)
\end{itemize}
\end{exemplebox}

\begin{exercice}[Entraînement : Verbes de modalité]
\begin{enumerate}
\item Conjugue au présent : \all{ich \_\_\_\_\_\_\_\_ (können)}, \all{er \_\_\_\_\_\_\_\_ (müssen)}, \all{wir \_\_\_\_\_\_\_\_ (dürfen)}, \all{ihr \_\_\_\_\_\_\_\_ (wollen)}, \all{Sie \_\_\_\_\_\_\_\_ (sollen)}.
\item Complète avec le modal approprié (\all{dürfen, müssen, können, sollen, wollen}) à la bonne forme :
  \begin{enumerate}
  \item In einer Demokratie \_\_\_\_\_\_\_\_ alle Bürger wählen gehen.
  \item Man \_\_\_\_\_\_\_\_ andere nicht bedrohen.
  \item Journalisten \_\_\_\_\_\_\_\_ die Wahrheit berichten.
  \item Wir \_\_\_\_\_\_\_\_ unsere Rechte kennen.
  \item Ich \_\_\_\_\_\_\_\_ gerne Journalist werden.
  \end{enumerate}
\item Traduis en allemand (Thème) :
  \begin{enumerate}
  \item Les citoyens doivent respecter la loi.
  \item On ne doit pas insulter les autres.
  \item Nous pouvons manifester pacifiquement.
  \item Tu as le droit de donner ton avis.
  \item Il veut devenir avocat.
  \end{enumerate}
\end{enumerate}
\end{exercice}

\begin{corrige}[Exercice : Verbes de modalité]
\begin{enumerate}
\item \all{ich kann, er muss, wir dürfen, ihr wollt, Sie sollen}.
\item \begin{enumerate}
  \item \all{können} (capacité/droit) ou \all{dürfen} (permission légale) — \all{können} insiste sur la possibilité effective.
  \item \all{darf} (interdiction : \all{nicht dürfen}).
  \item \all{sollen} (devoir moral/expectative) ou \all{müssen} (obligation déontologique).
  \item \all{müssen} (obligation forte).
  \item \all{will} (volonté).
  \end{enumerate}
\item \begin{enumerate}
  \item \all{Die Bürger müssen das Gesetz respektieren.} (ou \all{achten})
  \item \all{Man darf andere nicht beleidigen.}
  \item \all{Wir können friedlich demonstrieren.} (ou \all{dürfen})
  \item \all{Du darfst deine Meinung sagen.}
  \item \all{Er will Anwalt werden.}
  \end{enumerate}
\end{enumerate}
\end{corrige}

\coursec{Grammaire 2 : La concession (obwohl / trotzdem)}

Exprimer une concession, c'est admettre un fait qui va à l'encontre de l'attente ou de l'argument principal (« bien que..., pourtant... »). En allemand, on distingue la \cle{subordonnée} (\all{obwohl}) et l'\cle{adverbe de liaison} (\all{trotzdem}).

\begin{propriete}[\all{obwohl} — Subordonnée de concession]
\begin{itemize}
\item \all{obwohl} = « bien que, quoique, même si ». Introduit une \textbf{subordonnée} : le verbe conjugué est \textbf{rejeté à la fin}.
\item L'ordre des mots : \all{obwohl} + Sujet + ... + \textbf{Verbe conjugué à la fin}.
\item La subordonnée peut être en début de phrase (position 1) ou à la fin.
\end{itemize}

\end{propriete}

\begin{exemplebox}[Structure \all{obwohl}]
\textbf{Subordonnée en premier (accent sur la concession) :}
\all{Obwohl die Freiheiten Grenzen haben, sind sie sehr wichtig.}
« Bien que les libertés aient des limites, elles sont très importantes. »
(Verbe principal \all{sind} en position 2 après la virgule.)

\textbf{Subordonnée en second (accent sur l'idée principale) :}
\all{Die Freiheiten sind sehr wichtig, obwohl sie Grenzen haben.}
« Les libertés sont très importantes, bien qu'elles aient des limites. »
\end{exemplebox}


\begin{propriete}[\all{trotzdem} — Adverbe de concession (conjonction adverbiale)]
\begin{itemize}
\item \all{trotzdem} = « pourtant, néanmoins, malgré tout ». C'est un \textbf{adverbe}, pas une conjonction de subordination.
\item Il occupe une \textbf{place dans la phrase principale} (souvent position 1 ou 3) et \textbf{ne renvoie pas le verbe à la fin}. Le verbe reste en \textbf{position 2}.
\item Souvent précédé d'un point ou d'un point-virgule ; lie deux phrases principales.
\end{itemize}

\end{propriete}

\begin{exemplebox}[Structure \all{trotzdem}]
\all{Die Freiheiten haben Grenzen. \textbf{Trotzdem} sind sie sehr wichtig.}
« Les libertés ont des limites. \textbf{Pourtant}, elles sont très importantes. »
(Verbe \all{sind} en position 2 après \all{trotzdem}.)

\all{Die Freiheiten haben Grenzen; sie sind \textbf{trotzdem} sehr wichtig.}
« Les libertés ont des limites ; elles sont \textbf{néanmoins} très importantes. »
(Verbe \all{sind} en position 2, \all{trotzdem} en position 3 après le sujet \all{sie}.)
\end{exemplebox}


\begin{aretenir}[Choisir \all{obwohl} ou \all{trotzdem} ?]
\begin{center}
\begin{tabularx}{\linewidth}{|X|X|}\hline
\rowcolor{popblueL} \textbf{Si tu veux une phrase complexe (1 phrase, 2 verbes)} & \textbf{Si tu veux deux phrases simples (2 phrases, 2 verbes)} \\ \hline
Utilise \all{obwohl} + verbe à la fin. & Utilise \all{trotzdem} + verbe en position 2. \\ \hline
\all{Obwohl es regnet, gehen wir spazieren.} & \all{Es regnet. Trotzdem gehen wir spazieren.} \\ \hline
\end{tabularx}
\end{center}
\textbf{Astuce de traduction :} « Bien que / quoique / même si » → \all{obwohl}. « Pourtant / néanmoins / malgré tout / cependant » (en début de phrase principale) → \all{trotzdem}.
\end{aretenir}

\begin{attention}[Erreurs classiques]
\begin{itemize}
\item \textbf{Verbe pas à la fin avec \all{obwohl}} : *\all{Obwohl die Freiheiten haben Grenzen...} → \textbf{Faux}. Correct : \all{Obwohl die Freiheiten Grenzen \textbf{haben}...}
\item \textbf{Verbe à la fin avec \all{trotzdem}} : *\all{Die Freiheiten haben Grenzen, trotzdem sie sind wichtig.} → \textbf{Faux}. \all{trotzdem} n'est pas une conjonction de subordination. Correct : \all{...trotzdem \textbf{sind} sie wichtig.}
\item \textbf{Confusion \all{obwohl} / \all{weil}} : \all{weil} = cause (parce que), \all{obwohl} = concession (bien que). \all{Ich bleibe zu Hause, weil es regnet} vs \all{Ich gehe raus, obwohl es regnet}.
\end{itemize}
\end{attention}

\begin{exercice}[Entraînement : Concession]
\begin{enumerate}
\item Relie les deux phrases par \all{obwohl} (deux possibilités : subordonnée au début ou à la fin).
  \begin{enumerate}
  \item Die Pressefreiheit ist wichtig. Es gibt Risiken für Journalisten.
  \item Man darf seine Meinung sagen. Man darf nicht beleidigen.
  \end{enumerate}
\item Relie les deux phrases par \all{trotzdem}.
  \begin{enumerate}
  \item In manchen Ländern gibt es Zensur. Die Menschen kämpfen für Pressefreiheit.
  \item Die Demonstration war verboten. Viele Bürger kamen.
  \end{enumerate}
\item Traduis en allemand :
  \begin{enumerate}
  \item Bien que la liberté ait des limites, elle est un droit humain.
  \item La manifestation est interdite. Pourtant, les gens viennent.
  \item Même si c'est difficile, nous devons défendre la démocratie.
  \end{enumerate}
\end{enumerate}
\end{exercice}

\begin{corrige}[Exercice : Concession]
\begin{enumerate}
\item \begin{enumerate}
  \item \all{Obwohl die Pressefreiheit wichtig ist, gibt es Risiken für Journalisten.} / \all{Es gibt Risiken für Journalisten, obwohl die Pressefreiheit wichtig ist.}
  \item \all{Obwohl man seine Meinung sagen darf, darf man nicht beleidigen.} / \all{Man darf nicht beleidigen, obwohl man seine Meinung sagen darf.}
  \end{enumerate}
\item \begin{enumerate}
  \item \all{In manchen Ländern gibt es Zensur. Trotzdem kämpfen die Menschen für Pressefreiheit.}
  \item \all{Die Demonstration war verboten. Trotzdem kamen viele Bürger.}
  \end{enumerate}
\item \begin{enumerate}
  \item \all{Obwohl die Freiheit Grenzen hat, ist sie ein Menschenrecht.}
  \item \all{Die Demonstration ist verboten. Trotzdem kommen die Leute.}
  \item \all{Obwohl es schwer ist, müssen wir die Demokratie verteidigen.}
  \end{enumerate}
\end{enumerate}
\end{corrige}

\coursec{Méthode du commentaire et commentaire modèle}

Au baccalauréat, l'épreuve de \cle{commentaire de texte} (ou \cle{analyse de texte}) exige une réponse structurée en allemand, montrant que tu as compris le texte, que tu sais l'analyser et que tu peux t'en servir pour exprimer ton opinion personnelle.

\begin{methode}[Structure type du commentaire écrit (3 parties)]
\begin{enumerate}
\item \textbf{Introduction} (\all{Einleitung}) — 3 à 4 phrases :
  \begin{itemize}
  \item \textbf{Accroche} : Présente le thème général (ex. \all{Das Thema Freiheiten ist aktuell...}).
  \item \textbf{Référence du texte} : Auteur (si connu), titre, source, date (ici : texte de cours \all{„Freiheiten in einer Demokratie“}).
  \item \textbf{Sujet / Thèse} : De quoi parle le texte ? Quel est le message principal ? (ex. \all{Der Text zeigt, dass Freiheiten in einer Demokratie unverzichtbar sind, aber Grenzen haben.}).
  \item \textbf{Annonce du plan} : \all{Im Folgenden werde ich... analysieren und meine Meinung darlegen.}
  \end{itemize}
\item \textbf{Corps / Développement} (\all{Hauptteil}) — 2 à 3 paragraphes arguments :
  \begin{itemize}
  \item \textbf{Paragraphe 1 : Analyse du contenu.} Reprends les 4 libertés, explique leur rôle (contr
\end{itemize}
\end{enumerate}
\end{methode}


\coursec{Lexique thématique : droits, libertés et devoirs}

La lecture du texte \all{„Freiheiten in einer Demokratie“} nous a permis de comprendre les grandes idées : dans une démocratie, les citoyens jouissent de libertés fondamentales, mais ces libertés ont des limites et s'accompagnent de devoirs. Pour parler de tout cela en allemand, il faut maintenant maîtriser le vocabulaire précis du thème. C'est l'objet de cette section : nous allons construire, observer et mémoriser le champ lexical des \cle{libertés}, des \cle{droits} et des \cle{devoirs}, avec pour chaque mot son article, son pluriel et sa traduction. Méthode de travail : pour chaque mot nouveau, apprenez toujours le \textbf{bloc complet} : article + nom + pluriel + traduction, par exemple \all{das Recht, die Rechte} = le droit.

\courssub{Observation : le vocabulaire dans son contexte}

Avant toute liste, lisons un court texte original qui emploie naturellement les mots que nous allons étudier. Observez les mots en gras : ce sont nos cibles.

\begin{textebox}[„Rechte und Pflichten — ein Schüler aus Yaoundé berichtet“]
In meiner Schule in Yaoundé sprechen wir oft über die \textbf{Freiheit}. Unsere Lehrerin erklärt: „Die \textbf{Meinungsfreiheit} ist ein wichtiges \textbf{Recht}. Jeder \textbf{Bürger} darf seine Meinung sagen, aber er muss auch die Meinung der anderen respektieren.“

In Kamerun gibt es viele Zeitungen und Radiosender. Die \textbf{Pressefreiheit} erlaubt den Journalisten, über Politik, Wirtschaft und Sport zu berichten. Aber die Pressefreiheit hat auch \textbf{Grenzen}: Ein Gesetz \textbf{verbietet} zum Beispiel Lügen und Beleidigungen.

Am Wochenende gehe ich mit meiner Familie in die Kirche; mein Freund Ahmad geht in die Moschee. Die \textbf{Religionsfreiheit} schützt uns beide. Der \textbf{Staat} erlaubt allen Bürgern, ihre Religion frei auszuüben.

Manchmal organisieren die Schüler eine Versammlung, um über Probleme in der Schule zu diskutieren. Das ist die \textbf{Versammlungsfreiheit}. Der Direktor \textbf{erlaubt} diese Treffen, aber wir müssen ruhig bleiben und die Regeln respektieren.

„Freiheit bedeutet nicht: Alles ist erlaubt“, sagt unsere Lehrerin. „Jedes Recht hat eine \textbf{Pflicht}. Ihr dürft eure Meinung äußern, aber ihr müsst die Gesetze respektieren. Und wenn ihr älter seid, dürft ihr wählen — das ist ein großes Recht in einer \textbf{Demokratie}.“

Ich finde: Für die Freiheit muss man kämpfen, aber man muss sie auch verantwortungsvoll nutzen.
\end{textebox}

\textbf{Glossaire :} der Bürger (-) = le citoyen ; die Grenze (-n) = la limite ; die Beleidigung (-en) = l'insulte ; die Versammlung (-en) = la réunion ; das Treffen (-) = la rencontre ; die Regel (-n) = la règle ; die Verantwortung = la responsabilité ; wählen = voter, choisir ; berichten = rapporter, informer ; ausüben (séparable : \all{übt ... aus}) = exercer, pratiquer ; verantwortungsvoll = de manière responsable ; nutzen = utiliser.

\textbf{Questions de compréhension rapide :}
\begin{enumerate}
\item \all{Welche vier Freiheiten nennt der Text?} (Quelles quatre libertés le texte mentionne-t-il ?)
\item \all{Was verbietet ein Gesetz?} (Qu'est-ce qu'une loi interdit ?)
\item \all{Was sagt die Lehrerin über Rechte und Pflichten?} (Que dit la professeure à propos des droits et des devoirs ?)
\end{enumerate}

\textbf{Réponses attendues :} 1. \all{Der Text nennt die Meinungsfreiheit, die Pressefreiheit, die Religionsfreiheit und die Versammlungsfreiheit.} 2. \all{Ein Gesetz verbietet Lügen und Beleidigungen.} 3. \all{Sie sagt: Jedes Recht hat eine Pflicht.}

\courssub{Le champ lexical des libertés}

Le mot central est \all{die Freiheit} (\cle{la liberté}), dérivé de l'adjectif \all{frei} (\cle{libre}). L'allemand adore les noms composés : en ajoutant un nom devant \all{-freiheit}, on obtient le nom de chaque liberté fondamentale.

\begin{propriete}[Les noms composés en -freiheit]
Tous les noms composés avec \all{-freiheit} sont \textbf{féminins} (le genre d'un nom composé est toujours celui du dernier élément) :
\begin{itemize}
\item \all{die Meinungsfreiheit} = la liberté d'expression (de \all{die Meinung}, l'opinion)
\item \all{die Pressefreiheit} = la liberté de la presse (de \all{die Presse})
\item \all{die Religionsfreiheit} = la liberté de religion (de \all{die Religion})
\item \all{die Versammlungsfreiheit} = la liberté de réunion (de \all{die Versammlung})
\end{itemize}
Ces noms composés s'emploient souvent \textbf{sans article} quand on parle du principe en général : \all{Meinungsfreiheit ist ein Grundrecht.} (La liberté d'expression est un droit fondamental.) Avec l'article, on insiste sur la liberté concrète d'un pays ou d'une situation : \all{Die Pressefreiheit ist in diesem Land wichtig.} (La liberté de la presse est importante dans ce pays.)
\end{propriete}

\begin{popfigure}
\begin{center}
\begin{tikzpicture}[
  grow=down,
  level 1/.style={sibling distance=5.5cm, level distance=1.6cm},
  level 2/.style={sibling distance=2.9cm, level distance=1.9cm},
  every node/.style={draw, rounded corners, align=center, font=\small, inner sep=4pt},
  edge from parent/.style={draw=popmuted, thick}
]
\node[fill=popgoldL, font=\small\bfseries, align=center]{frei\\ (libre, adj.)}
  child { node[fill=popblueL, font=\small\bfseries, align=center]{die Freiheit\\ (la liberté)}
    child { node[fill=popgreenL, align=center]{die Meinungs-\\freiheit\\ (expression)} }
    child { node[fill=popgreenL, align=center]{die Presse-\\freiheit\\ (presse)} }
    child { node[fill=popgreenL, align=center]{die Religions-\\freiheit\\ (religion)} }
    child { node[fill=popgreenL, align=center]{die Versammlungs-\\freiheit\\ (réunion)} }
  };
\end{tikzpicture}
\end{center}
\legende{Arbre de la famille de mots : de l'adjectif \all{frei} aux quatre libertés fondamentales.}
\end{popfigure}

\courssub{Le champ lexical des droits et des devoirs}

Une liberté est aussi un \cle{droit} (\all{das Recht}) garanti par la loi, et chaque droit implique un \cle{devoir} (\all{die Pflicht}). Voici les acteurs et les notions de ce petit théâtre démocratique :

\begin{definition}[Droits, devoirs et acteurs]
\begin{itemize}
\item \all{das Recht (-e)} = le droit ; \all{das Recht auf + Akkusativ} = le droit à quelque chose : \all{das Recht auf Meinungsfreiheit} (le droit à la liberté d'expression), \all{das Recht auf Bildung} (le droit à l'éducation).
\item \all{die Pflicht (-en)} = le devoir ; \all{die Pflicht haben, etwas zu tun} = avoir le devoir de faire quelque chose : \all{Jeder Bürger hat die Pflicht, die Gesetze zu respektieren.} (Chaque citoyen a le devoir de respecter les lois.)
\item \all{der Bürger (-)} / \all{die Bürgerin (-nen)} = le citoyen / la citoyenne.
\item \all{der Staat (-en)} = l'État ; \all{die Demokratie (-n)} = la démocratie ; \all{das Gesetz (-e)} = la loi.
\end{itemize}
\end{definition}

\begin{attention}[Faux amis et pièges de genre]
\begin{itemize}
\item \all{das Recht} = le droit (juridique), mais \all{recht haben} = avoir raison ! \all{Du hast recht.} = Tu as raison. Ne confondez pas non plus \all{das Recht} (le droit, notion abstraite ou ensemble de règles) avec \all{das Gesetz} (la loi concrète votée par le parlement).
\item \all{der Staat} se prononce « chtat » et prend un \textbf{-es} au génitif singulier : \all{die Bürger des Staates} (les citoyens de l'État).
\item En français, « la liberté » est féminin : bonne nouvelle, \all{die Freiheit} aussi. Mais attention : \all{das Recht} et \all{das Gesetz} sont \textbf{neutres}, alors que « le droit » et « la loi » ne vous donnent aucun indice de genre en français.
\item \all{die Pflicht} est féminin comme « le devoir » est masculin : ne vous fiez jamais au genre français !
\end{itemize}
\end{attention}

\courssub{Les verbes clés : erlauben, verbieten et les autres}

\begin{definition}[erlauben et verbieten]
\all{erlauben} = permettre, autoriser ; \all{verbieten} = interdire. Les deux se construisent avec un \textbf{datif de la personne} et un \textbf{accusatif de la chose} :\\
\all{Der Staat erlaubt den Bürgern (Dat.) Demonstrationen (Akk.).} = L'État autorise les citoyens à organiser des manifestations.\\
\all{Das Gesetz verbietet den Journalisten (Dat.) Lügen (Akk.).} = La loi interdit aux journalistes les mensonges.\\
On peut aussi ajouter une \textbf{infinitive avec \all{zu}} : \all{Der Staat erlaubt den Bürgern, ihre Religion frei auszuüben.} = L'État permet aux citoyens de pratiquer librement leur religion.
\end{definition}

\begin{exemplebox}[Exemples traduits]
\begin{itemize}
\item \all{Das Gesetz verbietet die Zensur.} = La loi interdit la censure.
\item \all{Die Schule erlaubt den Schülern, eine eigene Zeitung zu haben.} = L'école autorise les élèves à avoir leur propre journal.
\item \all{Die Bürger kämpfen für die Freiheit.} = Les citoyens luttent pour la liberté.
\item \all{Der Staat schützt die Menschenrechte.} = L'État protège les droits de l'homme.
\item \all{Wir müssen die Gesetze respektieren.} = Nous devons respecter les lois.
\item \all{Die Bürger wählen alle vier Jahre.} = Les citoyens votent tous les quatre ans.
\end{itemize}
\end{exemplebox}

\begin{attention}[verbieten est un verbe fort !]
\all{verbieten} se conjugue comme un verbe fort : \all{ich verbiete} (présent), \all{er verbot} (Präteritum), \all{er hat verboten} (Perfekt). Ne dites jamais \emph{« er hat verbietet »} ! Même famille que \all{bieten — bot — hat geboten}. En revanche, \all{erlauben} est un verbe faible régulier : \all{er erlaubte — er hat erlaubt}. Attention aussi à la place du verbe : dans \all{Das Gesetz verbietet die Zensur}, le verbe conjugué est bien en deuxième position.
\end{attention}

\courssub{Expressions utiles et familles de mots}

\begin{aretenir}[Expressions à réutiliser à l'écrit et à l'oral]
\begin{itemize}
\item \all{das Recht auf + Akk.} : \all{das Recht auf freie Meinungsäußerung} (le droit à la libre expression de l'opinion)
\item \all{seine Meinung sagen / äußern} = dire / exprimer son opinion
\item \all{für die Freiheit kämpfen} = lutter pour la liberté
\item \all{ein Gesetz respektieren / brechen} = respecter / violer une loi
\item \all{Rechte und Pflichten haben} = avoir des droits et des devoirs
\item \all{eine Versammlung organisieren} = organiser une réunion
\end{itemize}
\end{aretenir}

Les familles de mots permettent d'apprendre trois mots pour le prix d'un :

\begin{center}
\begin{tabularx}{\linewidth}{|X|X|X|}
\hline
\rowcolor{popblueL} Famille & Membres & Exemple traduit \\
\hline
\all{frei} & \all{frei} (libre) → \all{die Freiheit} (la liberté) → \all{die Meinungsfreiheit} & \all{Die Presse ist frei.} = La presse est libre. \\
\hline
\all{das Recht} & \all{das Recht} (le droit) → \all{gerecht} (juste) → \all{die Gerechtigkeit} (la justice) & \all{Das Urteil ist gerecht.} = Le jugement est juste. \\
\hline
\all{die Pflicht} & \all{die Pflicht} (le devoir) → \all{verpflichten} (obliger) → \all{verpflichtet} (obligé) & \all{Das Gesetz verpflichtet alle Bürger.} = La loi oblige tous les citoyens. \\
\hline
\end{tabularx}
\end{center}

\courssub{Tableau récapitulatif du vocabulaire}

\begin{center}
\begin{tabularx}{\linewidth}{|l|l|X|}
\hline
\rowcolor{popblueL} Deutsch (Artikel + Plural) & Français & Beispiel \\
\hline
\all{die Freiheit (-en)} & la liberté & \all{Die Freiheit ist ein Recht.} \\
\hline
\all{die Meinung (-en)} & l'opinion & \all{Ich sage meine Meinung.} \\
\hline
\all{das Recht (-e)} & le droit & \all{das Recht auf Bildung} \\
\hline
\all{die Pflicht (-en)} & le devoir & \all{Wählen ist auch eine Pflicht.} \\
\hline
\all{das Gesetz (-e)} & la loi & \all{Das Gesetz schützt uns.} \\
\hline
\all{die Grenze (-n)} & la limite & \all{Die Freiheit hat Grenzen.} \\
\hline
\all{der Staat (-en)} & l'État & \all{Der Staat erlaubt Proteste.} \\
\hline
\all{der Bürger (-) / die Bürgerin (-nen)} & le citoyen / la citoyenne & \all{Die Bürger wählen.} \\
\hline
\all{die Demokratie (-n)} & la démocratie & \all{Kamerun ist eine Demokratie.} \\
\hline
\all{erlauben / verbieten} & permettre / interdire & \all{Das Gesetz verbietet die Zensur.} \\
\hline
\end{tabularx}
\end{center}

\textbf{Complément utile au Bac :} \all{die Menschenrechte} (pluriel, les droits de l'homme), \all{die Gleichheit} (l'égalité), \all{die Zensur} (la censure). Ces trois mots reviennent très souvent dans les textes d'examen sur la citoyenneté : \all{Die Zensur verletzt die Menschenrechte.} = La censure viole les droits de l'homme.

\begin{savaistu}[La démocratie directe en Suisse]
En Suisse, les citoyens votent régulièrement par référendum, plusieurs fois par an, sur des questions précises (impôts, environnement, armée...) : c'est la démocratie directe, \all{die direkte Demokratie}, un exercice concret des droits politiques. Les Suisses disent : \all{Wir stimmen ab.} (Nous votons.) L'allemand possède d'ailleurs le mot \all{die Volksabstimmung} (le vote du peuple) pour désigner ce référendum. En Allemagne, les libertés fondamentales sont inscrites dans la Constitution, \all{das Grundgesetz} (littéralement « la Loi fondamentale ») : ses premiers articles protègent précisément la dignité humaine, la liberté d'expression, la liberté de la presse, la liberté de religion et la liberté de réunion.
\end{savaistu}

\begin{exoresolu}[Complétez avec le bon mot du lexique]
Énoncé — Complétez avec : \all{Meinungsfreiheit}, \all{verbietet}, \all{Pflicht}, \all{erlaubt}, \all{Gesetz}.\\
1. \all{Der Staat \_\_\_ den Bürgern Demonstrationen.}\\
2. \all{Die \_\_\_ ist ein Grundrecht: Jeder darf seine Meinung sagen.}\\
3. \all{Ein neues \_\_\_ schützt die Presse.}\\
4. \all{Das Gesetz \_\_\_ die Zensur.}\\
5. \all{Respekt ist auch eine \_\_\_.}
\tcblower
Solution détaillée :\\
1. \all{erlaubt} — le sujet \all{der Staat} exige la 3\textsuperscript{e} personne du singulier ; le sens est « autorise ». = L'État autorise les citoyens à organiser des manifestations.\\
2. \all{Meinungsfreiheit} — indice : \all{seine Meinung sagen}. = La liberté d'expression est un droit fondamental.\\
3. \all{Gesetz} — neutre (\all{ein neues}), c'est la loi qui protège. = Une nouvelle loi protège la presse.\\
4. \all{verbietet} — verbe fort, mais ici au présent : \all{das Gesetz verbietet}. Attention, au Perfekt le participe est irrégulier : \all{hat verboten}, jamais \emph{« hat verbietet »}.\\
5. \all{Pflicht} — le respect n'est pas seulement un droit, c'est aussi un devoir.
\end{exoresolu}

\begin{exercice}[À vous de jouer]
Traduisez en allemand : 1. La liberté de la presse est un droit important. 2. La loi interdit la censure. 3. L'école autorise les élèves à exprimer leur opinion. 4. Chaque citoyen a des droits et des devoirs. 5. Nous luttons pour la liberté et l'égalité.
\end{exercice}

\begin{corrige}[Exercice « À vous de jouer »]
1. \all{Die Pressefreiheit ist ein wichtiges Recht.} 2. \all{Das Gesetz verbietet die Zensur.} 3. \all{Die Schule erlaubt den Schülern, ihre Meinung zu äußern.} (datif \all{den Schülern} + infinitive avec \all{zu}) 4. \all{Jeder Bürger hat Rechte und Pflichten.} 5. \all{Wir kämpfen für die Freiheit und die Gleichheit.}
\end{corrige}

Ce lexique est désormais votre boîte à outils : gardez-le à portée de main, car les verbes de modalité (\all{dürfen}, \all{müssen}...) que nous étudierons dans la section grammaire vont nous permettre de le mettre en phrases complètes et nuancées.

\coursec{Grammaire 1 : les verbes de modalité}

Dans la section précédente, nous avons découvert le lexique des droits, des libertés et des devoirs : \all{die Freiheit}, \all{das Recht}, \all{die Pflicht}, \all{erlauben}, \all{verbieten}. Or, pour parler de ce qui est permis, obligatoire ou possible dans une démocratie, l'allemand dispose d'un outil grammatical central : les \cle{verbes de modalité} (\all{die Modalverben}). Ce sont eux qui permettent de dire « on a le droit de », « on doit », « on peut ». Reprenons deux phrases de notre texte support et observons-les de près.

\courssub{Observation : que remarque-t-on ?}

\begin{exemplebox}[Deux phrases du texte support]
\all{Jeder darf seine Meinung sagen.} « Chacun a le droit d'exprimer son opinion. »\par
\all{Wir müssen für unsere Rechte kämpfen.} « Nous devons lutter pour nos droits. »
\end{exemplebox}

Observons la construction de ces deux phrases :

\begin{itemize}
\item Le verbe conjugué (qui porte la marque de la personne) est \all{darf} (3\up{e} pers. sing. de \all{dürfen}) et \all{müssen} (1\up{re} pers. plur. de \all{müssen}) : il occupe la \textbf{deuxième position} de la phrase, comme tout verbe conjugué allemand.
\item Le verbe qui exprime l'action réelle (\all{sagen}, \all{kämpfen}) reste à \textbf{l'infinitif} et se place \textbf{en toute fin de phrase}.
\item Entre les deux, on trouve tous les compléments : \all{seine Meinung}, \all{für unsere Rechte}.
\end{itemize}

\begin{propriete}[La règle d'or du verbe de modalité]
Dans une phrase avec verbe de modalité, le verbe \textbf{conjugué} est le \cle{modal} (2\up{e} position) et le verbe lexical reste à \textbf{l'infinitif en fin de phrase}. Le modal et l'infinitif forment une sorte de « cadre » (\all{die Satzklammer}) qui enserre le reste de la phrase.\par
\all{Die Presse darf frei berichten.} « La presse a le droit d'informer librement. »
\end{propriete}

\begin{popfigure}
\begin{center}
\begin{tikzpicture}[
box/.style={draw=popblue, fill=popblueL, rounded corners=3pt, align=center, inner sep=5pt, font=\small},
ibox/.style={draw=poporange, fill=poporangeL, rounded corners=3pt, align=center, inner sep=5pt, font=\small},
mid/.style={draw=popmuted, fill=popcream, rounded corners=3pt, align=center, inner sep=5pt, font=\small}]
\node[box, align=center] (s) at (0,0){\textbf{Sujet}\\ \all{Der Bürger}};
\node[box, fill=poppinkL, draw=poppink, align=center] (m) at (3.4,0){\textbf{MODAL conjugué}\\ \all{darf} (2\up{e} position)};
\node[mid, align=center] (c) at (7.2,0){compléments\\ \all{frei} \ldots};
\node[ibox, align=center] (i) at (10.6,0){\textbf{INFINITIF}\\ \all{wählen} (fin)};
\draw[-{Stealth}, very thick, poppink] (m.north) to[bend left=35] node[above, font=\footnotesize, text=popdark]{le cadre verbal} (i.north);
\end{tikzpicture}
\end{center}
\legende{Structure de la phrase avec modal : le verbe conjugué ouvre le cadre, l'infinitif le ferme.}
\end{popfigure}

\courssub{Les cinq verbes de modalité et leur sens}

\begin{definition}[Les verbes de modalité]
Les \cle{Modalverben} expriment l'attitude du locuteur envers l'action : permission, obligation, capacité, conseil, volonté. Dans le thème des libertés, ils sont indispensables :\par
\begin{itemize}
\item \all{dürfen} = avoir le droit de, être autorisé à (\textbf{permission}) ;
\item \all{müssen} = devoir, être obligé de (\textbf{obligation}) ;
\item \all{können} = pouvoir, être capable de (\textbf{capacité, possibilité}) ;
\item \all{sollen} = devoir (au sens de recommandation, de devoir moral ou d'ordre reçu) ;
\item \all{wollen} = vouloir (\textbf{volonté, intention}).
\end{itemize}
\end{definition}

\begin{exemplebox}[Les modaux dans le thème des libertés]
\all{Man darf niemanden beleidigen.} « Il est interdit d'insulter qui que ce soit. »\par
\all{Wir müssen die Gesetze respektieren.} « Nous devons respecter les lois. »\par
\all{Die Bürger können ihre Religion frei wählen.} « Les citoyens peuvent choisir librement leur religion. »\par
\all{Jeder soll die Meinung der anderen tolerieren.} « Chacun est censé tolérer l'opinion des autres. »\par
\all{Die Jugendlichen wollen sich frei ausdrücken.} « Les jeunes veulent s'exprimer librement. »
\end{exemplebox}

\courssub{Conjugaison au présent}

Les modaux sont des verbes irréguliers : la voyelle radicale change aux trois personnes du singulier (\all{ich darf}, mais \all{wir dürfen}). Remarque aussi que la 1\up{re} et la 3\up{e} personne du singulier sont \textbf{identiques} (\all{ich darf, er darf}) — sans le \all{-t} habituel de la 3\up{e} personne.

\begin{center}
\begin{tabular}{|l|l|l|l|l|l|}
\hline
\rowcolor{popblueL} & \all{dürfen} & \all{müssen} & \all{können} & \all{sollen} & \all{wollen} \\
\hline
\all{ich} & darf & muss & kann & soll & will \\
\hline
\all{du} & darfst & musst & kannst & sollst & willst \\
\hline
\all{er/sie/es} & darf & muss & kann & soll & will \\
\hline
\all{wir} & dürfen & müssen & können & sollen & wollen \\
\hline
\all{ihr} & dürft & müsst & könnt & sollt & wollt \\
\hline
\all{sie/Sie} & dürfen & müssen & können & sollen & wollen \\
\hline
\end{tabular}
\end{center}

\begin{aretenir}[À retenir]
\begin{itemize}
\item Voyelle modifiée au singulier : \all{ü $\rightarrow$ a} (\all{dürfen/darf}, \all{müssen/muss}), \all{ö $\rightarrow$ a} (\all{können/kann}).
\item \all{sollen} ne change \textbf{pas} de voyelle : \all{ich soll, er soll}.
\item \all{wollen} : \all{o $\rightarrow$ i} au singulier : \all{ich will, du willst, er will}.
\item Pas de \all{-t} à la 3\up{e} personne du singulier : on dit \all{er muss} (jamais *\all{er musst}*).
\item Attention à \all{du musst} (avec \all{-ss-}, issu de \all{müssen}) : l'Umlaut disparaît au singulier.
\end{itemize}
\end{aretenir}

\courssub{La place du modal : phrase déclarative et subordonnée}

Dans la phrase déclarative, le modal conjugué est en 2\up{e} position et l'infinitif en fin :

\all{Der Bürger darf frei wählen.} « Le citoyen a le droit de voter librement. »

Mais dans une \textbf{subordonnée} (introduite par \all{dass}, \all{weil}, \all{obwohl}\ldots), le verbe conjugué passe en fin de phrase : c'est donc le \textbf{modal conjugué} qui occupe la dernière place, juste après l'infinitif.

\begin{propriete}[Le modal dans la subordonnée]
Subordonnée : sujet + compléments + \textbf{infinitif + modal conjugué} (en fin).\par
\all{\ldots dass der Bürger frei wählen darf.} « \ldots que le citoyen ait le droit de voter librement. »
\end{propriete}

\begin{exemplebox}[Modal en fin de subordonnée]
\all{Obwohl jeder seine Meinung sagen darf, gibt es Grenzen.} « Bien que chacun ait le droit d'exprimer son opinion, il y a des limites. » (le modal conjugué \all{darf} se place en fin de subordonnée.)\par
\all{Ich glaube, dass wir für die Pressefreiheit kämpfen müssen.} « Je pense que nous devons lutter pour la liberté de la presse. »\par
\all{Es ist gut, dass die Bürger ihre Religion frei wählen können.} « Il est bon que les citoyens puissent choisir librement leur religion. »
\end{exemplebox}

\courssub{Modal + négation : un piège capital}

La négation change radicalement de sens selon le modal employé. C'est l'un des pièges les plus fréquents au baccalauréat.

\begin{center}
\begin{tabularx}{\linewidth}{|X|X|X|}
\hline
\rowcolor{popblueL} Français & Deutsch & Remarque \\
\hline
il est interdit de / on ne doit pas & \all{man darf nicht} & interdiction \\
\hline
on n'est pas obligé de & \all{man muss nicht} & absence d'obligation \\
\hline
on ne peut pas & \all{man kann nicht} & impossibilité \\
\hline
\end{tabularx}
\end{center}

\begin{attention}[« Tu ne dois pas fumer » : attention au faux sens !]
\all{Man darf nicht} signifie « c'est interdit » ; \all{man muss nicht} signifie « ce n'est pas obligatoire ». Ne traduis donc JAMAIS « tu ne dois pas fumer » par *\all{du musst nicht rauchen}* (qui voudrait dire « tu n'es pas obligé de fumer » !), mais par \all{du darfst nicht rauchen}.\par
\all{Du musst nicht kommen.} « Tu n'es pas obligé de venir. »\par
\all{Du darfst hier nicht rauchen.} « Il t'est interdit de fumer ici. »
\end{attention}

\courssub{Le prétérit des modaux (utile au Bac)}

Pour raconter l'histoire des droits et des libertés (par exemple l'évolution du droit de vote), on emploie le prétérit des modaux, très fréquent à l'écrit. Bonne nouvelle : au prétérit, les modaux perdent leur Umlaut et se conjuguent comme des verbes faibles.

\begin{center}
\begin{tabular}{|l|l|l|}
\hline
\rowcolor{popblueL} Infinitif & Präteritum & Sens au prétérit \\
\hline
\all{dürfen} & \all{durfte} & avait le droit de \\
\hline
\all{müssen} & \all{musste} & devait, a dû \\
\hline
\all{können} & \all{konnte} & pouvait \\
\hline
\all{sollen} & \all{sollte} & devait (recommandation) \\
\hline
\all{wollen} & \all{wollte} & voulait \\
\hline
\end{tabular}
\end{center}

\begin{exemplebox}[Les modaux au prétérit]
\all{Früher durften die Frauen nicht wählen.} « Autrefois, les femmes n'avaient pas le droit de voter. »\par
\all{Die Journalisten mussten vorsichtig schreiben.} « Les journalistes devaient écrire avec prudence. »\par
\all{Niemand konnte frei reisen.} « Personne ne pouvait voyager librement. »\par
\all{Die Bürger wollten mehr Freiheit.} « Les citoyens voulaient plus de liberté. »
\end{exemplebox}

\courssub{Comparaison français / allemand}

\begin{center}
\begin{tabularx}{\linewidth}{|X|X|X|}
\hline
\rowcolor{popblueL} Français & Deutsch & Remarque \\
\hline
Chacun a le droit de parler. & \all{Jeder darf sprechen.} & même logique : verbe conjugué + infinitif \\
\hline
Nous devons respecter les lois. & \all{Wir müssen die Gesetze respektieren.} & en allemand, l'infinitif est rejeté en fin \\
\hline
Je peux exprimer mon opinion. & \all{Ich kann meine Meinung ausdrücken.} & le complément passe avant l'infinitif \\
\hline
Chacun doit pouvoir voter. & \all{Jeder muss wählen können.} & deux verbes à l'infinitif en fin : le modal lexical (\all{können}) ferme la phrase \\
\hline
\end{tabularx}
\end{center}

Le français et l'allemand conjuguent tous deux « pouvoir/devoir » suivi d'un infinitif ; la grande différence, c'est que l'allemand rejette l'infinitif \textbf{en toute fin de phrase}, après tous les compléments.

\begin{attention}[Erreur fréquente des francophones]
Ne place jamais l'infinitif juste après le modal, comme en français : *\all{Ich darf sagen meine Meinung}* est FAUX. L'infinitif va en fin de phrase : \all{Ich darf meine Meinung sagen.} De même, n'oublie pas la majuscule aux noms : \all{die Meinung}, \all{die Gesetze}, \all{die Freiheit}.
\end{attention}

\courssub{Texte d'application : les modaux en contexte}

\begin{textebox}[Freiheit hat auch Regeln]
In einer Demokratie darf jeder Bürger seine Meinung frei sagen. Man kann die Zeitungen lesen, die man will, und man darf sich friedlich versammeln. Auch die Religion kann jeder frei wählen: Niemand muss einer Kirche angehören, aber jeder darf es. Doch die Freiheit hat Grenzen. Man darf andere Menschen nicht beleidigen, und man muss die Gesetze respektieren. Journalisten können frei berichten, aber sie sollen die Wahrheit sagen und dürfen keine Lügen verbreiten. Früher war das anders: Viele Menschen durften nicht wählen, und die Presse konnte nicht frei schreiben. Wer Kritik äußern wollte, musste vorsichtig sein. Heute wollen die jungen Leute ihre Rechte kennen und nutzen. „Wir müssen unsere Freiheiten schützen“, sagt eine Schülerin aus Yaoundé, „aber wir sollen auch die Rechte der anderen achten.“ Ein Experte erklärt: „Freiheit bedeutet nicht, dass man alles tun darf. Sie bedeutet, dass jeder seine Rechte kennt und die Pflichten nicht vergisst.“ Deshalb müssen die Bürger informiert bleiben und dürfen ihre Verantwortung nicht vergessen. Nur so kann eine Gesellschaft frei und gerecht bleiben.
\end{textebox}

\textbf{Glossaire :} \all{sich versammeln} (se rassembler) ; \all{die Kirche, -n} (l'église) ; \all{angehören + Dat.} (appartenir à) ; \all{beleidigen} (insulter) ; \all{die Lüge, -n} (le mensonge) ; \all{verbreiten} (répandre) ; \all{Kritik äußern} (exprimer une critique) ; \all{achten} (respecter) ; \all{die Pflicht, -en} (le devoir) ; \all{die Verantwortung} (la responsabilité) ; \all{gerecht} (juste).

\textbf{Questions de compréhension :}
\begin{enumerate}
\item \all{Was darf jeder Bürger in einer Demokratie?}
\item \all{Was sollen Journalisten tun?}
\item \all{Wie war die Situation früher?}
\item \all{Was bedeutet Freiheit für den Experten?}
\end{enumerate}

\courssub{Exercice résolu}

\begin{exoresolu}[Modalverben einsetzen und übersetzen]
Complète avec le modal qui convient au présent, puis traduis la phrase en français :\par
1. Jeder Bürger \ldots\ frei wählen. (avoir le droit de)\par
2. Wir \ldots\ die Gesetze respektieren. (devoir, obligation)\par
3. Ihr \ldots\ eure Religion frei wählen. (pouvoir)\par
4. Du \ldots\ hier nicht rauchen. (c'est interdit)\par
5. Früher \ldots\ die Frauen nicht wählen. (prétérit, ne pas avoir le droit)
\tcblower
\textbf{Solution détaillée :}\par
1. \all{Jeder Bürger darf frei wählen.} « Chaque citoyen a le droit de voter librement. » — permission : \all{dürfen}, 3\up{e} pers. sing. = \all{darf}.\par
2. \all{Wir müssen die Gesetze respektieren.} « Nous devons respecter les lois. » — obligation : \all{müssen}.\par
3. \all{Ihr könnt eure Religion frei wählen.} « Vous pouvez choisir librement votre religion. » — possibilité : \all{können}, 2\up{e} pers. plur. = \all{könnt}.\par
4. \all{Du darfst hier nicht rauchen.} « Il t'est interdit de fumer ici. » — interdiction = \all{dürfen + nicht} ; *\all{musst nicht}* signifierait « tu n'es pas obligé ».\par
5. \all{Früher durften die Frauen nicht wählen.} « Autrefois, les femmes n'avaient pas le droit de voter. » — prétérit de \all{dürfen} : \all{durften}. Dans tous les cas, l'infinitif reste en fin de phrase.
\end{exoresolu}

\begin{exercice}[À toi de jouer]
Traduis en allemand : 1. Chacun a le droit d'exprimer son opinion. 2. Nous devons lutter pour nos droits. 3. Les citoyens peuvent se rassembler pacifiquement (\all{sich friedlich versammeln}). 4. Il est interdit de répandre des mensonges. 5. Bien que chacun puisse voter librement, peu de jeunes votent.
\end{exercice}

\begin{corrige}[Exercice « À toi de jouer »]
1. \all{Jeder darf seine Meinung sagen (ausdrücken).}\par
2. \all{Wir müssen für unsere Rechte kämpfen.}\par
3. \all{Die Bürger können sich friedlich versammeln.}\par
4. \all{Man darf keine Lügen verbreiten.}\par
5. \all{Obwohl jeder frei wählen kann, wählen wenige junge Leute.} (modal conjugué \all{kann} en fin de subordonnée.)
\end{corrige}

\coursec{Grammaire 2 : la concession (obwohl, trotzdem)}

Dans la section précédente, nous avons étudié les \cle{verbes de modalité} (\all{dürfen, müssen, können}...), qui permettent d'exprimer ce qui est permis, obligatoire ou possible dans une démocratie : \all{Man darf demonstrieren.} « On a le droit de manifester. » Mais une liberté n'est jamais absolue : elle rencontre des limites, des obstacles, des contradictions. Pour exprimer cette opposition entre deux idées — « les libertés sont importantes, \emph{et pourtant} elles ont des limites » —, l'allemand dispose de deux outils que le Bac adore et que les francophones confondent souvent : la conjonction \cle{obwohl} et l'adverbe \cle{trotzdem}. C'est l'objet de cette deuxième leçon de grammaire.

\courssub{Observation : deux constructions pour une même idée}

Reprenons deux phrases du texte support \emph{Freiheiten in einer Demokratie} :

\begin{exemplebox}[Deux phrases du texte]
\all{Obwohl diese Freiheiten wichtig sind, haben sie auch Grenzen.}\\
« Bien que ces libertés soient importantes, elles ont aussi des limites. »

\all{Die Freiheiten sind wichtig. Trotzdem haben sie Grenzen.}\\
« Les libertés sont importantes. Pourtant, elles ont des limites. »
\end{exemplebox}

Observons attentivement. Les deux phrases expriment \textbf{exactement la même opposition} entre deux idées : (1) les libertés sont importantes ; (2) elles ont des limites. Pourtant, la \textbf{construction grammaticale} est totalement différente :

\begin{itemize}
\item Avec \all{obwohl}, les deux idées sont réunies en \textbf{une seule phrase} : une subordonnée (\all{Obwohl diese Freiheiten wichtig sind}) suivie d'une principale (\all{haben sie auch Grenzen}). On remarque que le verbe conjugué de la subordonnée, \all{sind}, est rejeté \textbf{à la toute fin}, et que le verbe de la principale, \all{haben}, vient \textbf{immédiatement après la virgule}.
\item Avec \all{trotzdem}, on a \textbf{deux phrases indépendantes}. \all{Trotzdem} occupe la première position de la seconde phrase et provoque l'\textbf{inversion} : le verbe \all{haben} passe devant le sujet \all{sie}.
\end{itemize}

Retenons dès maintenant l'idée directrice : \all{obwohl} est une \textbf{conjonction de subordination} (comme \all{weil}, \all{dass}, \all{wenn}), tandis que \all{trotzdem} est un \textbf{adverbe de liaison} (comme \all{deshalb}, \all{danach}). Toute la différence de construction découle de cette différence de nature.

\begin{popfigure}
\begin{tikzpicture}[
  box/.style={draw=popblue, thick, rounded corners, fill=popblueL, align=center, font=\small, inner sep=6pt},
  obox/.style={draw=poporange, thick, rounded corners, fill=poporangeL, align=center, font=\small, inner sep=6pt},
  lab/.style={font=\bfseries\small},
  arr/.style={-{Stealth[length=3mm]}, thick, popdark}
]
\node[lab, popblue] (t1) at (0,0) {OBWOHL (conjonction)};
\node[box, align=center] (s1) at (-2.6,-1.6){Subordonnée\\ verbe conjugué \textbf{en fin}\\ \all{Obwohl es regnet,}};
\node[box, align=center] (p1) at (2.6,-1.6){Principale\\ verbe \textbf{en position 1}\\ \all{gehen wir spazieren.}};
\draw[arr] (s1) -- (p1);
\node[lab, poporange] (t2) at (0,-3.6) {TROTZDEM (adverbe)};
\node[obox, align=center] (p2) at (-2.6,-5.2){Première phrase principale\\ \all{Es regnet.}};
\node[obox, align=center] (p3) at (2.6,-5.2){Seconde phrase : verbe \textbf{en position 2}\\ inversion : \all{Trotzdem gehen wir.}};
\draw[arr] (p2) -- (p3);
\end{tikzpicture}
\legende{Deux constructions pour une même concession : subordination avec \all{obwohl}, juxtaposition avec inversion après \all{trotzdem}.}
\end{popfigure}

\courssub{obwohl : la concession en subordonnée}

\begin{propriete}[obwohl + subordonnée]
\cle{obwohl} « bien que, quoique » est une \textbf{conjonction de subordination}. Elle introduit une subordonnée dans laquelle le \textbf{verbe conjugué se place en dernière position}.

Deux schémas possibles :
\begin{enumerate}
\item Principale + subordonnée : \all{Die Presse muss die Wahrheit sagen, obwohl sie frei ist.}
\item Subordonnée en tête + principale inversée : \all{Obwohl die Presse frei ist, muss sie die Wahrheit sagen.} — quand la subordonnée occupe la première position, elle compte comme la \textbf{position 1} de la phrase ; le verbe de la principale vient donc \textbf{immédiatement après la virgule} (inversion sujet-verbe).
\end{enumerate}

\all{Obwohl die Presse frei ist, muss sie die Wahrheit sagen.} = « Bien que la presse soit libre, elle doit dire la vérité. »
\end{propriete}

\begin{attention}[obwohl + indicatif, pas de subjonctif !]
En français, « bien que » exige le \textbf{subjonctif} : « bien que la presse \emph{soit} libre ». En allemand, \all{obwohl} est suivi de l'\textbf{indicatif} : \all{obwohl sie frei \emph{ist}}. Ne cherchez donc jamais à calquer le subjonctif français ; conjuguez simplement le verbe à l'indicatif, au temps logique voulu : \all{Obwohl es gestern geregnet hat, sind wir gekommen.} « Bien qu'il ait plu hier, nous sommes venus. » C'est un point où l'allemand est plus facile que le français !
\end{attention}

\begin{exemplebox}[obwohl en contexte « libertés »]
\all{Obwohl Demonstrationen erlaubt sind, gibt es Regeln.} « Bien que les manifestations soient autorisées, il y a des règles. »

\all{Obwohl jeder seine Meinung sagen darf, darf niemand beleidigen.} « Bien que chacun puisse dire son opinion, personne n'a le droit d'insulter. »

\all{Die Schüler diskutieren laut, obwohl der Lehrer schon da ist.} « Les élèves discutent bruyamment, bien que le professeur soit déjà là. »

\all{Obwohl das Thema schwierig ist, bleiben alle konzentriert.} « Bien que le sujet soit difficile, tous restent concentrés. »
\end{exemplebox}

\courssub{trotzdem : la concession en phrase principale}

\begin{propriete}[trotzdem + inversion]
\cle{trotzdem} « pourtant, néanmoins, quand même » est un \textbf{adverbe de liaison}. Il se place dans une phrase principale (souvent une nouvelle phrase, après un point). Placé en \textbf{première position}, il provoque l'\textbf{inversion} : le verbe conjugué reste en deuxième position, donc \textbf{devant le sujet}.

\all{Die Presse ist frei. Trotzdem muss sie die Wahrheit sagen.} = « La presse est libre. Pourtant, elle doit dire la vérité. »

On peut aussi placer \all{trotzdem} plus loin dans la phrase (position 3), sans inversion : \all{Die Presse muss trotzdem die Wahrheit sagen.}
\end{propriete}

\begin{exemplebox}[trotzdem en contexte « libertés »]
\all{Das Thema ist schwierig. Trotzdem diskutieren die Schüler darüber.} « Le sujet est difficile. Pourtant, les élèves en discutent. »

\all{Es regnet in Yaoundé. Trotzdem gehen die Demonstranten auf die Straße.} « Il pleut à Yaoundé. Pourtant, les manifestants descendent dans la rue. »

\all{Der Artikel ist kritisch. Trotzdem wird er veröffentlicht.} « L'article est critique. Il est pourtant publié. »

\all{Die Prüfung war hart. Trotzdem haben viele Schüler bestanden.} « L'examen était dur. Pourtant, beaucoup d'élèves ont réussi. »
\end{exemplebox}

\courssub{Correspondances : transformer obwohl en trotzdem}

Une même réalité peut s'exprimer des deux façons. Au Bac, on vous demandera souvent de \textbf{transformer} une phrase. Voici l'équivalence fondamentale :

\begin{center}
\begin{tabularx}{\linewidth}{|X|X|}
\hline
\rowcolor{popblueL} \textbf{Avec obwohl (subordonnée)} & \textbf{Avec trotzdem (deux phrases)} \\
\hline
\all{Obwohl es verboten ist, demonstrieren sie.} & \all{Es ist verboten. Trotzdem demonstrieren sie.} \\
\hline
\all{Obwohl sie müde sind, lernen die Schüler.} & \all{Sie sind müde. Trotzdem lernen die Schüler.} \\
\hline
\all{Obwohl die Zeitung kritisiert wird, erscheint sie weiter.} & \all{Die Zeitung wird kritisiert. Trotzdem erscheint sie weiter.} \\
\hline
\end{tabularx}
\end{center}

Remarquez bien le sens : avec \all{obwohl}, l'obstacle est dans la subordonnée ; avec \all{trotzdem}, l'obstacle est dans la \textbf{première} phrase et le résultat surprenant dans la seconde. L'ordre logique s'inverse : \all{Obwohl A, B} = \all{A. Trotzdem B.}

\begin{methode}[Transformer une phrase de concession en 3 étapes]
\begin{enumerate}
\item \textbf{Repérer les deux idées opposées} : quelle est la concession (l'obstacle) ? quelle est la conséquence surprenante ? Ex. : obstacle = « il pleut » ; résultat = « nous sortons ».
\item \textbf{Choisir l'outil} : \all{obwohl} si l'on veut une seule phrase avec subordonnée ; \all{trotzdem} si l'on veut deux phrases principales.
\item \textbf{Vérifier la place du verbe} : avec \all{obwohl}, verbe conjugué à la fin de la subordonnée, et verbe de la principale juste après la virgule ; avec \all{trotzdem}, inversion sujet-verbe dans la seconde phrase. Relire à voix haute pour contrôler.
\end{enumerate}
\end{methode}

\begin{exoresolu}[Transformation : de obwohl à trotzdem]
Énoncé : transformez la phrase suivante en utilisant \all{trotzdem} : \all{Obwohl die Meinungsfreiheit wichtig ist, hat sie Grenzen.}
\tcblower
\textbf{Étape 1} : les deux idées opposées sont « la liberté d'expression est importante » (concession) et « elle a des limites » (résultat).\\
\textbf{Étape 2} : avec \all{trotzdem}, on fait deux phrases : la concession devient la première phrase, le résultat la seconde.\\
\textbf{Étape 3} : \all{trotzdem} en tête provoque l'inversion : le verbe \all{hat} passe devant le sujet \all{sie}.\\
\textbf{Solution} : \all{Die Meinungsfreiheit ist wichtig. Trotzdem hat sie Grenzen.}\\
« La liberté d'expression est importante. Pourtant, elle a des limites. »\\
Vérification inverse : \all{Obwohl die Meinungsfreiheit wichtig ist, hat sie Grenzen.} — verbe \all{ist} en fin de subordonnée, verbe \all{hat} juste après la virgule : correct.
\end{exoresolu}

\courssub{Complément Bac : trotz + génitif et zwar... aber}

\begin{propriete}[Autres connecteurs de concession]
\begin{itemize}
\item \cle{trotz} + \textbf{génitif} « malgré » : préposition suivie d'un nom au génitif. \all{Trotz des Verbots demonstrieren die Studenten.} « Malgré l'interdiction, les étudiants manifestent. » (\all{das Verbot} $\rightarrow$ \all{des Verbots}). Attention : \all{trotz} est suivi d'un \textbf{nom}, pas d'une phrase ; pour une phrase, il faut \all{obwohl} ou \all{trotzdem}.
\item \cle{zwar... aber} « certes... mais » : \all{zwar} admet le premier fait, \all{aber} introduit la restriction. \all{Die Presse ist zwar frei, aber sie muss die Wahrheit sagen.} « La presse est certes libre, mais elle doit dire la vérité. » \all{aber} est une conjonction de coordination : pas d'inversion après elle (\all{aber sie muss}).
\end{itemize}
\end{propriete}

\begin{center}
\begin{tabularx}{\linewidth}{|l|X|X|}
\hline
\rowcolor{popblueL} \textbf{Français} & \textbf{Deutsch} & \textbf{Construction} \\
\hline
bien que + subjonctif & \all{obwohl} + indicatif & subordonnée, verbe en fin \\
\hline
pourtant, néanmoins & \all{trotzdem} & phrase principale, inversion \\
\hline
malgré + nom & \all{trotz} + génitif & préposition + nom \\
\hline
certes... mais & \all{zwar... aber} & coordination, sans inversion \\
\hline
\end{tabularx}
\end{center}

\courssub{Entraînement guidé}

\begin{textebox}[Freiheit braucht Mut]
\all{Obwohl viele Menschen ihre Rechte kennen, schweigen sie oft. In einer Demokratie darf jeder seine Meinung sagen. Trotzdem haben viele Angst, laut zu sprechen. Eine Journalistin aus Douala erklärt: „Obwohl es manchmal gefährlich ist, schreibe ich die Wahrheit.“ Sie kämpft trotz des Drucks weiter. Die Presse ist zwar frei, aber sie trägt auch eine große Verantwortung. Obwohl Demonstrationen erlaubt sind, müssen die Bürger Regeln respektieren. Freiheit ist schwierig. Trotzdem lohnt es sich, jeden Tag dafür zu kämpfen.}

\textbf{Glossaire} : \all{das Recht, -e} le droit ; \all{schweigen} se taire ; \all{die Angst, die Ängste} la peur ; \all{gefährlich} dangereux ; \all{der Druck} la pression ; \all{die Verantwortung} la responsabilité ; \all{respektieren} respecter ; \all{es lohnt sich} cela vaut la peine.
\end{textebox}

\begin{exercice}[À vous]
\begin{enumerate}
\item Soulignez dans le texte ci-dessus tous les connecteurs de concession et indiquez leur nature (conjonction, adverbe, préposition).
\item Traduisez : « Bien que la liberté de la presse existe, elle est menacée. »
\item Transformez avec \all{trotzdem} : \all{Obwohl es verboten ist, demonstrieren die Studenten.}
\item Complétez : \all{Obwohl die Schüler wenig Zeit haben, ... (lernen / sie / viel)}.
\end{enumerate}
\end{exercice}

\begin{corrige}[Exercice « À vous »]
1. \all{obwohl} (conjonction de subordination), \all{trotzdem} (adverbe), \all{trotz des Drucks} (préposition + génitif), \all{zwar... aber} (adverbe + conjonction de coordination).\\
2. \all{Obwohl die Pressefreiheit existiert, ist sie bedroht.} (indicatif après \all{obwohl} ; verbe de la principale juste après la virgule).\\
3. \all{Es ist verboten. Trotzdem demonstrieren die Studenten.}\\
4. \all{Obwohl die Schüler wenig Zeit haben, lernen sie viel.} — le verbe \all{lernen} vient immédiatement après la virgule, puis le sujet \all{sie}.
\end{corrige}

\begin{attention}[Les trois pièges du francophone]
\begin{enumerate}
\item Ne jamais utiliser \all{obwohl} comme un adverbe : *\all{Obwohl, wir kämpfen.} est impossible. Si vous voulez dire « pourtant », utilisez \all{trotzdem}.
\item Après \all{obwohl} en tête de phrase, le verbe de la principale vient \textbf{tout de suite après la virgule} : \all{Obwohl es regnet, gehen wir} (et non *\all{Obwohl es regnet, wir gehen}).
\item Après \all{trotzdem}, n'oubliez pas l'inversion : \all{Trotzdem müssen wir kämpfen} (et non *\all{Trotzdem wir müssen kämpfen}). Le verbe reste toujours en deuxième position.
\end{enumerate}
\end{attention}

\begin{aretenir}[L'essentiel de la concession]
\all{obwohl} « bien que » = conjonction de subordination + indicatif : verbe conjugué en fin de subordonnée ; si la subordonnée est en tête, inversion dans la principale. \all{trotzdem} « pourtant » = adverbe : première position + inversion sujet-verbe. Équivalence : \all{Obwohl A, B} = \all{A. Trotzdem B.} Compléments : \all{trotz} + génitif « malgré », \all{zwar... aber} « certes... mais ».
\end{aretenir}

\coursec{Méthode du commentaire et commentaire modèle}

Après avoir étudié la concession avec \all{obwohl} et \all{trotzdem}, vous disposez maintenant de tous les outils grammaticaux de la leçon : les verbes de modalité (\all{dürfen, müssen, können, sollen}) et les connecteurs de concession. Il est temps de les réemployer dans une tâche d'expression écrite typique du baccalauréat : le \cle{commentaire de texte} (\all{die Textanalyse} / \all{der Kommentar}). Cette section vous donne une méthode infaillible en cinq étapes, un plan type en trois parties, les expressions utiles, puis un commentaire modèle rédigé et analysé.

\courssub{Les cinq étapes du commentaire de texte au Bac}

Observer d'abord comment un bon candidat procède. Le commentaire n'est pas un résumé mot à mot : c'est une \cle{analyse organisée} qui montre que vous avez compris le texte et que vous savez en parler avec vos propres mots.

\begin{methode}[Commenter un texte au Bac : la méthode en 5 étapes]
\begin{enumerate}
\item \textbf{Présenter le texte} (\all{den Text vorstellen}) : indiquez le type de texte (article, discours, lettre...), le thème et la source. Exemple : \all{Der Text ist ein Artikel über die Freiheiten in einer Demokratie.}
\item \textbf{Dégager l'idée principale} (\all{die Hauptidee}) : en une ou deux phrases, dites de quoi il s'agit vraiment. Exemple : \all{Der Autor zeigt, dass die Freiheiten wichtig, aber nicht grenzenlos sind.}
\item \textbf{Analyser la structure} (\all{die Struktur}) : pour notre texte, repérez les trois mouvements : a) la présentation des libertés, b) leurs limites, c) l'appel à l'action final.
\item \textbf{Relever les procédés} (\all{die Stilmittel}) : énumérations, connecteurs (\all{aber, deshalb, obwohl}), verbes de modalité (\all{dürfen, müssen}). Citez brièvement le texte entre guillemets pour justifier : \all{Der Autor schreibt: „Man darf niemanden beleidigen.“}
\item \textbf{Donner son avis argumenté} (\all{die eigene Meinung}) : terminez toujours par votre opinion personnelle avec \all{meiner Meinung nach} ou \all{ich bin der Meinung, dass...}, si possible en comparant avec la situation au Cameroun.
\end{enumerate}
\end{methode}

\begin{aretenir}[Les trois règles d'or du commentaire]
\begin{itemize}
\item Suivez \textbf{toujours} le plan en trois parties : \all{Einleitung} → \all{Hauptteil} → \all{Schluss}.
\item \textbf{Citez brièvement} le texte entre guillemets „…“ pour justifier votre analyse, sans jamais recopier de longs passages.
\item \textbf{Terminez par votre avis personnel} avec \all{meiner Meinung nach} : un commentaire sans opinion personnelle perd des points.
\end{itemize}
\end{aretenir}

\courssub{Le plan type : Einleitung, Hauptteil, Schluss}

Le commentaire allemand suit une architecture rigide, plus stricte qu'en français. Retenez les trois parties et leur contenu :

\begin{center}
\begin{tabularx}{\linewidth}{|l|X|X|}
\hline
\rowcolor{popblueL} \textbf{Partie} & \textbf{Contenu} & \textbf{Français} \\
\hline
\all{die Einleitung, -en} & Présentation : type de texte, thème, source, idée principale & l'introduction \\
\hline
\all{der Hauptteil, -e} & Analyse : structure du texte, idées secondaires, procédés (modaux, connecteurs, énumérations) avec citations courtes & le développement \\
\hline
\all{der Schluss, „-e} & Conclusion : bilan et opinion personnelle argumentée & la conclusion \\
\hline
\end{tabularx}
\end{center}

\begin{popfigure}
\begin{tikzpicture}[
  node distance=0.55cm and 1.1cm,
  box/.style={draw=popblue, thick, rounded corners=3pt, fill=popblueL, align=center, text width=4.2cm, minimum height=1.1cm, font=\small},
  sub/.style={draw=poporange, thick, rounded corners=3pt, fill=poporangeL, align=center, text width=4.6cm, minimum height=1.1cm, font=\small},
  fin/.style={draw=popgreen, thick, rounded corners=3pt, fill=popgreenL, align=center, text width=4.2cm, minimum height=1.1cm, font=\small},
  arr/.style={-{Stealth[length=3mm]}, thick, popdark}
]
\node[box, align=center] (e){\textbf{Einleitung}\\ \emph{Der Text handelt von...}\\ type, thème, idée principale};
\node[sub, below=of e, align=center] (h1){\textbf{Hauptteil 1 : die Ideen}\\ les libertés, leurs limites,\\ l'appel à l'action};
\node[sub, below=of h1, align=center] (h2){\textbf{Hauptteil 2 : die Stilmittel}\\ modaux, \emph{obwohl}, énumérations\\ + citations courtes „...“};
\node[fin, below=of h2, align=center] (s){\textbf{Schluss}\\ \emph{Meiner Meinung nach...}\\ avis personnel + comparaison};
\draw[arr] (e) -- (h1);
\draw[arr] (h1) -- (h2);
\draw[arr] (h2) -- (s);
\end{tikzpicture}
\legende{Organigramme du plan de commentaire : de l'introduction à l'avis personnel.}
\end{popfigure}

\courssub{Les expressions utiles du commentaire}

Voici la boîte à outils linguistique du candidat. Apprenez ces formules par cœur : elles garantissent un allemand correct et idiomatique.

\begin{center}
\begin{tabularx}{\linewidth}{|X|X|}
\hline
\rowcolor{popblueL} \textbf{Deutsch} & \textbf{Français} \\
\hline
\all{Der Text handelt von...} & Le texte traite de... \\
\hline
\all{Der Autor zeigt, dass...} & L'auteur montre que... \\
\hline
\all{Der Autor nennt vier Freiheiten.} & L'auteur cite quatre libertés. \\
\hline
\all{Am Ende fordert er die Bürger auf, ...} & À la fin, il invite les citoyens à... \\
\hline
\all{Meiner Meinung nach...} & À mon avis... \\
\hline
\all{Ich bin der Meinung, dass...} & Je suis d'avis que... \\
\hline
\all{Dieser Text ist sehr wichtig, denn...} & Ce texte est très important, car... \\
\hline
\end{tabularx}
\end{center}

\begin{exemplebox}[Exemples traduits]
\all{Der Autor zeigt, dass Freiheit Grenzen hat.} « L'auteur montre que la liberté a des limites. »

\all{Ich bin der Meinung, dass die Pressefreiheit sehr wichtig ist.} « Je suis d'avis que la liberté de la presse est très importante. »

\all{Der Text handelt von den Rechten und Pflichten der Bürger.} « Le texte traite des droits et des devoirs des citoyens. »

\all{Obwohl die Freiheiten Grenzen haben, sind sie ein kostbares Gut.} « Bien que les libertés aient des limites, elles sont un bien précieux. »
\end{exemplebox}

Remarquez la grammaire à l'œuvre : après \all{dass} et \all{obwohl}, le verbe conjugué va \textbf{à la fin} de la subordonnée (\all{... dass Freiheit Grenzen hat} ; \all{... dass die Pressefreiheit wichtig ist}). C'est exactement la règle de la concession étudiée dans la section précédente, réemployée ici dans une production personnelle.

\begin{attention}[Piège : l'inversion après « meiner Meinung nach »]
\all{Meiner Meinung nach} occupe la première position de la phrase : le verbe conjugué doit donc venir \textbf{en deuxième position}, avant le sujet.

Correct : \all{Meiner Meinung nach \textbf{ist} dieser Text wichtig.}

Faux : *\all{Meiner Meinung nach dieser Text ist wichtig.}

C'est la même inversion qu'après \all{trotzdem} : \all{Trotzdem \textbf{müssen} wir für unsere Rechte kämpfen.} En français, l'ordre ne change pas (« À mon avis, ce texte est important ») : ne calquez jamais l'ordre français !
\end{attention}

\courssub{Commentaire modèle rédigé}

Lisez maintenant un commentaire complet (8 à 10 lignes), tel qu'un excellent élève de Terminale pourrait le rédiger à l'examen.

\begin{textebox}[Commentaire modèle : „Freiheiten in einer Demokratie“]
Der Text „Freiheiten in einer Demokratie“ handelt von den wichtigen Rechten der Bürger. Der Autor nennt vier Freiheiten: die Meinungsfreiheit, die Pressefreiheit, die Religionsfreiheit und die Versammlungsfreiheit. Er zeigt aber auch, dass die Freiheiten Grenzen haben: man darf niemanden beleidigen. Am Ende fordert er die Bürger auf, für ihre Rechte zu kämpfen. Meiner Meinung nach ist dieser Text sehr wichtig, denn auch in Kamerun müssen die jungen Leute ihre Rechte kennen. Freiheit ist wirklich ein kostbares Gut.
\end{textebox}

\textbf{Traduction :} « Le texte "Les libertés dans une démocratie" traite des droits importants des citoyens. L'auteur cite quatre libertés : la liberté d'expression, la liberté de la presse, la liberté de religion et la liberté de réunion. Mais il montre aussi que les libertés ont des limites : on ne doit insulter personne. À la fin, il invite les citoyens à lutter pour leurs droits. À mon avis, ce texte est très important, car au Cameroun aussi, les jeunes doivent connaître leurs droits. La liberté est vraiment un bien précieux. »

\textbf{Glossaire :} \all{der Bürger, -} (le citoyen) ; \all{die Grenze, -n} (la limite, la frontière) ; \all{beleidigen} (insulter, offenser) ; \all{auffordern (+ zu + Inf.)} (inviter à, encourager à) ; \all{kämpfen für} (lutter pour) ; \all{kostbar} (précieux) ; \all{das Gut, die Güter} (le bien).

\courssub{Analyse du modèle : connecteurs et grammaire réemployés}

Décortiquons le modèle pour comprendre pourquoi il est réussi. Chaque phrase joue un rôle précis dans le plan :

\begin{center}
\begin{tabularx}{\linewidth}{|X|X|}
\hline
\rowcolor{popblueL} \textbf{Élément du modèle} & \textbf{Fonction} \\
\hline
\all{Der Text ... handelt von...} & Einleitung : présentation du thème \\
\hline
\all{Der Autor nennt vier Freiheiten: ...} & Hauptteil : énumération des idées \\
\hline
\all{Er zeigt \textbf{aber} auch, dass...} & connecteur d'opposition + subordonnée avec \all{dass} (verbe final) \\
\hline
\all{man \textbf{darf} niemanden beleidigen} & verbe de modalité (la règle et l'interdiction) \\
\hline
\all{fordert er die Bürger \textbf{auf}, ... zu kämpfen} & verbe à particule séparable \all{auffordern} \\
\hline
\all{\textbf{Meiner Meinung nach ist} dieser Text...} & Schluss : avis personnel avec inversion correcte \\
\hline
\all{\textbf{denn} auch in Kamerun müssen...} & connecteur de cause (verbe en 2e position) + modal \all{müssen} \\
\hline
\end{tabularx}
\end{center}

Observez les \cle{connecteurs} qui donnent du rythme au texte : \all{aber} (opposition), \all{denn} (cause), \all{außerdem} (addition), \all{deshalb} (conséquence), \all{meiner Meinung nach} (opinion). Un commentaire sans connecteurs est une suite de phrases isolées : le correcteur sanctionne ce manque de cohérence. Notez aussi que le modèle réemploie les deux points de grammaire de la leçon : les verbes de modalité (\all{darf, müssen}) et la logique de concession (l'idée « les libertés ont des limites, \emph{mais} elles restent précieuses », qu'on aurait pu exprimer avec \all{obwohl} : \all{Obwohl die Freiheiten Grenzen haben, sind sie kostbar.}).

\begin{exoresolu}[Transformer une phrase simple en phrase de commentaire]
Énoncé : Vous voulez dire dans votre commentaire : « L'auteur montre que les jeunes doivent connaître leurs droits. » Écrivez la phrase allemande correcte avec \all{zeigen, dass} et le verbe modal \all{müssen}.
\tcblower
Solution détaillée : la subordonnée introduite par \all{dass} place le verbe conjugué à la fin ; le modal \all{müssen} est conjugué et l'infinitif \all{kennen} se place juste avant lui.

\all{Der Autor zeigt, dass die jungen Leute ihre Rechte kennen müssen.}

Vérification : sujet \all{die jungen Leute} ; verbe conjugué \all{müssen} en dernière position ; infinitif \all{kennen} juste avant ; l'ordre « infinitif + modal » à la fin est la règle des subordonnées avec verbe de modalité. Erreur à éviter : *\all{... dass die jungen Leute müssen ihre Rechte kennen} (calque de l'ordre français « doivent connaître leurs droits »).
\end{exoresolu}

\begin{attention}[Les trois erreurs fatales du commentaire]
\begin{enumerate}
\item \textbf{Recopier le texte} : le correcteur veut vos phrases, pas celles de l'auteur. Ne citez que de courts extraits (3 à 6 mots) entre guillemets.
\item \textbf{Oublier l'opinion personnelle} : le \all{Schluss} avec \all{meiner Meinung nach} est obligatoire. Sans lui, le commentaire est incomplet.
\item \textbf{Phrases sans verbe conjugué en 2e position} : chaque phrase principale allemande doit avoir son verbe conjugué en deuxième position. Après \all{meiner Meinung nach}, \all{trotzdem}, \all{deshalb} ou \all{außerdem}, c'est le verbe qui suit immédiatement, puis le sujet : \all{Deshalb müssen wir kämpfen.} (et non *\all{Deshalb wir müssen kämpfen.}).
\end{enumerate}
\end{attention}

\begin{exercice}[À vous de commenter]
Rédigez un commentaire de 6 à 8 lignes sur le texte „Freiheiten in einer Demokratie“ en suivant le plan Einleitung – Hauptteil – Schluss. Consignes : utilisez au moins deux expressions de la boîte à outils (\all{Der Text handelt von...}, \all{Der Autor zeigt, dass...}), un verbe de modalité, une phrase avec \all{obwohl} ou \all{trotzdem}, et terminez par votre avis personnel comparant avec la situation des jeunes au Cameroun.
\end{exercice}

\begin{corrige}[Exercice : éléments de correction]
Un bon commentaire contient : une Einleitung (\all{Der Text handelt von den Freiheiten in einer Demokratie.}) ; un Hauptteil avec les quatre libertés et leurs limites, par exemple \all{Obwohl man seine Meinung sagen darf, darf man niemanden beleidigen.} ; un Schluss avec inversion correcte : \all{Meiner Meinung nach müssen auch die Jugendlichen in Kamerun ihre Rechte kennen.} Vérifiez : verbe conjugué en 2e position dans chaque phrase principale, verbe final après \all{dass} et \all{obwohl}, aucun long passage recopié.
\end{corrige}

Vous savez désormais analyser un texte sur les libertés et rédiger un commentaire structuré, avec connecteurs, modaux et avis personnel. La section suivante, « Production guidée écrite et orale », vous fera passer du commentaire à la prise de parole : débat et jeu de rôle sur les droits et devoirs des citoyens.

\coursec{Production guidée écrite et orale}

Après avoir appris à analyser et à commenter le texte „Freiheiten in einer Demokratie“, il est temps de passer à l'étape la plus active de la leçon : \cle{produire} à ton tour. Tu vas d'abord rédiger un court paragraphe en suivant une trame précise, puis participer à un débat oral en binôme. Ces deux activités réinvestissent tout ce que tu as appris dans cette leçon : le lexique des libertés, les verbes de modalité et les connecteurs de concession \all{obwohl} et \all{trotzdem}. C'est exactement le type de tâche qui t'attend au baccalauréat.

\courssub{Production écrite guidée : « Welche Freiheit ist für dich am wichtigsten? »}

\textbf{La consigne}

\begin{exercice}[Sujet de production écrite]
Rédige un paragraphe de 5 à 6 phrases en allemand répondant à la question : \all{Welche Freiheit ist für dich am wichtigsten? Warum?} (Quelle liberté est la plus importante pour toi ? Pourquoi ?)

Contraintes obligatoires :
\begin{itemize}
\item utiliser au moins \textbf{2 verbes de modalité} (\all{dürfen, müssen, können, sollen}) ;
\item écrire au moins \textbf{1 subordonnée avec} \all{obwohl} ;
\item écrire au moins \textbf{1 phrase avec} \all{trotzdem} ;
\item employer au moins \textbf{5 mots du lexique de la leçon} (par exemple : \all{die Freiheit, -en ; die Meinungsfreiheit, -en ; die Pressefreiheit, -en ; das Recht, -e ; die Pflicht, -en ; erlauben ; verbieten ; die Demokratie, -n ; die Religion, -en ; die Versammlung, -en}).
\end{itemize}
\end{exercice}

\textbf{La trame fournie}

Pour t'aider, voici la trame à suivre phrase par phrase. Chaque ligne correspond à une fonction précise du paragraphe :

\begin{center}
\begin{tabularx}{\linewidth}{|l|X|l|}
\hline
\rowcolor{popblueL} \textbf{Phrase} & \textbf{Trame allemande} & \textbf{Fonction} \\
\hline
1 & \all{Meiner Meinung nach ist die ... am wichtigsten, denn ...} & opinion + justification \\
\hline
2 & \all{Obwohl ..., dürfen wir ...} & concession + droit \\
\hline
3 & \all{Wir müssen ...} & devoir \\
\hline
4 & \all{Trotzdem ...} & conclusion nuancée \\
\hline
\end{tabularx}
\end{center}

\begin{popfigure}
\begin{tikzpicture}[node distance=0.6cm, every node/.style={font=\small}]
\node[draw=popblue, fill=popblueL, rounded corners, text width=0.9\linewidth, align=center] (p1) {\textbf{Phrase 1 : Opinion} \\ \all{Meiner Meinung nach ist die Meinungsfreiheit am wichtigsten, denn ...}};
\node[draw=popgreen, fill=popgreenL, rounded corners, text width=0.9\linewidth, align=center, below=of p1] (p2) {\textbf{Phrase 2 : Concession} \\ \all{Obwohl ..., dürfen wir unsere Ideen frei sagen.}};
\node[draw=poporange, fill=poporangeL, rounded corners, text width=0.9\linewidth, align=center, below=of p2] (p3) {\textbf{Phrase 3 : Devoir} \\ \all{Wir müssen die Meinung der anderen respektieren.}};
\node[draw=poppurple, fill=poppurpleL, rounded corners, text width=0.9\linewidth, align=center, below=of p3] (p4) {\textbf{Phrase 4 : Conclusion} \\ \all{Trotzdem hat jede Freiheit ihre Grenzen.}};
\draw[-{Stealth}, thick, popdark] (p1) -- (p2);
\draw[-{Stealth}, thick, popdark] (p2) -- (p3);
\draw[-{Stealth}, thick, popdark] (p3) -- (p4);
\end{tikzpicture}
\legende{La trame de production : de l'opinion à la conclusion en quatre étapes}
\end{popfigure}

\begin{methode}[Réussir la production guidée en quatre étapes]
\begin{enumerate}
\item \textbf{Suivre la trame phrase par phrase.} Ne change pas l'ordre : opinion, justification avec \all{denn}, concession avec \all{obwohl}, obligation avec \all{müssen}, conclusion avec \all{trotzdem}. La trame est ton squelette ; tu ajoutes seulement la chair (le contenu).
\item \textbf{Vérifier la place de chaque verbe.} Le modal conjugué se place en \textbf{2\textsuperscript{e} position} et l'infinitif \textbf{en fin de phrase} : \all{Wir \underline{müssen} die Rechte der anderen \underline{respektieren}.} Dans la subordonnée avec \all{obwohl}, le verbe conjugué va \textbf{à la fin} : \all{Obwohl alle Menschen frei \underline{sind}, ...}
\item \textbf{Compter les contraintes.} Avant de rendre ta copie, vérifie : 2 modaux ? 1 \all{obwohl} ? 1 \all{trotzdem} ? 5 mots du lexique ?
\item \textbf{Relire en soulignant} les mots du lexique utilisés et en corrigeant les majuscules des noms, les terminaisons verbales et l'orthographe.
\end{enumerate}
\end{methode}

\begin{exemplebox}[Modèle rédigé et traduit]
Voici un paragraphe complet qui respecte toutes les contraintes :

\all{Meiner Meinung nach ist die Meinungsfreiheit am wichtigsten, denn jeder Mensch muss seine Ideen sagen können. Obwohl diese Freiheit manchmal Konflikte schafft, dürfen wir in einer Demokratie unsere Meinung frei äußern. Wir müssen aber die Rechte und die Pflichten der anderen respektieren. Niemand darf andere Menschen beleidigen. Trotzdem ist die Meinungsfreiheit ein Grundrecht, das wir schützen müssen.}

« À mon avis, la liberté d'expression est la plus importante, car chacun doit pouvoir exprimer ses idées. Bien que cette liberté crée parfois des conflits, nous pouvons exprimer librement notre opinion dans une démocratie. Mais nous devons respecter les droits et les devoirs des autres. Personne ne doit insulter autrui. Néanmoins, la liberté d'expression est un droit fondamental que nous devons protéger. »

Vérification des contraintes : modaux \all{muss ... können}, \all{dürfen}, \all{müssen}, \all{darf}, \all{müssen} ; subordonnée \all{Obwohl diese Freiheit manchmal Konflikte schafft} ; phrase avec \all{Trotzdem} ; lexique : \all{Meinungsfreiheit, Demokratie, Rechte, Pflichten, Grundrecht}.
\end{exemplebox}

\begin{attention}[Les trois pièges de la trame]
\begin{itemize}
\item Après \all{denn}, l'ordre est \textbf{direct} (conjonction de coordination) : \all{..., denn jeder Mensch \underline{muss} seine Ideen sagen können.} Ne fais jamais passer le verbe à la fin après \all{denn} !
\item Après \all{obwohl}, le verbe conjugué va \textbf{à la fin} : \all{Obwohl diese Freiheit wichtig \underline{ist}, ...} et non \emph{Obwohl diese Freiheit ist wichtig}.
\item \all{Trotzdem} occupe la \textbf{1\textsuperscript{re} position} et provoque l'inversion : \all{Trotzdem \underline{ist} die Meinungsfreiheit ein Grundrecht.} (sujet après le verbe).
\end{itemize}
\end{attention}

\begin{exoresolu}[Compléter la trame]
Complète ce paragraphe avec les éléments qui manquent : \all{Meiner Meinung nach ist die Pressefreiheit am wichtigsten, denn die Journalisten \underline{\hspace{1.5cm}} (devoir-informer) die Bürger \underline{\hspace{1.5cm}}. Obwohl manche Zeitungen Fehler \underline{\hspace{1.5cm}} (faire), \underline{\hspace{1.5cm}} (pouvoir-nous) wir sie lesen. Wir \underline{\hspace{1.5cm}} (devoir) aber kritisch \underline{\hspace{1.5cm}} (rester). Trotzdem \underline{\hspace{1.5cm}} (être) die Pressefreiheit wichtig für die Demokratie.}
\tcblower
\all{Meiner Meinung nach ist die Pressefreiheit am wichtigsten, denn die Journalisten \textbf{müssen} die Bürger \textbf{informieren}. Obwohl manche Zeitungen Fehler \textbf{machen}, \textbf{dürfen} wir sie lesen. Wir \textbf{müssen} aber kritisch \textbf{bleiben}. Trotzdem \textbf{ist} die Pressefreiheit wichtig für die Demokratie.}

Points de grammaire vérifiés : modal en 2\textsuperscript{e} position + infinitif final (\all{müssen ... informieren}) ; verbe conjugué en fin de subordonnée (\all{Fehler machen}) ; inversion après \all{trotzdem} (\all{Trotzdem ist ...}).
\end{exoresolu}

\courssub{Production orale : débat en binômes « Darf man in der Schule alles sagen? »}

Passons maintenant à l'oral. Le débat est un exercice classique de la classe de Terminale : il entraîne la prise de parole spontanée, l'argumentation et l'écoute de l'autre.

\begin{experience}[Débat en binômes : Peut-on tout dire à l'école?]
\textbf{Sujet :} \all{Darf man in der Schule alles sagen?}

\textbf{Rôles :}
\begin{itemize}
\item \textbf{Élève A (pour)} : tu défends la liberté d'expression totale à l'école. Arguments possibles : \all{Die Meinungsfreiheit ist ein Recht. / Schüler müssen ihre Meinung sagen dürfen. / Man lernt durch Diskussionen.}
\item \textbf{Élève B (contre)} : tu penses qu'il faut des limites. Arguments possibles : \all{Man darf niemanden beleidigen. / Wir müssen die Lehrer respektieren. / Obwohl die Freiheit wichtig ist, hat sie Grenzen.}
\end{itemize}

\textbf{Déroulement :} 2 minutes de préparation, 3 minutes de débat, puis changement de rôle. Le professeur circule et note les structures utilisées.
\end{experience}

\textbf{Exprimer l'accord et le désaccord}

Pour débattre, tu as besoin de formules de réaction. Observe ce tableau et mémorise-le :

\begin{center}
\begin{tabularx}{\linewidth}{|X|X|X|}
\hline
\rowcolor{popblueL} \textbf{Deutsch} & \textbf{Français} & \textbf{Emploi} \\
\hline
\all{Ich stimme dir zu.} & Je suis d'accord avec toi. & accord fort \\
\hline
\all{Du hast Recht.} & Tu as raison. & accord \\
\hline
\all{Das stimmt, aber ...} & C'est vrai, mais ... & accord partiel \\
\hline
\all{Ich bin anderer Meinung.} & Je suis d'un autre avis. & désaccord poli \\
\hline
\all{Das sehe ich anders.} & Je vois cela différemment. & désaccord poli \\
\hline
\all{Ich stimme dir nicht zu, denn ...} & Je ne suis pas d'accord, car ... & désaccord + argument \\
\hline
\end{tabularx}
\end{center}

\begin{attention}[Du ou Sie?]
Dans le débat oral, emploie le \textbf{\all{du}} avec tes camarades : \all{Ich stimme \underline{dir} zu.} / \all{Was meinst \underline{du}?} Le \textbf{\all{Sie}} ne s'emploie qu'avec le professeur : \all{Ich bin anderer Meinung, Frau Lehrerin.} Erreur typique du francophone : dire \all{Ich stimme Sie zu} à un camarade — le pronom personnel au datif de \all{du} est \all{dir}, pas \all{Sie}!
\end{attention}

\begin{textebox}[Exemple de débat entre deux élèves]
\all{A: Meiner Meinung nach darf man in der Schule alles sagen, denn die Meinungsfreiheit ist ein Grundrecht.}

\all{B: Das sehe ich anders. Obwohl die Meinungsfreiheit wichtig ist, dürfen wir andere Schüler nicht beleidigen.}

\all{A: Das stimmt, aber Schüler müssen ihre Probleme offen besprechen können.}

\all{B: Ich stimme dir zu, aber es gibt Regeln. Wir müssen die Lehrer und die Mitschüler respektieren.}

\all{A: Du hast Recht. Trotzdem finde ich: Die Schule muss ein Ort der freien Diskussion bleiben.}

\all{B: Ja, aber mit Respekt. Freiheit bedeutet auch Verantwortung.}

\textbf{Glossaire :} \all{beleidigen} = insulter ; \all{offen besprechen} = discuter ouvertement ; \all{die Regel, -n} = la règle ; \all{der Mitschüler, -} = le camarade de classe ; \all{die Verantwortung} = la responsabilité ; \all{bedeuten} = signifier.
\end{textebox}

\courssub{Correction collective et grille d'évaluation}

Après les productions, la classe procède à une correction collective. Deux ou trois paragraphes sont lus au tableau, et deux binômes rejouent leur débat devant la classe. Voici la grille utilisée pour évaluer :

\begin{center}
\begin{tabularx}{\linewidth}{|X|l|X|}
\hline
\rowcolor{popblueL} \textbf{Critère} & \textbf{Points} & \textbf{Ce qui est vérifié} \\
\hline
Contenu et pertinence & 4 & réponse à la question, arguments clairs \\
\hline
Grammaire : modaux & 4 & place du modal (2\textsuperscript{e} position) et infinitif final \\
\hline
Grammaire : concession & 4 & verbe final après \all{obwohl}, inversion après \all{trotzdem} \\
\hline
Lexique de la leçon & 4 & au moins 5 mots du thème, articles corrects \\
\hline
Prononciation (oral) & 4 & clarté, intonation, « ch » et « r » corrects \\
\hline
\end{tabularx}
\end{center}

\begin{methode}[S'auto-évaluer avant la correction]
Avant de rendre ton texte ou de passer à l'oral, pose-toi quatre questions :
\begin{enumerate}
\item Ai-je répondu exactement à la question posée ?
\item Chaque modal est-il conjugué en 2\textsuperscript{e} position avec l'infinitif à la fin ?
\item Mon \all{obwohl} envoie-t-il bien le verbe à la fin, et mon \all{trotzdem} provoque-t-il l'inversion ?
\item Ai-je souligné mes 5 mots de lexique et vérifié les majuscules des noms ?
\end{enumerate}
\end{methode}

\begin{savaistu}[L'expression écrite au baccalauréat]
Au Bac, l'expression écrite est notée sur trois critères principaux : la \textbf{pertinence du contenu}, la \textbf{correction grammaticale} et la \textbf{richesse du vocabulaire}. Réutiliser les structures de la leçon — les verbes de modalité, \all{obwohl}, \all{trotzdem}, les formules d'opinion comme \all{Meiner Meinung nach} — fait gagner des points précieux, car le correcteur voit immédiatement que tu maîtrises le programme de Terminale. Un texte simple mais correct vaut toujours mieux qu'un texte ambitieux plein de fautes!
\end{savaistu}

\courssub{Réinvestissement : devoir à la maison}

\begin{exercice}[Devoir à la maison : résumé du texte support]
Rédige un résumé du texte „Freiheiten in einer Demokratie“ en \textbf{3 phrases} seulement. Consignes :
\begin{itemize}
\item phrase 1 : le thème général du texte (\all{Der Text handelt von ...}) ;
\item phrase 2 : les libertés présentées (utilise \all{die Meinungsfreiheit, die Pressefreiheit, die Religionsfreiheit, die Versammlungsfreiheit}) ;
\item phrase 3 : l'idée principale de la conclusion (utilise un modal ou \all{obwohl}/\all{trotzdem}).
\end{itemize}
\end{exercice}

\begin{corrige}[Devoir à la maison — modèle de résumé]
\all{Der Text handelt von den wichtigen Freiheiten in einer Demokratie. Er erklärt die Meinungsfreiheit, die Pressefreiheit, die Religionsfreiheit und die Versammlungsfreiheit. Obwohl diese Freiheiten Grundrechte sind, müssen die Bürger auch ihre Pflichten kennen.}

« Le texte traite des libertés importantes dans une démocratie. Il explique la liberté d'expression, la liberté de la presse, la liberté de religion et la liberté de réunion. Bien que ces libertés soient des droits fondamentaux, les citoyens doivent aussi connaître leurs devoirs. »
\end{corrige}

\begin{aretenir}[L'essentiel de la production guidée]
\begin{itemize}
\item Une production guidée réussie suit la \textbf{trame} à la lettre et respecte toutes les \textbf{contraintes} (modaux, \all{obwohl}, \all{trotzdem}, lexique).
\item Modal conjugué en 2\textsuperscript{e} position, infinitif en fin ; verbe final après \all{obwohl} ; inversion après \all{trotzdem}.
\item À l'oral : \all{du} avec les camarades, \all{Sie} avec le professeur ; réagis avec \all{Ich stimme dir zu.} / \all{Das sehe ich anders.}
\item Relire, souligner le lexique, vérifier les majuscules : c'est la signature d'un bon élève de Terminale.
\end{itemize}
\end{aretenir}

\newpage

\coursec{Activité d'intégration}

\coursec{Production guidée écrite et orale : Activité d'intégration}

\courssub{Objectifs de l'activité}

À l'issue de cette activité d'intégration, l'élève sera capable de :
\begin{itemize}
\item \cle{réinvestir} le lexique thématique des libertés et des devoirs (\all{die Freiheit, die Meinungsfreiheit, die Pressefreiheit, die Religionsfreiheit, die Versammlungsfreiheit, das Recht, die Pflicht, erlauben, verbieten}) dans une production écrite authentique ;
\item \cle{employer correctement} les verbes de modalité (\all{dürfen, müssen, können, sollen}) à la deuxième place de la phrase déclarative, avec l'infinitif renvoyé en fin de proposition ;
\item \cle{articuler} ses idées grâce aux connecteurs de concession \all{obwohl} (subordonnée, verbe en fin) et \all{trotzdem} (phrase principale, verbe en deuxième position) ;
\item \cle{rédiger} en groupe un court article de journal scolaire (8 à 10 lignes) respectant les contraintes du genre journalistique : titre, accroche, développement, appel à l'action ;
\item \cle{présenter} oralement sa production en deux minutes et \cle{évaluer} les productions des autres groupes selon des critères simples.
\end{itemize}

\courssub{Situation}

Le lycée bilingue de Yaoundé prépare la \all{Woche der Demokratie} (la Semaine de la démocratie), organisée chaque année au mois de mars. À cette occasion, le club d'allemand lance le premier numéro de son journal scolaire germanophone, intitulé \all{„Unsere Rechte, unsere Pflichten“} (« Nos droits, nos devoirs »). Le journal sera affiché au mur de la salle d'allemand et distribué aux classes de Seconde et de Première. Le rédacteur en chef du club a publié l'appel à contributions suivant :

\begin{textebox}[Appel à contributions du club d'allemand]
\textbf{Liebe Schülerinnen und Schüler der Terminale,}\par
unser Deutschclub sucht Artikel für die erste Ausgabe unserer Schülerzeitung „Unsere Rechte, unsere Pflichten“. Schreibt einen kurzen Artikel (8 bis 10 Zeilen) über eine Freiheit: Meinungsfreiheit, Pressefreiheit, Religionsfreiheit oder Versammlungsfreiheit. Erklärt, was wir Jugendliche in Kamerun dürfen und was wir müssen. Vergesst nicht: mindestens ein Satz mit \emph{obwohl} und ein Satz mit \emph{trotzdem}! Beendet euren Artikel mit einem Appell an die Leser. Die besten Artikel hängen wir an der Wand des Deutschsaals auf. Viel Erfolg!\par
\emph{Euer Redaktionsteam}
\end{textebox}

\textbf{Glossaire :} \all{die Ausgabe, -n} (le numéro, l'édition) ; \all{die Schülerzeitung, -en} (le journal scolaire) ; \all{die Zeile, -n} (la ligne) ; \all{vergessen (vergisst, vergaß, hat vergessen)} (oublier) ; \all{aufhängen} (accrocher, afficher) ; \all{der Deutschsaal, die Deutschsäle} (la salle d'allemand).

Votre classe de Terminale A relève le défi. Par groupes de quatre, vous rédigez un article, vous l'affichez au mur de la classe, puis vous le présentez oralement en deux minutes. La classe vote ensuite pour le meilleur article, qui sera envoyé au rédacteur en chef du club.

\courssub{Consignes}

\textbf{Étape 1 — Compréhension et choix du sujet (10 min).} Lisez attentivement l'appel à contributions. En groupe, choisissez \cle{une seule} liberté parmi les quatre proposées : \all{die Meinungsfreiheit}, \all{die Pressefreiheit}, \all{die Religionsfreiheit} ou \all{die Versammlungsfreiheit}. Notez trois idées que vous voulez défendre sur cette liberté dans votre lycée ou au Cameroun.

\textbf{Étape 2 — Remue-méninges lexical (10 min).} Complétez collectivement une fiche de vocabulaire : pour votre liberté, trouvez au moins cinq mots ou expressions de la leçon (avec article et pluriel pour les noms) et deux verbes de modalité adaptés. Exemple : pour la liberté d'expression, on peut noter \all{die Meinung, -en ; die Meinung äußern ; erlauben ; verbieten ; man darf...}.

\textbf{Étape 3 — Rédaction de l'article (20 min).} Rédigez votre article de 8 à 10 lignes en respectant la structure suivante :
\begin{enumerate}
\item un \cle{titre} court et accrocheur ;
\item une \cle{phrase d'introduction} qui présente la liberté choisie ;
\item un \cle{développement} : ce que les élèves \all{dürfen} (ont le droit de faire) et ce qu'ils \all{müssen} ou \all{sollen} (doivent faire) ;
\item au moins une phrase avec \all{obwohl} et une phrase avec \all{trotzdem} ;
\item un \cle{appel à l'action} final, par exemple : \all{Wir müssen für unsere Rechte kämpfen!} (« Nous devons nous battre pour nos droits ! »).
\end{enumerate}

\textbf{Étape 4 — Relecture grammaticale (10 min).} Vérifiez dans votre texte : la place du verbe conjugué (deuxième position), l'infinitif en fin de phrase avec les verbes de modalité, le verbe en fin de subordonnée après \all{obwohl}, la majuscule à tous les noms, les Umlaute (\all{ä, ö, ü}).

\textbf{Étape 5 — Présentation orale et vote (15 min).} Affichez votre article au mur. Chaque groupe présente le sien en deux minutes : un élève lit le titre et l'introduction, un autre le développement, un troisième l'appel à l'action. La classe vote selon trois critères : richesse du vocabulaire, correction grammaticale, force de l'appel à l'action.

\textbf{Variante rapide (si le temps manque) :} rédigez en groupe une affiche avec un slogan, par exemple \all{Freiheit ist ein kostbares Gut!} (« La liberté est un bien précieux ! »), accompagné de deux phrases avec un verbe de modalité et d'une phrase avec \all{obwohl}.

\courssub{Ressources à mobiliser}

\begin{aretenir}[Lexique utile pour l'article]
\begin{itemize}
\item \all{die Freiheit, -en} (la liberté) ; \all{die Meinungsfreiheit} (la liberté d'expression) ; \all{die Pressefreiheit} (la liberté de la presse) ; \all{die Religionsfreiheit} (la liberté de religion) ; \all{die Versammlungsfreiheit} (la liberté de réunion).
\item \all{das Recht, -e} (le droit) ; \all{die Pflicht, -en} (le devoir) ; \all{der Bürger, -} (le citoyen) ; \all{die Demokratie, -n} (la démocratie).
\item \all{erlauben} (autoriser), \all{verbieten (verbietet, verbot, hat verboten)} (interdire), \all{die Meinung äußern} (exprimer son opinion), \all{respektieren} (respecter), \all{kämpfen für + Akkusativ} (se battre pour).
\end{itemize}
\end{aretenir}

\begin{propriete}[Rappel : les verbes de modalité]
Le verbe de modalité conjugué occupe la \cle{deuxième position} ; l'infinitif du verbe principal est renvoyé \cle{en fin de phrase} : \all{Wir dürfen unsere Meinung frei äußern.} (« Nous avons le droit d'exprimer librement notre opinion. »).

Conjugaison au présent :
\begin{center}
\begin{tabular}{|l|l|l|}
\hline
\rowcolor{popblueL} \textbf{Personne} & \textbf{dürfen} & \textbf{müssen} \\ \hline
ich & darf & muss \\ \hline
du & darfst & musst \\ \hline
er/sie/es & darf & muss \\ \hline
wir & dürfen & müssen \\ \hline
ihr & dürft & müsst \\ \hline
sie/Sie & dürfen & müssen \\ \hline
\end{tabular}
\end{center}
\emph{Remarque :} \all{können} (pouvoir, savoir) et \all{sollen} (devoir, conseil) suivent le même schéma (pas de \emph{-t} à la 3\textsuperscript{e} pers. sg. pour \all{kann}, \all{soll}).
\end{propriete}

\begin{propriete}[Rappel : la concession]
\all{obwohl} introduit une \cle{subordonnée} : le verbe conjugué va \cle{en fin de proposition} : \all{Obwohl die Meinungsfreiheit wichtig ist, müssen wir andere respektieren.} (« Bien que la liberté d'expression soit importante, nous devons respecter les autres. »).

\all{trotzdem} est un \cle{adverbe} (conjonctif) : il occupe la première position (topique) et le verbe reste en deuxième position : \all{Die Freiheit ist begrenzt, trotzdem ist sie ein kostbares Gut.} (« La liberté est limitée, pourtant c'est un bien précieux. »).
\end{propriete}

\begin{attention}[Pièges fréquents]
\begin{itemize}
\item Ne dites jamais \all{*obwohl... trotzdem} dans la même phrase : les deux connecteurs expriment la même concession, il faut choisir l'un ou l'autre.
\item Attention à la fausse amie : \all{die Freiheit} signifie « la liberté », pas « la fraîcheur » (\all{die Frische}).
\item Après \all{dürfen} et \all{müssen}, l'infinitif se place \cle{toujours en toute fin de phrase} (derrière les compléments) : \all{Wir müssen für unsere Rechte kämpfen}, et non \all{*Wir müssen kämpfen für unsere Rechte}.
\item Majuscule obligatoire à \cle{tous les noms} : \all{die Meinungsfreiheit, das Grundrecht, die Grenze}.
\end{itemize}
\end{attention}

\courssub{Proposition de production et corrigé commenté}

\begin{exoresolu}[Article modèle : la liberté d'expression]
\textbf{Énoncé.} Rédigez un article de 8 à 10 lignes sur la liberté d'expression pour le journal \all{„Unsere Rechte, unsere Pflichten“}, avec au moins deux verbes de modalité, une phrase avec \all{obwohl}, une avec \all{trotzdem} et un appel à l'action final.
\tcblower
\textbf{Production modèle :}\par
\textbf{Reden ist Silber, Schweigen ist Gold? Nein: Reden ist ein Recht!}\par
\all{Die Meinungsfreiheit ist ein wichtiges Recht in einer Demokratie. Auch wir Schüler in Kamerun haben das Recht, unsere Meinung frei zu äußern: in der Klasse, in der Schülerzeitung und im Internet. Wir dürfen zum Beispiel über die Probleme unseres Gymnasiums diskutieren und Kritik äußern. Aber wir müssen auch die Meinung der anderen respektieren und wir dürfen niemanden beleidigen. Obwohl die Meinungsfreiheit ein Grundrecht ist, hat sie Grenzen. Trotzdem dürfen wir nicht schweigen, wenn wir Unrecht sehen. Liebe Mitschüler, nutzt eure Stimme: Wir müssen für unsere Rechte kämpfen!}\par
\textbf{Traduction :} « Parler est d'argent, se taire est d'or ? Non : parler est un droit ! La liberté d'expression est un droit important dans une démocratie. Nous aussi, les élèves du Cameroun, nous avons le droit d'exprimer librement notre opinion : en classe, dans le journal scolaire et sur Internet. Nous pouvons par exemple discuter des problèmes de notre lycée et exprimer des critiques. Mais nous devons aussi respecter l'opinion des autres et nous n'avons pas le droit d'insulter personne. Bien que la liberté d'expression soit un droit fondamental, elle a des limites. Pourtant, nous ne devons pas nous taire quand nous voyons une injustice. Chers camarades, utilisez votre voix : nous devons nous battre pour nos droits ! »\par
\textbf{Commentaire des choix de langue :}
\begin{itemize}
\item Le titre détourne le proverbe allemand \all{Reden ist Silber, Schweigen ist Gold} (« Parler est d'argent, se taire est d'or ») : c'est une accroche journalistique typique qui capte l'attention du lecteur.
\item Les verbes de modalité sont correctement placés : \all{wir dürfen... diskutieren und äußern}, \all{wir müssen... respektieren}, avec les infinitifs en fin de phrase.
\item La subordonnée de concession respecte la place du verbe : \all{Obwohl die Meinungsfreiheit ein Grundrecht ist} — le verbe \all{ist} est en fin de proposition, et la principale qui suit commence par le verbe (\all{hat}).
\item \all{trotzdem} occupe la première position et le verbe \all{dürfen} reste en deuxième position : \all{Trotzdem dürfen wir nicht schweigen}.
\item Le lexique de la leçon est réinvesti : \all{das Recht, die Meinung äußern, respektieren, beleidigen, das Grundrecht, die Grenze, kämpfen für}.
\item L'appel à l'action final s'adresse directement aux lecteurs (\all{Liebe Mitschüler}), ce qui est conforme au genre de l'article de journal scolaire.
\end{itemize}
\end{exoresolu}

\courssub{Bilan de l'activité}

Cette activité d'intégration a permis de mobiliser simultanément les quatre savoir-faire de la leçon dans une tâche authentique et motivante : la \cle{compréhension écrite} (lecture de l'appel à contributions), le \cle{lexique thématique} (droits, libertés et devoirs), la \cle{grammaire} (verbes de modalité et connecteurs de concession) et la \cle{production écrite et orale} (article puis présentation). Le travail en groupe a favorisé la négociation de sens et la correction collective, tandis que le vote final a développé l'esprit critique des élèves. Pour prolonger l'activité à la maison, chaque élève peut rédiger individuellement un second article sur une autre liberté (presse, religion ou réunion) en suivant la même grille de relecture : place du verbe, infinitif final, majuscules des noms, connecteurs de concession. Les meilleurs textes pourront alimenter le prochain numéro du journal du club d'allemand et être présentés lors de la \all{Woche der Demokratie}.

\newpage

\coursec{Méthodes et exercices résolus}

\courssub{Méthode 1 : Comprendre un texte authentique sur les libertés fondamentales}

\begin{methode}[Compréhension écrite : démarche en trois lectures]
\textbf{Objectif :} Extraire l'information globale, détaillée et l'intention de l'auteur d'un texte de presse ou d'essai sur la démocratie.
\begin{enumerate}
    \item \textbf{1re lecture (globale) :} Identifier le \cle{type de texte} (article, extrait de loi, éditorial), la \cle{source}, le \cle{thème général} (ici : les libertés en démocratie) et la \cle{thèse} défendue. Repérer les \cle{mots-clés} récurrents (Freiheit, Recht, Staat, Bürger).
    \item \textbf{2e lecture (détaillée) :} Survoler les questions. Chercher les \cle{informations explicites} (chiffres, noms, définitions données dans le texte). Souligner les \cle{connecteurs logiques} (deshalb, jedoch, zwar... aber) pour suivre l'argumentation. Noter le \cle{vocabulaire inconnu} déductible du contexte.
    \item \textbf{3e lecture (vérification) :} Relire les passages correspondant aux questions. Répondre \textbf{en allemand} par des phrases complètes ou des éléments de texte (citation avec guillemets „…“). Vérifier l'accord des temps (souvent \cle{Präteritum} pour le récit, \cle{Präsens} pour les vérités générales).
\end{enumerate}
\textbf{Reflexe Bac :} Ne jamais recopier de longs paragraphes. Citer le numéro de ligne (Z. 3--5). Utiliser ses propres mots (Umformulierung) pour montrer la compréhension.
\end{methode}

\begin{exoresolu}[Compréhension détaillée : « Freiheiten in einer Demokratie »]
\textbf{Énoncé :} Lisez le texte suivant (extrait simulé d'un manuel civique allemand) et répondez aux questions en allemand.

\begin{textebox}[Texte support : Freiheiten in einer Demokratie]
In einer Demokratie sind Grundrechte nicht verhandelbar. Die Meinungsfreiheit (Artikel 5 GG) garantiert jedem Bürger das Recht, seine Meinung frei zu äußern und zu verbreiten. Doch diese Freiheit findet ihre Grenze dort, wo die Rechte anderer verletzt werden. Hassrede und Hetze sind verboten. Ebenso schützt die Pressefreiheit die Medien vor staatlicher Zensur; Journalisten dürfen Missstände aufdecken, ohne Angst vor Repressionen zu haben. Die Religionsfreiheit (Artikel 4 GG) sichert das Recht, einer Glaubensgemeinschaft anzugehören oder keiner. Schließlich erlaubt die Versammlungsfreiheit (Artikel 8 GG), sich friedlich und ohne Waffen zu versammeln. Wer diese Rechte nutzt, trägt Verantwortung für den gesellschaftlichen Zusammenhalt.
\end{textebox}

\textbf{Fragen :}
\begin{enumerate}
    \item Was garantiert Artikel 5 des Grundgesetzes (GG)?
    \item Wo findet die Meinungsfreiheit ihre Grenze?
    \item Welche Rolle spielt die Pressefreiheit für Journalisten?
    \item Nennen Sie zwei Bedingungen für die Ausübung der Versammlungsfreiheit laut Text.
    \item Erklären Sie den letzten Satz mit eigenen Worten: „Wer diese Rechte nutzt, trägt Verantwortung...“
\end{enumerate}

\tcblower
\textbf{Solution détaillée :}

\textbf{Analyse grammaticale et lexicale :}
\begin{itemize}
    \item \all{Grundrechte} (n. pl., neutre) : droits fondamentaux. \all{nicht verhandelbar} : non négociables (adjectif + \cle{nicht}).
    \item \all{garantiert} (verbe \all{garantieren}, Präteritum 3e pers. sg. ou Partizip II) : garantit. Sujet : \all{die Meinungsfreiheit} (féminin, sg.).
    \item \all{äußern und zu verbreiten} : infinitifs (verbes à particule séparable : \all{sich äußern}, \all{verbreiten}) dépendant de \all{das Recht}.
    \item \all{findet... ihre Grenze} : verbe \all{finden} (irrégulier : findet, fand, gefunden) + possessif \all{ihr} (féminin, accusatif) accordé à \all{die Grenze}.
    \item \all{verletzt werden} : \cle{Passiv présent} (\all{werden} + Partizip II \all{verletzt}) = sont violés.
    \item \all{Hassrede und Hetze} (féminins) : discours de haine / incitation à la haine. \all{sind verboten} : Passif présent.
    \item \all{aufdecken} (séparable : \all{deckt auf}) : révéler. \all{ohne Angst... zu haben} : construction \all{ohne... zu} + infinitif.
    \item \all{sichert} : sécurise/garantit. \all{anzugehören} (séparable, datif : \all{einer Glaubensgemeinschaft}) : appartenir à.
    \item \all{erlaubt} : permet. \all{sich... zu versammeln} (réfléchi, séparable) : se réunir. \all{friedlich und ohne Waffen} : adverbes/préposition.
    \item \all{trägt Verantwortung} : porte la responsabilité (collocation : \all{Verantwortung tragen}).
\end{itemize}

\textbf{Réponses rédigées (niveau Bac) :}
\begin{enumerate}
    \item \all{Artikel 5 des Grundgesetzes garantiert die Meinungsfreiheit. Jeder Bürger hat das Recht, seine Meinung frei zu äußern und zu verbreiten.}
    \item \all{Die Meinungsfreiheit findet ihre Grenze dort, wo die Rechte anderer verletzt werden. Hassrede und Hetze sind verboten.}
    \item \all{Die Pressefreiheit schützt die Medien vor staatlicher Zensur. Journalisten dürfen Missstände aufdecken, ohne Angst vor Repressionen zu haben.}
    \item \all{Die Versammlung muss friedlich sein und ohne Waffen stattfinden.}
    \item \all{Der letzte Satz bedeutet, dass die Ausübung von Rechten Pflichten mit sich bringt. Wer seine Freiheiten nutzt, muss darauf achten, den gesellschaftlichen Zusammenhalt nicht zu gefährden.}
\end{enumerate}
\textbf{Conclusion :} Les réponses reprennent le vocabulaire du texte (Meinungsfreiheit, Zensur, Versammlungsfreiheit) et transforment la structure (passif $\to$ actif, nominalisation $\to$ verbe).
\end{exoresolu}

\courssub{Méthode 2 : Maîtriser le lexique thématique (droits, libertés, devoirs)}

\begin{methode}[Vocabulaire : familles de mots, articles, pluriels, collocations]
\textbf{Objectif :} Construire un répertoire lexical actif et orthographié correctement (majuscules, Umlaute, genres).
\begin{enumerate}
    \item \textbf{Tableau de dérivation :} Pour chaque nom central, noter le \cle{verbe}, l'\cle{adjectif}, l'\cle{adverbe}, le \cle{genre}, le \cle{pluriel} et la \cle{préposition} régie.
    \item \textbf{Collocations (Blocs figés) :} Apprendre les couples \textbf{verbe + nom} ou \textbf{adjectif + nom} (ex. \all{ein Recht haben}, \all{ein Verbot aussprechen}, \all{grundlegend für}). Cela évite les calques du français (*\all{machen ein Recht}).
    \item \textbf{Faux amis / Pièges :} \all{die Freiheit} (liberté) $\neq$ \all{die Freiwilligkeit} (volontariat). \all{erlauben} (permettre, + datif pers. + infinitif) vs \all{verbieten} (interdire, + datif pers. + infinitif). \all{das Recht} (le droit / la raison) vs \all{das Gesetz} (la loi).
    \item \textbf{Entraînement :} Transformer des phrases : Nominalisation (\all{verbieten} $\to$ \all{das Verbot}), Verbalisation (\all{in Freiheit leben} $\to$ \all{frei sein}), Synonymes (\all{erlauben} $\leftrightarrow$ \all{gestatten}, \all{verbieten} $\leftrightarrow$ \all{untersagen}).
\end{enumerate}
\textbf{Structure à retenir :} \all{jemandem (Dat) etwas (Akk) erlauben / verbieten / gestatten / untersagen}.
\end{methode}

\begin{exoresolu}[Exercices lexicaux : dérivation, collocations, thème]
\textbf{Énoncé :}
\begin{enumerate}
    \item Complétez le tableau de dérivation. Donnez le genre, le pluriel (noms) et la préposition si nécessaire.
    \begin{center}
    \begin{tabularx}{\linewidth}{|X|X|X|X|X|}\hline
    \rowcolor{popblueL} \textbf{Nom (Art. Pl.)} & \textbf{Verbe} & \textbf{Adjectif} & \textbf{Adverbe} & \textbf{Français} \\ \hline
    \all{die Freiheit, die Freiheiten} & \all{freien / befreien} & \all{frei} & \all{frei} & la liberté / libre \\ \hline
    \all{das Recht, die Rechte} &  &  &  & le droit / juste \\ \hline
     & \all{verbieten} &  &  & interdire \\ \hline
     & \all{erlauben} &  &  & permettre \\ \hline
    \all{die Pflicht, die Pflichten} &  &  &  & le devoir / obligatoire \\ \hline
    \all{die Versammlung, die Versammlungen} &  &  &  & la réunion / assemblée \\ \hline
    \end{tabularx}
    \end{center}
    \item Choisissez le bon verbe (\all{erlauben}, \all{verbieten}, \all{gestatten}, \all{untersagen}) et mettez-le à la forme correcte (Präsens).
    \begin{enumerate}
        \item Der Lehrer \_\_\_\_\_\_\_\_\_ den Schülern, das Handy zu benutzen. (interdit)
        \item Die Verfassung \_\_\_\_\_\_\_\_\_ jedem Bürger, seine Meinung zu sagen. (garantit/permet formel)
        \item Die Polizei \_\_\_\_\_\_\_\_\_ die Demonstration, weil sie gewalttätig wurde. (a interdit)
        \item Man \_\_\_\_\_\_\_\_\_ hier nicht rauchen. (il est interdit de / on ne permet pas)
    \end{enumerate}
    \item \textbf{Thème (Français $\to$ Allemand) :} Traduisez en utilisant le vocabulaire du tableau.
    \begin{enumerate}
        \item La liberté de la presse est un droit fondamental.
        \item Les citoyens ont le devoir de respecter les lois.
        \item Il est interdit de crier dans la bibliothèque.
        \item La réunion a été interdite par la police.
    \end{enumerate}
\end{enumerate}

\tcblower
\textbf{Solution détaillée :}

\textbf{1. Tableau complété :}
\begin{center}
\begin{tabularx}{\linewidth}{|X|X|X|X|X|}\hline
\rowcolor{popblueL} \textbf{Nom (Art. Pl.)} & \textbf{Verbe} & \textbf{Adjectif} & \textbf{Adverbe} & \textbf{Français} \\ \hline
\all{die Freiheit, die Freiheiten} & \all{freien / befreien} & \all{frei} & \all{frei} & la liberté / libre \\ \hline
\all{das Recht, die Rechte} & \all{recht haben / berechtigen} & \all{rechtmäßig / rechtlich} & \all{rechtlich} & le droit / juste \\ \hline
\all{das Verbot, die Verbote} & \all{verbieten} & \all{verboten} & \all{verboten} & l'interdiction / interdit \\ \hline
\all{die Erlaubnis, die Erlaubnisse} & \all{erlauben} & \all{erlaubt} & \all{erlaubterweise} & la permission / permis \\ \hline
\all{die Pflicht, die Pflichten} & \all{pflichten (selten)} & \all{pflichtig / pflichtbewusst} & \all{pflichtgemäß} & le devoir / obligatoire \\ \hline
\all{die Versammlung, die Versammlungen} & \all{sich versammeln} & \all{versammelt} &  & la réunion / assemblée \\ \hline
\end{tabularx}
\end{center}
\textit{Notes :} \all{recht haben} (avoir raison) $\neq$ \all{berechtigt sein} (être en droit). \all{das Recht} (droit/justice) vs \all{das Gesetz} (loi). \all{pflichtig} (tenu de, + Génitif/Datif).

\textbf{2. Choix des verbes (Präsens) :}
\begin{enumerate}
    \item \all{verbietet} (\all{verbieten} + Dat \all{den Schülern} + Infinitif \all{zu benutzen}). \textit{Ordre : verbe en 2e position, infinitif à la fin.}
    \item \all{gestattet} ou \all{garantiert}. \all{gestatten} est plus formel/administratif (\all{die Verfassung gestattet}). \all{erlauben} possible mais moins fort pour une constitution.
    \item \all{untersagt} (\all{untersagen} = interdire formellement, + Dat \all{der Demonstration} + Infinitif). \all{verboten} (Partizip II) ne convient pas seul comme verbe conjugué.
    \item \all{darf... nicht} ou \all{ist... untersagt/verboten}. \textbf{Meilleure réponse :} \all{Man darf hier nicht rauchen.} (Impersonnel \all{man} + modal \all{dürfen} + \all{nicht}). Ou passif : \all{Hier ist das Rauchen verboten.}
\end{enumerate}

\textbf{3. Thème (Traduction) :}
\begin{enumerate}
    \item \all{Die Pressefreiheit ist ein Grundrecht.} (\all{Pressefreiheit} f., \all{Grundrecht} n.).
    \item \all{Die Bürger haben die Pflicht, die Gesetze zu respektieren.} (\all{die Pflicht haben} + Infinitif \all{zu respektieren}. \all{respektieren} / \all{achten} / \all{einhalten}).
    \item \all{In der Bibliothek ist es verboten, zu schreien.} (Structure impersonnelle \all{es ist verboten} + \all{zu}-Infinitif) ou \all{Man darf in der Bibliothek nicht schreien.}
    \item \all{Die Versammlung wurde von der Polizei verboten.} (\cle{Passif procès} \all{werden} + Partizip II \all{verboten} + \all{von} + Dat \all{der Polizei}). \textit{Alternatif :} \all{Die Polizei hat die Versammlung verboten.} (Actif, Perfekt).
\end{enumerate}
\textbf{Conclusion :} La maîtrise du genre (\all{die Freiheit} vs \all{das Recht}), du pluriel (\all{-en}, \all{-e}, \all{-en}, \all{Umlaut}) et des régimes (\all{verbieten jdm etw}) est indispensable pour le thème.
\end{exoresolu}

\courssub{Méthode 3 : Conjuguer et employer les six verbes de modalité}

\begin{methode}[Les verbes de modalité (Modalverben) : Formen, Bedeutungen, Syntax]
\textbf{Objectif :} Maîtriser \all{können, müssen, dürfen, sollen, wollen, möchten} aux temps du programme (Präsens, Präteritum, Perfekt, Futur I, Konjunktiv II).
\begin{enumerate}
    \item \textbf{Conjugaison (Präsens) :} Apprendre par cœur les \cle{irrégularités du singulier} (changement de voyelle pour \all{können, müssen, dürfen, sollen, wollen}) et la \cle{régularité du pluriel} (terminaisons \all{-en, -t, -en}). \all{möchten} (Konjunktiv II de \all{mögen}) se conjugue comme un verbe régulier au présent.
    \item \textbf{Syntaxe (Position) :} Verbe modal \textbf{conjugué en position 2} (phrase principale) ou \textbf{en position finale} (subordonnée). L'\textbf{infinitif du verbe principal à la fin de la phrase} (klammer structure).
    \item \textbf{Sens de base vs Sens épistémique :} Distinguer la valeur \textbf{déontique} (action : capacité, obligation, permission, volonté) de la valeur \textbf{épistémique} (degré de certitude : \all{Er muss krank sein} = Il doit être malade).
    \item \textbf{Perfekt :} \textbf{Double infinitif} (\all{Er hat arbeiten müssen}) auxiliaire \all{haben}. Pas de \all{ge-} devant le modal.
    \item \textbf{Négation :} \all{nicht müssen} = n'avoir pas l'obligation (pas devoir) $\neq$ \all{nicht dürfen} = interdiction (ne pas avoir le droit). \all{nicht sollen} = ne pas devoir (recommandation négative).
\end{enumerate}
\textbf{Tableau récapitulatif (Präsens / Präteritum / Partizip II) à connaître par cœur.}
\end{methode}

\begin{exoresolu}[Grammaire : Verbes de modalité (Conjugaison, Syntaxe, Sens, Thème)]
\textbf{Énoncé :}
\begin{enumerate}
    \item Conjuguez les verbes entre parenthèses au temps indiqué.
    \begin{enumerate}
        \item Ich \_\_\_\_\_\_\_\_\_ (können, Präteritum) gestern nicht kommen.
        \item Wir \_\_\_\_\_\_\_\_\_ (müssen, Perfekt) die Hausaufgaben machen.
        \item Die Schüler \_\_\_\_\_\_\_\_\_ (sollen, Futur I) mehr lesen.
        \item Er \_\_\_\_\_\_\_\_\_ (wollen, Konjunktiv II) gerne Arzt werden.
    \end{enumerate}
    \item Transformez les phrases à la forme indiquée (garde le sens).
    \begin{enumerate}
        \item \all{Du darfst hier nicht parken.} (Négation avec \textit{müssen}) $\to$ \_\_\_\_\_\_\_\_\_
        \item \all{Er kann gut Deutsch sprechen.} (Subordonnée avec \textit{weil}) $\to$ \_\_\_\_\_\_\_\_\_
        \item \all{Wir wollen das Buch lesen.} (Perfekt) $\to$ \_\_\_\_\_\_\_\_\_
    \end{enumerate}
    \item \textbf{Thème :} Traduisez en allemand.
    \begin{enumerate}
        \item Nous devons (obligation intérieure) apprendre le vocabulaire.
        \item Tu n'as pas le droit d'utiliser ton portable en classe.
        \item Il a pu (a réussi à) résoudre le problème.
        \item Vous devriez (conseil) aller chez le médecin.
    \end{enumerate}
\end{enumerate}

\tcblower
\textbf{Solution détaillée :}

\textbf{1. Conjugaison :}
\begin{enumerate}
    \item \all{konnte} (\all{können} : ich kann, du kannst, er kann / ich \textbf{konnte}, du \textbf{konntest}...). \textit{Perte de l'Umlaut au Präteritum.}
    \item \all{haben... machen müssen} (Perfekt : \all{haben} + Infinitif \all{machen} + Infinitif modal \all{müssen}). \textit{Règle du double infinitif.}
    \item \all{werden... lesen sollen} (Futur I : \all{werden} + Infinitif \all{lesen} + Infinitif modal \all{sollen}).
    \item \all{würde... gerne werden} ou \all{wollte} (Präteritum/Konjunktiv II forme simple). \textbf{Standard :} \all{Er wollte gerne Arzt werden.} (Präteritum de \all{wollen} = Konjunktiv II forme simple pour \all{wollen}). \textit{Note :} \all{würde werden} sonne mal (double \all{werden}). On préfère \all{wollte}.
\end{enumerate}

\textbf{2. Transformations :}
\begin{enumerate}
    \item \all{Du musst hier nicht parken.} (\all{nicht müssen} = pas l'obligation). \textit{Attention :} \all{Du musst nicht parken} = Tu n'es pas obligé de garer (tu peux si tu veux). \all{Du darfst nicht parken} = Il t'est interdit de garer.
    \item \all{..., weil er gut Deutsch sprechen kann.} (Subordonnée \all{weil} $\to$ verbe conjugué \all{kann} à la \textbf{fin}. Infinitif \all{sprechen} avant \all{kann}).
    \item \all{Wir haben das Buch lesen wollen.} (Perfekt double infinitif : \all{haben} + \all{lesen} + \all{wollen}).
\end{enumerate}

\textbf{3. Thème :}
\begin{enumerate}
    \item \all{Wir müssen den Vokabeln lernen.} ou \all{Wir müssen Vokabeln lernen.} (\all{müssen} = obligation intérieure/nécessité). \textit{Attention :} \all{lernen} + Akk (\all{die Vokabeln}), pas \all{zu} ici (verbe modal).
    \item \all{Du darfst dein Handy im Unterricht nicht benutzen.} (\all{dürfen} + \all{nicht} = interdiction. \all{im Unterricht} = en classe. \all{benutzen} / \all{nutzen}).
    \item \all{Er hat das Problem lösen können.} (\all{können} = capacité/réussite $\to$ Perfekt double infinitif \all{haben... lösen können}). \textit{Pas :} \all{Er konnte das Problem lösen} (Präteritum, correct aussi mais consigne Perfekt implicite ou Präteritum accepté). Si Prétérit : \all{Er konnte das Problem lösen.}
    \item \all{Sie sollten zum Arzt gehen.} (\all{sollen} Konjunktiv II \all{sollten} = conseil/recommandation). \all{gehen} (aller à pied) vs \all{fahren} (en voiture). \all{zum Arzt} = \all{zu dem Arzt}.
\end{enumerate}
\textbf{Conclusion :} La distinction \all{nicht müssen / nicht dürfen} et la structure du \cle{double infinitif} au Perfekt sont les deux points les plus testés au Bac.
\end{exoresolu}

\courssub{Méthode 4 : Exprimer la concession (obwohl / trotzdem)}

\begin{methode}[La concession : Struktur und Wortstellung]
\textbf{Objectif :} Relier deux idées contraires (Réalité vs Attente) sans les opposer frontalement (comme \all{aber/sondern}) mais en intégrant l'obstacle.
\begin{enumerate}
    \item \textbf{\all{obwohl} (conjonction de subordination) :} Introduit une \textbf{proposition subordonnée}. Le verbe conjugué \textbf{va à la fin}. La proposition \all{obwohl} peut être en \textbf{position 1} (sujet après verbe principal) ou en fin de phrase.
    \begin{itemize}
        \item \all{Obwohl es regnet, gehen wir spazieren.} (Pos 1 : verbe principal \all{gehen} en Pos 2).
        \item \all{Wir gehen spazieren, obwohl es regnet.} (Pos principale d'abord).
    \end{itemize}
    \item \textbf{\all{trotzdem} (adverbe de liaison / Konnektor) :} Relie deux \textbf{propositions principales} (ou phrases séparées). Il occupe une \textbf{place quelconque} (souvent Pos 1 ou Pos 3+), le verbe principal reste en \textbf{Pos 2}.
    \begin{itemize}
        \item \all{Es regnet. Trotzdem gehen wir spazieren.} (Pos 1 \all{Trotzdem} $\to$ \all{gehen} Pos 2).
        \item \all{Es regnet, wir gehen trotzdem spazieren.} (Pos 3+).
    \end{itemize}
    \item \textbf{\all{trotz} + Génitif (préposition) :} \all{Trotz des Regens gehen wir spazieren.} (Niveau supérieur, nominalisation).
    \item \textbf{Distinction \all{aber} vs \all{trotzdem} :} \all{aber} = simple opposition (mais). \all{trotzdem} = concession (malgré cela, néanmoins), insiste sur le fait que l'action principale se réalise \textbf{malgré} la cause énoncée avant.
\end{enumerate}
\textbf{Reflexe :} Si le sujet change $\to$ \all{obwohl} (subordonnée) ou \all{trotzdem} (phrase principale). Si même sujet $\to$ \all{trotz} + Nominalisation (élégant) ou \all{obwohl} + pronom.
\end{methode}

\begin{exoresolu}[Grammaire : Concession (obwohl / trotzdem / trotz) - Transformation \& Production]
\textbf{Énoncé :}
\begin{enumerate}
    \item Reliez les deux phrases par \textbf{\all{obwohl}} (deux possibilités de place).
    \begin{enumerate}
        \item Die Pressefreiheit ist ein Grundrecht. In manchen Ländern wird sie eingeschränkt.
        \item Die Journalisten arbeiten gefährlich. Sie decken Missstände auf.
    \end{enumerate}
    \item Reliez les deux phrases par \textbf{\all{trotzdem}}.
    \begin{enumerate}
        \item Die Demonstration war verboten. Viele Menschen kamen.
        \item Er hat Angst vor Repressionen. Er schreibt kritische Artikel.
    \end{enumerate}
    \item Transformez les phrases avec \textbf{\all{obwohl}} en phrases avec \textbf{\all{trotz} + Génitif} (nominalisation).
    \begin{enumerate}
        \item \all{Obwohl die Verfassung die Freiheit garantiert, gibt es Zensur.}
        \item \all{Obwohl er krank war, ging er zur Versammlung.}
    \end{enumerate}
    \item \textbf{Production courte :} Écrivez 3 phrases sur la situation des libertés dans votre pays (Cameroun / contexte local) en utilisant \all{obwohl}, \all{trotzdem}, \all{trotz}.
\end{enumerate}

\tcblower
\textbf{Solution détaillée :}

\textbf{1. Liaison avec \all{obwohl} (Subordonnée, verbe à la fin) :}
\begin{enumerate}
    \item \all{Obwohl die Pressefreiheit ein Grundrecht ist, wird sie in manchen Ländern eingeschränkt.} (Subordonnée en Pos 1 $\to$ inversion sujet/verbe principale : \all{wird sie...}).
    \item \all{In manchen Ländern wird die Pressefreiheit eingeschränkt, obwohl sie ein Grundrecht ist.} (Principale d'abord).
    \item \all{Obwohl die Journalisten gefährlich arbeiten, decken sie Missstände auf.} (\all{arbeiten} à la fin).
    \item \all{Die Journalisten decken Missstände auf, obwohl sie gefährlich arbeiten.}
\end{enumerate}

\textbf{2. Liaison avec \all{trotzdem} (Deux principales, verbe Pos 2) :}
\begin{enumerate}
    \item \all{Die Demonstration war verboten. Trotzdem kamen viele Menschen.} (\all{Trotzdem} Pos 1 $\to$ \all{kamen} Pos 2). Ou : \all{Die Demonstration war verboten, viele Menschen kamen trotzdem.}
    \item \all{Er hat Angst vor Repressionen. Trotzdem schreibt er kritische Artikel.} Ou : \all{Er hat Angst vor Repressionen, er schreibt trotzdem kritische Artikel.}
\end{enumerate}

\textbf{3. Nominalisation avec \all{trotz} + Génitif (Article \all{des} / \all{der} / \all{eines}) :}
\begin{enumerate}
    \item \all{Trotz der Garantie der Freiheit durch die Verfassung gibt es Zensur.} (\all{die Garantie} f., \all{die Freiheit} f., \all{die Verfassung} f. $\to$ \all{trotz der Garantie der Freiheit der Verfassung} -- éviter la chaîne de génitifs si possible, mieux : \all{Trotz der verfassungsmäßigen Garantie der Freiheit...}). \textit{Version standard attendue :} \all{Trotz der Garantie der Freiheit durch die Verfassung gibt es Zensur.}
    \item \all{Trotz seiner Krankheit ging er zur Versammlung.} (\all{die Krankheit} f. $\to$ \all{seiner Krankheit}. Possessif \all{sein} prend marqueur génitif \all{-r}).
\end{enumerate}

\textbf{4. Production guidée (Exemples types) :}
\begin{enumerate}
    \item \all{Obwohl die Verfassung in Kamerun die Meinungsfreiheit garantiert, gibt es manchmal Einschränkungen in der Praxis.}
    \item \all{Einige Journalisten haben Angst vor Verhaftung. Trotzdem berichten sie über Korruption.}
    \item \all{Trotz der Schwierigkeiten setzen sich die Jugendlichen für ihre Rechte ein.}
\end{enumerate}
\textbf{Conclusion :} \all{obwohl} structure la pensée (cause/conséquence intégrée), \all{trotzdem} marque l'argumentation orale/écrite (rebond), \all{trotz} + Génitif élève le style (écrit formel). Vérifier : \all{obwohl} $\to$ verbe \textbf{à la fin} ; \all{trotzdem} $\to$ verbe \textbf{Pos 2}.
\end{exoresolu}

\courssub{Méthode 5 : Rédiger un commentaire structuré / Production guidée}

\begin{methode}[Le commentaire de texte / La production écrite argumentée (Niveau Bac)]
\textbf{Objectif :} Produire un paragraphe cohérent (120--150 mots) ou répondre à une question de synthèse en structurant sa pensée.
\begin{enumerate}
    \item \textbf{Analyse du sujet :} Identifier le \cle{thème} (libertés), l'\cle{aspect} (rôle, limites, responsabilité), le \cle{format} (commentaire, article de journal, lettre au directeur, discours) et le \cle{registre} (formel \all{Sie} / standard).
    \item \textbf{Plan en 3 étapes (Dialectique simple) :}
        \begin{itemize}
            \item \textbf{These (These) :} Affirmation générale + Exemple du texte / connaissance. \textit{Einleitung : „In einer Demokratie sind Grundrechte unverzichtbar...“}
            \item \textbf{Antithese (Antithese) :} Nuance, limites, contre-exemple, concession. \textit{Verbindung : „Allerdings / Zwar... aber / Obwohl...“} (\cle{obwohl}, \cle{trotzdem}).
            \item \textbf{Synthese (Synthese) :} Réponse personnelle, ouverture, condition de l'exercice des libertés. \textit{Schluss : „Letztlich hängt die Freiheit von der Verantwortung jedes Einzelnen ab.“}
        \end{itemize}
    \item \textbf{Connecteurs logiques (Redemittel) :} \all{Zunächst, Darüber hinaus, Einerseits... andererseits, Daher, Folglich, Zusammenfassend}.
    \item \textbf{Relecture (Check-list) :} \cle{Ordre des mots} (Verbe Pos 2 / Fin), \cle{Cas} (Nom/Acc/Dat/Gén), \cle{Accords} (Sujet-Verbe, Adjectifs), \cle{Orthographe} (Majuscules noms, Umlaute, ß), \cle{Temps} (cohérence).
\end{enumerate}
\textbf{Structure type intro :} \all{Der Text „...“ thematisiert die Frage nach... / Meiner Meinung nach ist... entscheidend, weil...}
\end{methode}

\begin{exoresolu}[Production écrite : Commentaire argumenté]
\textbf{Énoncé :} \textit{« La liberté d'expression est le pilier de la démocratie, mais elle n'est pas illimitée. »} Rédigez un commentaire structuré (environ 130 mots) en allemand. Utilisez le vocabulaire de la leçon (\all{Meinungsfreiheit, Hassrede, Verantwortung, Grundgesetz, Zensur, obwohl, trotzdem, modalverben}) et au moins une phrase avec \all{obwohl} et une avec \all{trotzdem}.

\tcblower
\textbf{Solution détaillée :}

\textbf{Plan de rédaction :}
\begin{itemize}
    \item \textbf{These :} Définition et importance (Art. 5 GG, pilier).
    \item \textbf{Antithese :} Limites nécessaires (Hassrede, Schutz anderer), concession (\all{obwohl}).
    \item \textbf{Synthese :} Responsabilité du citoyen, équilibre (\all{trotzdem}).
\end{itemize}

\textbf{Texte modèle (Niveau Bac Excellent) :}
\begin{textebox}[Modelltext : Die Grenzen der Meinungsfreiheit]
Die Meinungsfreiheit ist das Fundament jeder Demokratie. Artikel 5 des Grundgesetzes garantiert jedem das Recht, seine Meinung frei zu äußern. \textbf{Ohne diese Freiheit} gäbe es keine öffentliche Debatte und keine Kontrolle der Macht. \textbf{Allerdings} ist sie nicht grenzenlos. \textbf{Obwohl} die Verfassung die freie Äußerung schützt, \textbf{muss} der Staat eingreifen, wenn Hassrede oder Hetze die Würde anderer verletzen. \textbf{Man darf} nicht alles sagen, was man denkt. \textbf{Trotzdem} bleiben kritische Stimmen unverzichtbar für den Fortschritt. Journalisten \textbf{sollen} Missstände aufdecken, \textbf{ohne} die Wahrheit zu verletzen. \textbf{Zusammenfassend} lässt sich sagen: Freiheit erfordert Verantwortung. Wer Rechte für sich beansprucht, \textbf{muss} auch die Pflichten gegenüber der Gemeinschaft akzeptieren. Nur so gelingt der gesellschaftliche Zusammenhalt.
\end{textebox}

\textbf{Analyse grammaticale du modèle (Points clés pour le Bac) :}
\begin{itemize}
    \item \textbf{Modaux variés :} \all{gäbe} (Konj II \all{geben} - irrealité), \all{muss} (obligation objective), \all{darf nicht} (interdiction), \all{sollen} (recommandation/expectation), \all{muss} (nécessité finale).
    \item \textbf{Concession :} \all{Obwohl die Verfassung... schützt, muss der Staat...} (Verbe \all{schützt} à la fin subordonnée ; \all{muss} Pos 2 principale). \all{Trotzdem bleiben...} (Adverbe Pos 1 $\to$ verbe Pos 2).
    \item \textbf{Structure infinitif :} \all{ohne die Wahrheit zu verletzen} (Préposition \all{ohne} + \all{zu}-Infinitif).
    \item \textbf{Vocabulaire précis :} \all{das Fundament}, \all{die öffentliche Debatte}, \all{die Menschenwürde} (impliqué par \all{Würde anderer}), \all{der gesellschaftliche Zusammenhalt}.
    \item \textbf{Orthographe :} Majuscules sur \textbf{tous} les noms (\all{Freiheit, Demokratie, Grundgesetz, Recht, Meinung, Staat, Hassrede, Würde, Journalisten, Missstände, Wahrheit, Fortschritt, Verantwortung, Rechte, Pflichten, Gemeinschaft, Zusammenhalt}).
\end{itemize}
\textbf{Conclusion :} Ce texte montre la maîtrise du thème (vocabulaire), de la grammaire (modaux, concession, ordre des mots) et de la cohérence argumentative (connecteurs).
\end{exoresolu}

\courssub{Attention : Les 5 erreurs qui coûtent des points au Bac}

\begin{attention}[Les 5 erreurs fatales sur le thème « Libertés \& Grammaire »]
\begin{enumerate}
    \item \textbf{Majuscules oubliées sur les noms (Großschreibung).} \textit{Exemple :} \all{*die meinungsfreiheit, *ein grundrecht, *die presse}. \textbf{Règle :} \cle{Tout nom (Substantiv) prend une majuscule.} \all{Die Meinungsfreiheit, ein Grundrecht, die Presse}. Vérifiez chaque mot qui peut être précédé de \all{der/die/das}.
    \item \textbf{Confusion \all{nicht müssen} vs \all{nicht dürfen}.} \textit{Erreur :} \all{*Du musst nicht rauchen} (pour dire « Tu n'as pas le droit de fumer »). \textbf{Correction :} \all{Du darfst nicht rauchen.} (\all{nicht müssen} = pas obligé ; \all{nicht dürfen} = interdit). \textit{Astuce :} \all{verbieten} $\leftrightarrow$ \all{nicht dürfen}.
    \item \textbf{Place du verbe avec \all{obwohl} (Subordonnée).} \textit{Erreur :} \all{*Obwohl er ist krank, kommt er.} \textbf{Correction :} \all{Obwohl er krank \textbf{ist}, kommt er.} Verbe conjugué \textbf{à la toute fin} de la subordonnée. Si \all{obwohl} en début de phrase $\to$ inversion en principale (\all{kommt er}).
    \item \textbf{Double infinitif au Perfekt mal formé.} \textit{Erreur :} \all{*Er hat gemusst arbeiten} ou \all{*Er ist arbeiten müssen}. \textbf{Correction :} \all{Er hat arbeiten \textbf{müssen}.} Auxiliaire \all{haben} (sauf verbes de mouvement/état), pas de \all{ge-} sur le modal, ordre : Infinitif principal + Infinitif modal.
    \item \textbf{Calque français sur « permettre/interdire à quelqu'un de faire qc ».} \textit{Erreur :} \all{*Er erlaubt mir zu gehen} (manque objet) ou \all{*Er verbietet mir gehen}. \textbf{Structure allemande :} \all{jemandem (Dat) etwas (Akk) erlauben / verbieten.} \all{Er erlaubt mir, zu gehen.} / \all{Er verbietet mir das Rauchen.} / \all{Er verbietet mir, zu rauchen.} (Dat \textbf{mir} + Infinitif avec \all{zu} / ou Nominalisation Akk).
\end{enumerate}
\textbf{Conseil final :} Relisez votre copie \textbf{uniquement pour ces 5 points} dans les 10 dernières minutes de l'épreuve.
\end{attention}

\newpage

\coursec{Exercices}

\courssub{Je vérifie mes connaissances}

\begin{exercice}[QCM : les verbes de modalité]
Entoure la bonne réponse.
\begin{enumerate}
\item \all{Die Bürger \_\_\_ frei ihre Meinung sagen.} \\
a) \all{dürfen} \quad b) \all{müssen} \quad c) \all{wollen}
\item \all{Jeder Journalist \_\_\_ die Wahrheit berichten.} (devoir moral) \\
a) \all{kann} \quad b) \all{soll} \quad c) \all{will}
\item \all{In einer Demokratie \_\_\_ man demonstrieren, wenn man es anmeldet.} (permission) \\
a) \all{darf} \quad b) \all{muss} \quad c) \all{soll}
\item \all{Wer wählen gehen \_\_\_, muss seinen Ausweis mitbringen.} (volonté) \\
a) \all{kann} \quad b) \all{will} \quad c) \all{darf}
\end{enumerate}
\end{exercice}

\begin{exercice}[Vrai ou faux ? Justifie ta réponse]
Réponds par \emph{vrai} ou \emph{faux} et justifie en une phrase en français.
\begin{enumerate}
\item Le verbe modal se conjugue et se place en deuxième position ; l'infinitif va à la fin de la phrase.
\item \all{Obwohl} est une conjonction de coordination : le verbe conjugué reste en deuxième position.
\item \all{Trotzdem} occupe une place dans la phrase : le verbe conjugué vient juste après, puis le sujet.
\item \all{Die Freiheit} est un nom masculin.
\end{enumerate}
\end{exercice}

\begin{exercice}[Vocabulaire : articles et pluriels]
Complète le tableau avec l'article et le pluriel de chaque nom.
\begin{center}
\begin{tabular}{|l|l|l|}
\hline
\rowcolor{popblueL} Nom & Article & Pluriel \\
\hline
Freiheit & \ldots & \ldots \\
\hline
Recht & \ldots & \ldots \\
\hline
Pflicht & \ldots & \ldots \\
\hline
Meinung & \ldots & \ldots \\
\hline
Demonstration & \ldots & \ldots \\
\hline
Zeitung & \ldots & \ldots \\
\hline
\end{tabular}
\end{center}
\end{exercice}

\begin{exercice}[Conjugaison à compléter]
Complète avec le verbe modal indiqué, conjugué au présent.
\begin{enumerate}
\item Die Bürger \_\_\_ frei wählen. (\all{dürfen})
\item Ihr \_\_\_ die Gesetze respektieren. (\all{müssen})
\item Wir \_\_\_ unsere Religion frei ausüben. (\all{können})
\item Du \_\_\_ deine Pflichten nicht vergessen. (\all{sollen})
\item Die Jugendlichen \_\_\_ an der Demonstration teilnehmen. (\all{wollen})
\item Man \_\_\_ hier nicht rauchen, es ist verboten. (\all{dürfen})
\end{enumerate}
\end{exercice}

\courssub{Je m'entraîne}

\begin{exercice}[Reconstruction de phrases]
Remets les mots dans l'ordre pour former une phrase correcte. Attention à la place du verbe.
\begin{enumerate}
\item Meinung / jeder / sagen / seine / darf
\item dürfen / in einer Demokratie / die Bürger / frei / wählen
\item obwohl / verboten / ist / sie / die Zeitung / erscheint
\item müssen / trotzdem / wir / respektieren / die Gesetze
\item seine / kann / jeder / Religion / wählen / frei
\end{enumerate}
\end{exercice}

\begin{exercice}[Transformation : la concession]
Relie les deux phrases d'abord avec \all{obwohl}, puis avec \all{trotzdem}. Écris les deux versions.
\begin{enumerate}
\item Die Presse ist frei. Sie muss die Wahrheit sagen.
\item Die Demonstration war verboten. Viele Leute sind gekommen.
\item Er hat das Recht zu wählen. Er geht nie wählen.
\end{enumerate}
\end{exercice}

\begin{exercice}[Texte à trous : modaux et vocabulaire]
Complète le texte avec les mots suivants : \all{dürfen -- Meinungsfreiheit -- müssen -- Recht -- Pflicht -- Pressefreiheit}.

\begin{textebox}[Freiheiten und Pflichten]
In einer Demokratie haben die Bürger viele Rechte. Die \_\_\_\_\_\_\_\_\_\_\_\_ erlaubt jedem, seine Meinung frei zu sagen. Dank der \_\_\_\_\_\_\_\_\_\_\_\_ können Journalisten kritische Artikel schreiben. Jeder Bürger hat das \_\_\_\_\_\_\_\_\_\_\_\_ zu wählen. Aber Freiheit bedeutet auch Verantwortung : alle \_\_\_\_\_\_\_\_\_\_\_\_ die Gesetze respektieren, das ist eine \_\_\_\_\_\_\_\_\_\_\_\_. Niemand \_\_\_\_\_\_\_\_\_\_\_\_ die Rechte der anderen verletzen.
\end{textebox}
\end{exercice}

\begin{exercice}[Appariements : mots et définitions]
Relie chaque mot allemand à sa définition française.
\begin{center}
\begin{tabular}{|l|l|}
\hline
\rowcolor{popblueL} Mot allemand & Définition française \\
\hline
1. \all{die Meinungsfreiheit} & a. interdire \\
\hline
2. \all{das Recht} & b. ce qu'on doit faire \\
\hline
3. \all{die Pflicht} & c. liberté de dire ce qu'on pense \\
\hline
4. \all{verbieten} & d. autoriser \\
\hline
5. \all{erlauben} & e. ce que la loi garantit au citoyen \\
\hline
\end{tabular}
\end{center}
\end{exercice}

\courssub{Je me prépare au Bac}

\begin{exercice}[Compréhension de l'écrit type Bac]
Lis le texte suivant, puis réponds aux questions.

\begin{textebox}[Pressefreiheit in Deutschland]
In Deutschland ist die Pressefreiheit ein Grundrecht. Sie steht im Grundgesetz, der Verfassung des Landes. Journalisten dürfen frei berichten, auch über die Politik und die Regierung. Der Staat darf keine Zeitung zensieren, obwohl manche Artikel die Politiker kritisieren. Trotzdem gibt es Grenzen : niemand darf lügen oder andere Menschen beleidigen. Wer das tut, muss vor Gericht erscheinen. Viele Deutsche lesen jeden Tag Zeitung oder informieren sich im Internet. Für die Demokratie ist eine freie Presse sehr wichtig, denn die Bürger müssen gut informiert sein, wenn sie wählen gehen. Ohne freie Medien kann es keine echte Demokratie geben, sagen viele Experten.
\end{textebox}

\textbf{A. Richtig oder falsch ?} Cite le texte pour justifier.
\begin{enumerate}
\item Die Pressefreiheit steht im Grundgesetz.
\item Der Staat darf Zeitungen zensieren.
\item Journalisten dürfen alles schreiben, auch Lügen.
\item Viele Deutsche informieren sich nur im Internet.
\end{enumerate}

\textbf{B. Questions ouvertes} (réponds en allemand, avec des phrases complètes).
\begin{enumerate}
\item Was dürfen Journalisten in Deutschland?
\item Welche Grenzen gibt es für die Presse?
\item Warum ist eine freie Presse wichtig für die Demokratie?
\end{enumerate}

\textbf{C. Questions de langue.}
\begin{enumerate}
\item Relève dans le texte deux verbes de modalité conjugués et donne leur infinitif.
\item Relève une phrase avec \all{obwohl} et transforme-la avec \all{trotzdem}.
\end{enumerate}
\end{exercice}

\begin{exercice}[Thème et version]
\textbf{A. Thème.} Traduis en allemand.
\begin{enumerate}
\item Bien que la liberté soit importante, elle a des limites. Pourtant, nous devons lutter pour nos droits.
\item Chaque citoyen peut exprimer son opinion, mais il doit respecter les lois.
\item Les journalistes ne doivent pas mentir, même si la presse est libre.
\end{enumerate}

\textbf{B. Version.} Traduis en français le paragraphe suivant.

\begin{textebox}[Religionsfreiheit in der Schweiz]
In der Schweiz können alle Menschen ihre Religion frei wählen und ausüben. Das steht in der Verfassung. Obwohl die meisten Schweizer Christen sind, leben dort auch Muslime, Juden und Menschen ohne Religion. Der Staat darf niemanden wegen seiner Religion bestrafen. Trotzdem müssen alle die Gesetze des Landes respektieren.
\end{textebox}
\end{exercice}

\begin{exercice}[Expression écrite type Bac]
Sujet : \all{Welche Freiheit ist für junge Leute in Kamerun am wichtigsten?}

Rédige un texte de \textbf{60 à 80 mots} en allemand. Tu dois :
\begin{itemize}
\item choisir une liberté (expression, presse, religion, réunion) et expliquer ton choix ;
\item employer \textbf{au moins un verbe de modalité} (\all{dürfen, müssen, können, sollen, wollen}) ;
\item employer \textbf{au moins une concession} avec \all{obwohl} ou \all{trotzdem} ;
\item donner un exemple concret de la vie des jeunes au Cameroun (école, marché, sport, réseaux sociaux).
\end{itemize}
\end{exercice}

\newpage

\coursec{Corrigés des exercices}

\begin{corrige}[Exercice 1 — QCM : les verbes de modalité]
\begin{enumerate}
\item \textbf{a) \all{dürfen}} — \all{Die Bürger dürfen ihre Meinung frei sagen.} « Les citoyens ont le droit d'exprimer librement leur opinion. » \all{dürfen} exprime la permission/le droit, ce qui correspond au sens de la phrase.
\item \textbf{b) \all{soll}} — \all{Jeder Journalist soll die Wahrheit berichten.} « Tout journaliste doit rapporter la vérité. » \all{sollen} exprime le devoir moral, la recommandation ou une attente.
\item \textbf{a) \all{darf}} — \all{In einer Demokratie darf man demonstrieren, wenn man es anmeldet.} « En démocratie, on a le droit de manifester si on la déclare. » La permission se rend par \all{dürfen} ; sujet \all{man}, donc 3\textsuperscript{e} personne du singulier : \all{darf}.
\item \textbf{b) \all{will}} — \all{Wer wählen gehen will, muss seinen Ausweis mitbringen.} « Celui qui veut aller voter doit apporter sa pièce d'identité. » La volonté se rend par \all{wollen} ; \all{wer} = 3\textsuperscript{e} personne du singulier : \all{will}.
\end{enumerate}
\textbf{Point de grammaire :} le modal conjugué est en deuxième position et l'infinitif du verbe principal va à la fin de la phrase : \all{Die Bürger \textbf{dürfen} ihre Meinung frei \textbf{sagen}.}
\end{corrige}

\begin{corrige}[Exercice 2 — Vrai ou faux ? Justifie ta réponse]
\begin{enumerate}
\item \textbf{Vrai.} Le verbe modal se conjugue et occupe la deuxième position de la phrase principale ; l'infinitif du verbe principal est rejeté à la fin : \all{Man \textbf{darf} hier nicht \textbf{rauchen}.}
\item \textbf{Faux.} \all{Obwohl} est une conjonction de \emph{subordination} : elle renvoie le verbe conjugué à la fin de la subordonnée. Correct : \all{Obwohl die Presse frei \textbf{ist}, ...} (et non \all{*obwohl die Presse \textbf{ist} frei}).
\item \textbf{Vrai.} \all{Trotzdem} est un adverbe (conjonctif) qui occupe la première position ; le verbe conjugué reste en deuxième position et le sujet suit (inversion) : \all{\textbf{Trotzdem müssen wir} die Gesetze respektieren.}
\item \textbf{Faux.} \all{Die Freiheit} est un nom \emph{féminin} (pluriel : \all{die Freiheiten}). Les noms en \all{-heit} sont toujours féminins, comme \all{die Meinungsfreiheit}, \all{die Pressefreiheit}, \all{die Religionsfreiheit}.
\end{enumerate}
\end{corrige}

\begin{corrige}[Exercice 3 — Vocabulaire : articles et pluriels]
\begin{center}
\begin{tabularx}{\linewidth}{|X|X|X|}
\hline
\rowcolor{popblueL} \textbf{Nom} & \textbf{Article} & \textbf{Pluriel} \\
\hline
\all{Freiheit} & \all{die} & \all{die Freiheiten} \\
\hline
\all{Recht} & \all{das} & \all{die Rechte} \\
\hline
\all{Pflicht} & \all{die} & \all{die Pflichten} \\
\hline
\all{Meinung} & \all{die} & \all{die Meinungen} \\
\hline
\all{Demonstration} & \all{die} & \all{die Demonstrationen} \\
\hline
\all{Zeitung} & \all{die} & \all{die Zeitungen} \\
\hline
\end{tabularx}
\end{center}
\textbf{Astuces :} les noms en \all{-heit}, \all{-ung}, \all{-ion}, \all{-keit}, \all{-schaft}, \all{-tät} sont féminins (\all{die}). \all{Das Recht} est un faux ami de genre : en français « le droit » est masculin, en allemand il est neutre (\all{das}).
\end{corrige}

\begin{corrige}[Exercice 4 — Conjugaison à compléter]
\begin{enumerate}
\item \all{Die Bürger \textbf{dürfen} frei wählen.} (3\textsuperscript{e} pers. pluriel : \all{dürfen})
\item \all{Ihr \textbf{müsst} die Gesetze respektieren.} (2\textsuperscript{e} pers. pluriel : \all{müsst} — attention à l'Umlaut \all{ü})
\item \all{Wir \textbf{können} unsere Religion frei ausüben.} (1\textsuperscript{re} pers. pluriel : \all{können})
\item \all{Du \textbf{sollst} deine Pflichten nicht vergessen.} (2\textsuperscript{e} pers. singulier : \all{sollst})
\item \all{Die Jugendlichen \textbf{wollen} an der Demonstration teilnehmen.} (3\textsuperscript{e} pers. pluriel : \all{wollen})
\item \all{Man \textbf{darf} hier nicht rauchen, es ist verboten.} (3\textsuperscript{e} pers. singulier : \all{darf}, sans Umlaut)
\end{enumerate}
\textbf{Rappel :} au singulier, les modaux changent de voyelle (sauf \all{sollen}) : \all{dürfen → darf}, \all{müssen → muss}, \all{können → kann}, \all{wollen → will}, \all{sollen → soll}. La 1\textsuperscript{re} et la 3\textsuperscript{e} personne du singulier sont identiques : \all{ich darf / er darf}.
\end{corrige}

\begin{corrige}[Exercice 5 — Reconstruction de phrases]
\begin{enumerate}
\item \all{Jeder darf seine Meinung sagen.} « Chacun a le droit d'exprimer son opinion. » Sujet \all{jeder} (3\textsuperscript{e} pers. sing.) → \all{darf} ; infinitif \all{sagen} à la fin.
\item \all{Die Bürger dürfen in einer Demokratie frei wählen.} « Les citoyens ont le droit de voter librement dans une démocratie. »
\item \all{Die Zeitung erscheint, obwohl sie verboten ist.} « Le journal paraît bien qu'il soit interdit. » Après \all{obwohl}, le verbe conjugué \all{ist} va à la fin de la subordonnée.
\item \all{Trotzdem müssen wir die Gesetze respektieren.} « Pourtant, nous devons respecter les lois. » \all{Trotzdem} en tête → inversion : verbe puis sujet.
\item \all{Jeder kann seine Religion frei wählen.} « Chacun peut choisir librement sa religion. »
\end{enumerate}
\textbf{Point de vigilance :} ne jamais écrire \all{*obwohl sie ist verboten} : la subordonnée introduite par \all{obwohl} exige le verbe en position finale.
\end{corrige}

\begin{corrige}[Exercice 6 — Transformation : la concession]
\begin{enumerate}
\item \all{Obwohl die Presse frei ist, muss sie die Wahrheit sagen.} \\
\all{Die Presse ist frei. Trotzdem muss sie die Wahrheit sagen.} \\
« Bien que la presse soit libre, elle doit dire la vérité. / La presse est libre. Pourtant, elle doit dire la vérité. »
\item \all{Obwohl die Demonstration verboten war, sind viele Leute gekommen.} \\
\all{Die Demonstration war verboten. Trotzdem sind viele Leute gekommen.} \\
« Bien que la manifestation ait été interdite, beaucoup de gens sont venus. »
\item \all{Obwohl er das Recht zu wählen hat, geht er nie wählen.} \\
\all{Er hat das Recht zu wählen. Trotzdem geht er nie wählen.} \\
« Bien qu'il ait le droit de voter, il ne va jamais voter. »
\end{enumerate}
\textbf{Règle :} avec \all{obwohl}, une seule phrase complexe : subordonnée (verbe final) + principale (inversion si la subordonnée précède). Avec \all{trotzdem}, deux phrases indépendantes ; \all{trotzdem} en position 1 provoque l'inversion sujet/verbe dans la seconde phrase.
\end{corrige}

\begin{corrige}[Exercice 7 — Texte à trous : modaux et vocabulaire]
\textbf{Texte complété :}
\begin{center}
\all{In einer Demokratie haben die Bürger viele Rechte. Die \textbf{Meinungsfreiheit} erlaubt jedem, seine Meinung frei zu sagen. Dank der \textbf{Pressefreiheit} können Journalisten kritische Artikel schreiben. Jeder Bürger hat das \textbf{Recht} zu wählen. Aber Freiheit bedeutet auch Verantwortung : alle \textbf{müssen} die Gesetze respektieren, das ist eine \textbf{Pflicht}. Niemand \textbf{darf} die Rechte der anderen verletzen.}
\end{center}
\textbf{Justifications :}
\begin{itemize}
\item \all{Die Meinungsfreiheit} (Nominatif) : c'est elle qui permet de dire librement son opinion ; article \all{die} car nom féminin en \all{-heit}.
\item \all{der Pressefreiheit} : datif féminin après la préposition \all{dank} (« grâce à »).
\item \all{das Recht zu wählen} : « le droit de voter » ; \all{das Recht} est neutre.
\item \all{alle müssen} : « tous doivent » — obligation, 3\textsuperscript{e} pers. pluriel.
\item \all{eine Pflicht} : « un devoir », nom féminin (prédicat après \all{sein}).
\item \all{Niemand darf} : « personne n'a le droit » — interdiction, 3\textsuperscript{e} pers. singulier (\all{darf}).
\end{itemize}
\end{corrige}

\begin{corrige}[Exercice 8 — Appariements : mots et définitions]
\begin{center}
\begin{tabularx}{\linewidth}{|X|X|}
\hline
\rowcolor{popblueL} \textbf{Mot allemand} & \textbf{Définition française} \\
\hline
1. \all{die Meinungsfreiheit} & \textbf{c.} liberté de dire ce qu'on pense \\
\hline
2. \all{das Recht} & \textbf{e.} ce que la loi garantit au citoyen \\
\hline
3. \all{die Pflicht} & \textbf{b.} ce qu'on doit faire \\
\hline
4. \all{verbieten} & \textbf{a.} interdire \\
\hline
5. \all{erlauben} & \textbf{d.} autoriser \\
\hline
\end{tabularx}
\end{center}
\textbf{Familles de mots :} \all{die Freiheit → frei (libre) → die Meinungsfreiheit, die Pressefreiheit, die Religionsfreiheit, die Versammlungsfreiheit}. Attention : \all{verbieten} est un verbe fort : \all{verbietet -- verbot -- hat verboten}.
\end{corrige}

\begin{corrige}[Exercice 9 — Compréhension de l'écrit type Bac]
\textbf{Barème indicatif (sur 20) :} A : 4 × 1,5 = 6 pts (réponse + citation) ; B : 3 × 2 = 6 pts ; C : 2 × 2 = 4 pts ; qualité de la langue : 4 pts.

\textbf{A. Richtig oder falsch ?}
\begin{enumerate}
\item \textbf{Richtig.} \all{„Sie steht im Grundgesetz, der Verfassung des Landes.“}
\item \textbf{Falsch.} \all{„Der Staat darf keine Zeitung zensieren.“} (L'État n'a \emph{pas} le droit de censurer.)
\item \textbf{Falsch.} \all{„Niemand darf lügen oder andere Menschen beleidigen.“} (Il existe des limites : interdiction de mentir et d'insulter.)
\item \textbf{Falsch.} \all{„Viele Deutsche lesen jeden Tag Zeitung oder informieren sich im Internet.“} (Ils s'informent aussi par les journaux, pas seulement par Internet.)
\end{enumerate}

\textbf{B. Questions ouvertes — réponses modèles :}
\begin{enumerate}
\item \all{Journalisten dürfen frei berichten, auch über die Politik und die Regierung.} « Les journalistes ont le droit d'informer librement, y compris sur la politique et le gouvernement. »
\item \all{Niemand darf lügen oder andere Menschen beleidigen. Wer das tut, muss vor Gericht erscheinen.} « Personne n'a le droit de mentir ou d'insulter autrui ; qui le fait doit comparaître devant le tribunal. »
\item \all{Eine freie Presse ist wichtig, denn die Bürger müssen gut informiert sein, wenn sie wählen gehen. Ohne freie Medien kann es keine echte Demokratie geben.} « Une presse libre est importante parce que les citoyens doivent être bien informés quand ils vont voter. Sans médias libres, il ne peut pas y avoir de vraie démocratie. »
\end{enumerate}

\textbf{C. Questions de langue :}
\begin{enumerate}
\item Exemples de modaux conjugués relevés : \all{dürfen} (inf. \all{dürfen}), \all{darf} (inf. \all{dürfen}), \all{muss} (inf. \all{müssen}), \all{müssen} (inf. \all{müssen}), \all{kann} (inf. \all{können}). Deux suffisent.
\item Phrase relevée : \all{„Der Staat darf keine Zeitung zensieren, obwohl manche Artikel die Politiker kritisieren.“} \\
Transformation avec \all{trotzdem} : \all{Manche Artikel kritisieren die Politiker. Trotzdem darf der Staat keine Zeitung zensieren.} \\
« Certains articles critiquent les hommes politiques. Pourtant, l'État n'a le droit de censurer aucun journal. »
\end{enumerate}
\textbf{Points de vigilance :} dans les réponses, reprendre la structure du texte avec verbe en deuxième position ; ne pas oublier l'inversion après \all{trotzdem} ; \all{das Grundgesetz} = « la Loi fondamentale » (la Constitution allemande).
\end{corrige}

\begin{corrige}[Exercice 10 — Thème et version]
\textbf{A. Thème — traductions modèles :}
\begin{enumerate}
\item \all{Obwohl die Freiheit wichtig ist, hat sie Grenzen. Trotzdem müssen wir für unsere Rechte kämpfen.} \\
Justification : \all{obwohl} renvoie \all{ist} en fin de subordonnée ; principale en inversion (\all{hat sie}) ; \all{trotzdem} provoque l'inversion (\all{müssen wir}) ; « lutter pour » = \all{kämpfen für} + accusatif.
\item \all{Jeder Bürger kann seine Meinung äußern, aber er muss die Gesetze respektieren.} \\
Justification : « exprimer son opinion » = \all{seine Meinung äußern} (avec \all{ß} après diphtongue \all{äu}) ; le modal \all{kann} en 2\textsuperscript{e} position, \all{äußern} à la fin ; \all{aber} (coordination) garde l'ordre normal ; \all{muss ... respektieren}.
\item \all{Journalisten dürfen nicht lügen, obwohl die Presse frei ist.} (ou : \all{..., auch wenn die Presse frei ist.}) \\
Justification : « ne pas avoir le droit / interdiction » = \all{nicht dürfen} ; « mentir » = \all{lügen} ; « même si » = \all{obwohl} ou \all{auch wenn}, verbe final (\all{ist}).
\end{enumerate}

\textbf{B. Version — traduction modèle :}

« En Suisse, tous les êtres humains peuvent choisir et pratiquer librement leur religion. Cela figure dans la Constitution. Bien que la plupart des Suisses soient chrétiens, y vivent aussi des musulmans, des juifs et des personnes sans religion. L'État n'a le droit de punir personne en raison de sa religion. Pourtant, tous doivent respecter les lois du pays. »

\textbf{Points de vigilance :} \all{die Religionsfreiheit} = « la liberté de religion » ; \all{ausüben} = « pratiquer, exercer » (verbe à particule inséparable) ; \all{jemanden wegen seiner Religion bestrafen} = « punir quelqu'un en raison de sa religion » (\all{wegen} + génitif) ; \all{die meisten} = « la plupart » ; \all{die Verfassung} = « la Constitution » (aussi \all{das Grundgesetz} en Allemagne).
\end{corrige}

\newpage

\coursec{Fiche bilan}

\begin{popfigure}
\begin{center}\tcbox[colback=popblue,colframe=popblue,coltext=white,arc=9pt,boxsep=4pt,left=6pt,right=6pt]{\bfseries\small AL25 Die Freiheiten}\end{center}
\vspace{-2pt}
\begin{tcbraster}[raster columns=3,raster equal height=rows,raster column skip=6pt,raster row skip=6pt,enhanced,blankest]
\begin{tcolorbox}[enhanced,colback=poporange,colframe=poporange,coltext=white,boxrule=0.9pt,arc=3pt,left=4pt,right=4pt,top=3pt,bottom=3pt,boxsep=1pt,halign=left]\bfseries\scriptsize Texte „Freiheiten in einer Demokratie“\end{tcolorbox}
\begin{tcolorbox}[enhanced,colback=popgreen,colframe=popgreen,coltext=white,boxrule=0.9pt,arc=3pt,left=4pt,right=4pt,top=3pt,bottom=3pt,boxsep=1pt,halign=left]\bfseries\scriptsize Lexique Freiheit, Recht, Pflicht\end{tcolorbox}
\begin{tcolorbox}[enhanced,colback=poppurple,colframe=poppurple,coltext=white,boxrule=0.9pt,arc=3pt,left=4pt,right=4pt,top=3pt,bottom=3pt,boxsep=1pt,halign=left]\bfseries\scriptsize Modalverben dürfen, müssen, können\end{tcolorbox}
\begin{tcolorbox}[enhanced,colback=poppink,colframe=poppink,coltext=white,boxrule=0.9pt,arc=3pt,left=4pt,right=4pt,top=3pt,bottom=3pt,boxsep=1pt,halign=left]\bfseries\scriptsize Concession obwohl, trotzdem\end{tcolorbox}
\begin{tcolorbox}[enhanced,colback=popgold,colframe=popgold,coltext=white,boxrule=0.9pt,arc=3pt,left=4pt,right=4pt,top=3pt,bottom=3pt,boxsep=1pt,halign=left]\bfseries\scriptsize Méthode commentaire de texte\end{tcolorbox}
\begin{tcolorbox}[enhanced,colback=popblue!70,colframe=popblue!70,coltext=white,boxrule=0.9pt,arc=3pt,left=4pt,right=4pt,top=3pt,bottom=3pt,boxsep=1pt,halign=left]\bfseries\scriptsize Production résumé, paragraphe\end{tcolorbox}
\end{tcbraster}

\legende{Carte mentale de la leçon AL25 : Les libertés}
\end{popfigure}

\begin{bilan}
\textbf{1. Lexique clé (à connaître avec article et pluriel)}
\begin{itemize}
\item \all{die Freiheit, -en} : la liberté ; \all{die Meinungsfreiheit} : la liberté d'expression ; \all{die Pressefreiheit} : la liberté de la presse ; \all{die Religionsfreiheit} : la liberté de religion ; \all{die Versammlungsfreiheit} : la liberté de réunion.
\item \all{das Recht, -e} : le droit ; \all{die Pflicht, -en} : le devoir ; \all{das Grundgesetz} : la Loi fondamentale (Constitution allemande) ; \all{die Demokratie, -n} : la démocratie ; \all{der Bürger, - / die Bürgerin, -nen} : le/la citoyen(ne).
\item \all{erlauben} : autoriser ; \all{verbieten (verbot, hat verboten)} : interdire ; \all{schützen} : protéger ; \all{respektieren} : respecter ; \all{sich äußern} : s'exprimer ; \all{demonstrieren} : manifester $\rightarrow$ \all{die Demonstration, -en} : la manifestation.
\item Familles de mots : \all{frei} (adj.) $\rightarrow$ \all{die Freiheit} $\rightarrow$ \all{freiwillig} (volontaire) ; \all{die Meinung} (l'avis) $\rightarrow$ \all{meinen} (penser, être d'avis) $\rightarrow$ \all{die Meinungsfreiheit} ; \all{das Recht} $\rightarrow$ \all{gerecht} (juste) $\rightarrow$ \all{die Gerechtigkeit} (la justice).
\end{itemize}

\textbf{2. Grammaire à connaître par cœur}

\begin{center}
\begin{tabularx}{\linewidth}{|l|X|X|}
\hline
\rowcolor{popblueL} \textbf{Point} & \textbf{Règle} & \textbf{Exemple} \\
\hline
Modalverben & Le modal se conjugue en 2\textsuperscript{e} position ; l'infinitif va à la fin de la phrase. & \all{Jeder Bürger \textbf{darf} seine Meinung frei \textbf{äußern}.} \\
\hline
dürfen / müssen & \all{dürfen} = avoir le droit ; \all{müssen} = devoir ; \all{nicht dürfen} = interdiction ; \all{nicht müssen} = ne pas être obligé. & \all{Niemand darf wegen seiner Religion diskriminiert werden.} \\
\hline
obwohl & Subordonnée de concession : le verbe conjugué va à la fin de la subordonnée. & \all{Obwohl die Pressefreiheit geschützt ist, bleibt sie fragil.} \\
\hline
trotzdem & Adverbe de concession : en 1\textsuperscript{re} position, il entraîne l'inversion verbe + sujet. & \all{Trotzdem gehen viele Menschen auf die Straße.} \\
\hline
\end{tabularx}
\end{center}

\textbf{Rappel : conjugaison des modalverben au Präsens}

\begin{center}
\begin{tabular}{|l|l|l|l|l|l|l|}
\hline
\rowcolor{popblueL} & \textbf{dürfen} & \textbf{können} & \textbf{müssen} & \textbf{sollen} & \textbf{wollen} & \textbf{mögen} \\
\hline
ich & darf & kann & muss & soll & will & mag \\
\hline
du & darfst & kannst & musst & sollst & willst & magst \\
\hline
er/sie/es & darf & kann & muss & soll & will & mag \\
\hline
wir & dürfen & können & müssen & sollen & wollen & mögen \\
\hline
ihr & dürft & könnt & müsst & sollt & wollt & mögt \\
\hline
sie/Sie & dürfen & können & müssen & sollen & wollen & mögen \\
\hline
\end{tabular}
\end{center}

\textbf{3. Méthodes express}
\begin{itemize}
\item \textbf{Compréhension} : lire le titre et le 1\textsuperscript{er} paragraphe (sens global) $\rightarrow$ repérer les mots-clés du thème $\rightarrow$ répondre en phrases complètes en allemand, sans recopier de longs passages du texte.
\item \textbf{Résumé} : une phrase par idée principale, au Präsens, sans citation longue, avec des connecteurs (\all{zuerst, außerdem, schließlich}).
\item \textbf{Concession} : vérifier la place du verbe (fin de la subordonnée après \all{obwohl} ; inversion verbe + sujet après \all{trotzdem}).
\item \textbf{Commentaire} : introduire (auteur, thème, problématique) $\rightarrow$ analyser les idées principales $\rightarrow$ donner son avis argumenté (\all{Meiner Meinung nach...}).
\end{itemize}
\end{bilan}

\begin{attention}[Pièges fréquents des francophones]
\begin{itemize}
\item Ne confonds pas \all{nicht dürfen} (il est interdit de) et \all{nicht müssen} (il n'est pas obligatoire de) : \all{Du darfst hier nicht rauchen.} « Il est interdit de fumer ici » ; \all{Du musst nicht kommen.} « Tu n'es pas obligé de venir ».
\item \all{das Recht} signifie « le droit » (juridique) ; « avoir raison » se dit \all{Recht haben} : \all{Du hast Recht.} « Tu as raison ».
\item Après \all{obwohl}, le verbe conjugué va à la fin : \all{Obwohl er das Recht hat, ...} et non \emph{Obwohl er hat das Recht}.
\item \all{trotzdem} est un adverbe : s'il ouvre la phrase, le sujet suit immédiatement le verbe : \all{Trotzdem demonstrieren die Bürger.}
\item Majuscule obligatoire à tous les noms : \all{die Freiheit, das Recht, die Pflicht}.
\end{itemize}
\end{attention}

\begin{aretenir}[Checklist avant le Bac]
\begin{itemize}
\item[$\square$] Je connais les articles et pluriels du lexique des libertés (\all{die Freiheit, -en ; das Recht, -e ; die Pflicht, -en}).
\item[$\square$] Je sais conjuguer les six verbes de modalité au Präsens (\all{dürfen, können, müssen, sollen, wollen, mögen}).
\item[$\square$] Je place correctement l'infinitif à la fin de la phrase avec un modal.
\item[$\square$] Je distingue \all{nicht dürfen} (interdiction) et \all{nicht müssen} (absence d'obligation).
\item[$\square$] Je mets le verbe conjugué à la fin dans une subordonnée introduite par \all{obwohl}.
\item[$\square$] Je fais l'inversion verbe + sujet après \all{trotzdem} en tête de phrase.
\item[$\square$] Je n'oublie pas la majuscule à tous les noms allemands (\all{die Freiheit, das Recht}).
\item[$\square$] Je sais résumer un texte en 3 à 5 phrases avec des connecteurs.
\item[$\square$] Je sais définir les quatre libertés : \all{Meinungs-, Presse-, Religions-, Versammlungsfreiheit}.
\item[$\square$] Je peux produire un court paragraphe d'opinion sur les libertés au Cameroun et dans les pays germanophones.
\end{itemize}
\end{aretenir}

\findecours

\end{document}
